\documentclass[uplatex,dvipdfmx,a4paper,11pt,openany]{jsbook}
%================================================
% Packages
%================================================
\usepackage{amsmath,amssymb,amsthm,mathtools}
\usepackage{bm}
\usepackage{physics}
\usepackage{mathrsfs}
\usepackage{booktabs}
\usepackage{array}
\usepackage{longtable}
\usepackage[dvipdfmx]{hyperref}
\usepackage{pxjahyper}
\usepackage{geometry}
\usepackage{enumitem}
\usepackage{xcolor}
\usepackage{tcolorbox}
\tcbuselibrary{theorems, skins, breakable}
% --- 定理系ボックスの共通スタイル ---
\tcbset{
    myboxstyle/.style={
        enhanced,
        breakable,
        colback=white,
        coltitle=white,
        fonttitle=\bfseries,
        boxrule=0.8pt,
        arc=1.5mm,
        attach boxed title to top left={xshift=3mm, yshift=-2mm},
        top=3mm,
        before skip=3mm, after skip=3mm
    }
}
% --- 定理 (Theorem) : レッド系 ---
\newtcbtheorem[number within=section]{mythm}{定理}{
    myboxstyle,
    colframe=red!75!black,
    boxed title style={
        colback=red!50,
        colframe=red!75!black,
        boxrule=0.8pt,
        arc=1.5mm
    }
}{thm}
% --- 定義 (Definition) : ブルー系 ---
\newtcbtheorem[use counter from=mythm]{mydef}{定義}{
    myboxstyle,
    colframe=blue!70!black,
    boxed title style={
        colback=blue!60!gray!80,
        colframe=blue!70!black,
        boxrule=0.8pt,
        arc=1.5mm
    }
}{def}
% --- 補題 (Lemma) : グリーン系 ---
\newtcbtheorem[use counter from=mythm]{mylem}{補題}{
    myboxstyle,
    colframe=green!50!black,
    boxed title style={
        colback=green!50!gray,
        colframe=green!50!black,
        boxrule=0.8pt,
        arc=1.5mm
    }
}{lem}
% --- 命題 (Proposition) : オレンジ系 ---
\newtcbtheorem[use counter from=mythm]{myprop}{命題}{
    myboxstyle,
    colframe=orange!80!black,
    boxed title style={
        colback=orange!80!black,
        colframe=orange!80!black,
        boxrule=0.8pt,
        arc=1.5mm
    }
}{prop}
% --- 例 (Example) : グレー系 ---
\newtcbtheorem[use counter from=mythm]{myex}{例}{
    myboxstyle,
    colframe=gray!70!black,
    boxed title style={
        colback=gray!70!black,
        colframe=gray!70!black,
        boxrule=0.8pt,
        arc=1.5mm
    }
}{ex}
% 用語：日本語と英語を併記（英語版では \renewcommand で切り替える）
%   \term{日本語}{English} … 初出の用語（太字＋英語）
%   \en{English}          … 英語表記のみ括弧で添える
\newcommand{\term}[2]{\textbf{#1}（#2）}
\newcommand{\en}[1]{（#1）}
% 計算過程などを書く注釈環境
\newenvironment{note}{\vspace{1ex}\par\footnotesize\noindent\textbf{【注】}}{\par\vspace{1ex}}

\geometry{margin=25mm}
\setlength{\parskip}{0.3em}
\setlength{\parindent}{1em}
\raggedbottom  % ページ下端を揃えるための余分な縦の伸びを防ぐ
\setlist{itemsep=0.15em, topsep=0.4em, parsep=0pt, partopsep=0pt}
\numberwithin{equation}{chapter}
%================================================
% Theorem environments
%================================================
\theoremstyle{definition}
\newtheorem{definition}{定義}[chapter]
\newtheorem{example}{例}[chapter]

\theoremstyle{plain}
\newtheorem{theorem}{定理}[chapter]
\newtheorem{proposition}{命題}[chapter]

\theoremstyle{remark}
\newtheorem{remark}{注意}[chapter]
%================================================
% Commands
%================================================
\newcommand{\Lag}{\mathcal{L}}
\newcommand{\SM}{\mathrm{SM}}
\newcommand{\SU}{\mathrm{SU}}
\newcommand{\U}{\mathrm{U}}
\newcommand{\SO}{\mathrm{SO}}
\newcommand{\diag}{\operatorname{diag}}
\newcommand{\hc}{\mathrm{h.c.}}
\newcommand{\PL}{P_{\mathrm{L}}}
\newcommand{\PR}{P_{\mathrm{R}}}

\title{
標準模型講義ノート\\
\large Standard Model of Particle Physics
}
\author{}
\date{\today}
\begin{document}
\maketitle
\tableofcontents
%================================================
% Part I
%================================================
\part{標準模型の全体像とゲージ理論}
%================================================
\chapter{標準模型の概要}
%================================================
本章では、後の章で詳しく導出する内容を先取りして、
標準模型（Standard Model, SM）の全体像を概観する。
ここで現れる式の多くは「結果」として提示し、
その導出は第2章以降で行う。
本章の目的は、標準模型が
「どのような場から成り」「どのような対称性を持ち」
「何個のパラメータで決まるか」を一望することにある。
\paragraph{記法の規約}
本ノートでは以下の規約を用いる。
\begin{itemize}
    \item \term{自然単位系}{natural units} $\hbar=c=1$ を用いる。換算には $\hbar c\simeq0.1973\ \mathrm{GeV\,fm}$ を使う。
    \item 計量\en{metric}は $\eta_{\mu\nu}=\diag(+1,-1,-1,-1)$ とする。
    \item $\gamma^5=i\gamma^0\gamma^1\gamma^2\gamma^3$、$\{\gamma^5,\gamma^\mu\}=0$、$(\gamma^5)^2=1$、$(\gamma^5)^\dagger=\gamma^5$。
    \item $\epsilon^{0123}=-\epsilon_{0123}=+1$ とする。
    \item ゲージ場の共変微分は $D_\mu=\partial_\mu-igA_\mu^aT^a$ の符号で定義する。
    \item 弱超電荷は $Q=T_3+Y/2$ となるように規格化する。
    \item 繰り返された添字については和をとる（ゲージ群の添字も含む）。
\end{itemize}
%------------------------------------------------
\section{標準模型とは何か}
%------------------------------------------------
\subsection{標準模型が記述するもの}
標準模型は、現在知られている素粒子と、
重力を除く3つの基本相互作用
\begin{itemize}
    \item 強い相互作用\en{strong interaction}：原子核を束縛し、クォークをハドロンに閉じ込める
    \item 弱い相互作用\en{weak interaction}：$\beta$崩壊を引き起こし、フレーバーを変える
    \item 電磁相互作用\en{electromagnetic interaction}：原子・分子・化学を支配する
\end{itemize}
を統一的に記述する、相対論的な量子場の理論\en{quantum field theory, QFT}である。

標準模型は、加速器実験の到達エネルギー
（現在 LHC で $\sqrt s=13.6\ \mathrm{TeV}$）までの範囲で、
ほぼすべての実験結果を説明する。
例えば電子の異常磁気能率\en{anomalous magnetic moment} $a_e=(g-2)/2$ は
理論と実験が10桁以上の精度で一致しており、
これは物理学全体の中でも最も精密に検証された予言の一つである。
\subsection{3つの構成原理（building principles）}
標準模型は次の3つの原理（仮定）からほぼ一意に構成される。
\begin{enumerate}
    \item \term{ゲージ対称性}{gauge symmetry}：
    理論はゲージ群
    \begin{equation}
    G_{\SM}=\SU(3)_C\times\SU(2)_L\times\U(1)_Y
    \end{equation}
    の下での局所変換\en{local transformation}に対して不変である。
    \item \term{場の内容}{field content}：
    フェルミオンは3世代あり、各世代は同一のゲージ量子数を持つ。
    スカラー場としては $\SU(2)_L$ 二重項\en{doublet}のヒッグス場\en{Higgs field}を1つだけ含む。
    \item \term{繰り込み可能性}{renormalizability}：
    ラグランジアンには質量次元\en{mass dimension}4以下の演算子のみを含める。
\end{enumerate}
これらを与えると、ラグランジアンに現れ得る項は有限個に限られ、
理論は有限個（19個）のパラメータで完全に指定される（\S\ref{sec:SM-parameters}）。
\paragraph{なぜゲージ対称性か}
ゲージ原理\en{gauge principle}の最も簡単な例は量子電磁力学\en{quantum electrodynamics, QED}である。
自由な Dirac 場のラグランジアン $\bar\psi(i\gamma^\mu\partial_\mu-m)\psi$ は
大域的な\en{global}位相変換 $\psi\to e^{i\alpha}\psi$ で不変であるが、
$\alpha=\alpha(x)$ とすると $\partial_\mu\alpha$ に比例する項が生じて不変性が失われる。
これを回復するには、変換
\begin{equation}
A_\mu\to A_\mu+\frac1e\partial_\mu\alpha
\end{equation}
に従うベクトル場 $A_\mu$ を導入し、微分を共変微分 $D_\mu=\partial_\mu-ieA_\mu$ に置き換えればよい。
このとき相互作用項 $e\bar\psi\gamma^\mu\psi A_\mu$ の形が
\textbf{対称性によって一意に決まる}。
\begin{note}
$\psi\to e^{i\alpha(x)}\psi$ のとき
$\partial_\mu\psi\to e^{i\alpha}\left(\partial_\mu\psi+i(\partial_\mu\alpha)\psi\right)$ なので
\[
\bar\psi i\gamma^\mu\partial_\mu\psi
\to
\bar\psi i\gamma^\mu\partial_\mu\psi-(\partial_\mu\alpha)\bar\psi\gamma^\mu\psi
\]
と余分な項が出る。一方、共変微分は
\[
D_\mu\psi
\to
e^{i\alpha}\Bigl(
\partial_\mu\psi+i(\partial_\mu\alpha)\psi
-ieA_\mu\psi-ie\cdot\frac1e(\partial_\mu\alpha)\psi
\Bigr)
=e^{i\alpha}D_\mu\psi
\]
と場と同じように変換するので、$\bar\psi i\gamma^\mu D_\mu\psi$ は不変である。
展開すると $\bar\psi i\gamma^\mu D_\mu\psi=\bar\psi i\gamma^\mu\partial_\mu\psi+e\bar\psi\gamma^\mu\psi A_\mu$。
非可換群への拡張は第2章で行う。
\end{note}
標準模型はこの考え方を非可換群 $\SU(3)$, $\SU(2)$ に拡張したものであり、
フェルミオンとゲージボソンの相互作用はすべてゲージ原理から導かれる。
\paragraph{なぜ繰り込み可能性か}
繰り込み可能な理論では、有限個のパラメータを実験で決めれば
すべての物理量を任意の精度で（摂動論の範囲で）予言できる。
現代的な見方では、標準模型はあるエネルギースケール $\Lambda$ 以下で有効な
\term{有効場の理論}{effective field theory, EFT}であり、
次元 $d>4$ の演算子は $(E/\Lambda)^{d-4}$ で suppress されるため、
低エネルギーでは自動的に繰り込み可能な項だけが残ると理解される。
したがって「繰り込み可能性」は、
標準模型が $\Lambda$ よりはるかに低いエネルギーで実験されていることの帰結とみなせる。
\subsection{標準模型の特徴的な性質}
\paragraph{カイラルなゲージ理論（chiral gauge theory）}
弱い相互作用は左巻きのフェルミオンにのみ作用する。
すなわち、左巻き成分と右巻き成分は
異なるゲージ表現に属する（\S\ref{sec:chirality}）。
このためパリティ\en{parity} $P$ と荷電共役\en{charge conjugation} $C$ は弱い相互作用で最大限に破れている\en{maximally violated}。
これは 1956年の Lee--Yang による提唱と、
1957年の Wu らによる ${}^{60}$Co の $\beta$崩壊実験で確立された事実である。
\paragraph{自発的対称性の破れ（spontaneous symmetry breaking, SSB）}
ゲージ対称性はフェルミオンやゲージボソンに裸の質量項\en{bare mass term}を許さない。
しかし、弱い相互作用が短距離力であることは
それを媒介するボソンが massive であることを意味する。
質量 $M$ の粒子の交換は湯川ポテンシャル\en{Yukawa potential}
\begin{equation}
V(r)\propto\frac{e^{-Mr}}{r}
\end{equation}
を与え、その到達距離\en{range}は $1/M$ である。
$M_W\simeq80\ \mathrm{GeV}$ に対しては
\begin{equation}
\frac{1}{M_W}=\frac{\hbar c}{M_Wc^2}
\simeq\frac{0.197\ \mathrm{GeV\,fm}}{80\ \mathrm{GeV}}
\simeq2.5\times10^{-3}\ \mathrm{fm}
\end{equation}
となり、陽子の大きさ（$\sim1\ \mathrm{fm}$）よりはるかに小さい。

この矛盾は、ヒッグス場が真空期待値\en{vacuum expectation value, VEV}
\begin{equation}
\langle H\rangle=\frac{1}{\sqrt2}\begin{pmatrix}0\\ v\end{pmatrix},
\qquad v\simeq 246\ \mathrm{GeV}
\end{equation}
を持つことで解決される。これにより
\begin{equation}
\SU(2)_L\times\U(1)_Y\longrightarrow\U(1)_{\mathrm{EM}}
\end{equation}
と電弱対称性\en{electroweak symmetry}が自発的に破れ、
$W^\pm$, $Z^0$ ボソンとフェルミオンが質量を獲得する（第II部）。
ラグランジアンは対称性を保ったまま、真空\en{vacuum}だけが対称性を破るため、
理論の繰り込み可能性は保たれる（'t Hooft, 1971）。
\paragraph{閉じ込めと漸近的自由性（confinement and asymptotic freedom）}
$\SU(3)_C$ の結合定数\en{coupling constant}は高エネルギーで小さく（漸近的自由性）、
低エネルギー $\Lambda_{\mathrm{QCD}}\sim200$--$300\ \mathrm{MeV}$ 付近で強くなる。
その結果、クォークとグルーオンは単独では観測されず、
カラー一重項\en{color singlet}のハドロン（陽子・中性子・中間子など）に閉じ込められる。
陽子の質量 $938\ \mathrm{MeV}$ の大部分は、
$u,d$ クォークの質量（数 MeV）ではなく QCD の束縛エネルギー\en{binding energy}に由来する。
\subsection{エネルギースケールと相互作用の強さ（energy scales and interaction strengths）}
標準模型にはいくつかの特徴的なエネルギースケールがある。
\begin{center}
\renewcommand{\arraystretch}{1.3}
\begin{tabular}{lcl}
\toprule
scale & 値 & 意味\\
\midrule
$\Lambda_{\mathrm{QCD}}$ & $\sim0.2\ \mathrm{GeV}$ & QCD が強結合になるスケール\\
$M_W,\ M_Z$ & $80,\ 91\ \mathrm{GeV}$ & 弱ボソンの質量\\
$v$ & $246\ \mathrm{GeV}$ & 電弱対称性の破れのスケール\\
$m_t$ & $173\ \mathrm{GeV}$ & 最も重い素粒子\\
$M_{\mathrm{Pl}}$ & $1.2\times10^{19}\ \mathrm{GeV}$ & 重力が強くなるスケール（標準模型の外）\\
\bottomrule
\end{tabular}
\end{center}
$Z$ ボソンの質量付近での各相互作用の強さは、おおよそ
\begin{equation}
\alpha_s(M_Z)\simeq0.118,\qquad
\alpha(M_Z)\simeq\frac{1}{128},\qquad
\alpha_W\equiv\frac{g^2}{4\pi}\simeq\frac{1}{30}
\end{equation}
である。弱い相互作用が低エネルギーで「弱い」のは
結合定数 $g$ が小さいからではなく（実際 $g>e$ である）、
$W$ ボソンの伝播関数\en{propagator}が
\begin{equation}
\frac{1}{q^2-M_W^2}\xrightarrow{|q^2|\ll M_W^2}-\frac{1}{M_W^2}
\end{equation}
によって suppress されるからである。
この suppression は Fermi 定数\en{Fermi constant}
\begin{equation}
\frac{G_F}{\sqrt2}=\frac{g^2}{8M_W^2}=\frac{1}{2v^2},
\qquad
G_F\simeq1.166\times10^{-5}\ \mathrm{GeV}^{-2}
\label{eq:GF-preview}
\end{equation}
にまとめられる（Fermi 理論との対応は第8章で導出する）。
\begin{note}
ここでは $M_W=gv/2$（第6章で導出）を認めて代入するだけである：
$\dfrac{g^2}{8M_W^2}=\dfrac{g^2}{8\cdot g^2v^2/4}=\dfrac{1}{2v^2}$。
したがって $G_F=\dfrac{1}{\sqrt2v^2}$。
数値的には $\sqrt2G_F=1.650\times10^{-5}\ \mathrm{GeV}^{-2}$ より
$v=(1.650\times10^{-5})^{-1/2}\ \mathrm{GeV}\simeq246.2\ \mathrm{GeV}$。
\end{note}
逆に解けば
\begin{equation}
v=(\sqrt2G_F)^{-1/2}\simeq246\ \mathrm{GeV}
\end{equation}
となり、ミューオンの寿命\en{muon lifetime}の測定から電弱スケール\en{electroweak scale}が決まる。
\begin{remark}
標準模型は 1960--70年代に
Glashow (1961), Weinberg (1967), Salam (1968) による電弱理論と、
Gross--Wilczek, Politzer (1973) による QCD の漸近的自由性の発見を経て確立された。
小林・益川 (1973) は CP の破れ\en{CP violation}を説明するために第3世代の存在を予言した。
その後、中性カレント\en{neutral current}の発見（1973年）、
$W^\pm$, $Z^0$ の発見（1983年）、
トップクォークの発見（1995年）、
タウニュートリノの直接観測（2000年）、
ヒッグス粒子の発見（2012年）により、
理論が予言する粒子はすべて観測されている。
\end{remark}
\subsection{標準模型が説明しないもの（beyond the SM）}
標準模型は極めて成功した理論であるが、次のような現象は説明できない（第14章）。
\begin{itemize}
    \item 重力\en{gravity}（一般相対論との統合）
    \item ニュートリノ質量\en{neutrino mass}（振動実験により確立）
    \item 暗黒物質と暗黒エネルギー\en{dark matter, dark energy}
    \item 宇宙のバリオン数非対称性\en{baryon asymmetry of the universe}
    \item フェルミオン質量の階層性と世代数が3である理由
    \item ヒッグス質量の微調整問題\en{fine-tuning}（階層性問題\en{hierarchy problem}）と強い CP 問題
\end{itemize}
%------------------------------------------------
\section{素粒子の分類：フェルミオンとボソン}
%------------------------------------------------
\subsection{スピンと統計（spin and statistics）}
スピン統計定理\en{spin-statistics theorem}により、粒子は
\begin{itemize}
    \item 半整数スピンを持ち Fermi--Dirac 統計に従う\textbf{フェルミオン}
    \item 整数スピンを持ち Bose--Einstein 統計に従う\textbf{ボソン}
\end{itemize}
に分類される。
場の理論では、フェルミオン場は反交換関係\en{anticommutation relation}で、
ボソン場は交換関係で量子化される。
フェルミオンは Pauli の排他原理\en{Pauli exclusion principle}に従うため同じ状態に1つしか入れず、
これが原子の殻構造や物質の安定性を支えている。
標準模型において、フェルミオンは物質を構成し、
ボソンは相互作用を媒介する（あるいは質量を与える）役割を担う。

すべての粒子には、質量とスピンが同じで
すべての内部量子数（電荷・カラー・フレーバー）が反対の\term{反粒子}{antiparticle}が存在する。
光子、$Z^0$、ヒッグス粒子は自分自身が反粒子である。
\subsection{フェルミオン}
標準模型のフェルミオンはすべてスピン $1/2$ であり、
強い相互作用を受ける\term{クォーク}{quark}と、
受けない\term{レプトン}{lepton}に分けられる。
これらは同じゲージ量子数を持つ3つの\term{世代}{generation}を構成し、
世代間の違いは質量（Yukawa 結合\en{Yukawa coupling}）のみである。
\begin{center}
\renewcommand{\arraystretch}{1.3}
\begin{tabular}{ccccc}
\toprule
 & $Q$ & $i=1$ & $i=2$ & $i=3$\\
\midrule
$u_i$ (up-type) & $+2/3$ & $u$ \ ($2.2$ MeV) & $c$ \ ($1.27$ GeV) & $t$ \ ($173$ GeV)\\
$d_i$ (down-type) & $-1/3$ & $d$ \ ($4.7$ MeV) & $s$ \ ($93$ MeV) & $b$ \ ($4.18$ GeV)\\
\midrule
$e_i$ (charged lepton) & $-1$ & $e$ \ ($0.511$ MeV) & $\mu$ \ ($106$ MeV) & $\tau$ \ ($1.78$ GeV)\\
$\nu_i$ (neutrino) & $0$ & $\nu_e$ & $\nu_\mu$ & $\nu_\tau$\\
\bottomrule
\end{tabular}
\end{center}
ここで $i=1,2,3$ は世代の添字であり、
$u_i=(u,c,t)$, $d_i=(d,s,b)$, $e_i=(e,\mu,\tau)$, $\nu_i=(\nu_e,\nu_\mu,\nu_\tau)$ と書く
（荷電レプトンを $\ell_i$ と書くことも多い）。
括弧内は質量の概算値である
（軽いクォークの質量は $\overline{\mathrm{MS}}$ スキームでの値であり、
クォークは閉じ込められているため直接測定できない）。
\paragraph{質量の階層（mass hierarchy）}
質量は電子の $0.5\ \mathrm{MeV}$ からトップの $173\ \mathrm{GeV}$ まで
5桁以上にわたって分布する。
後で見るように、フェルミオン質量は Yukawa 結合 $y_f$ によって
$m_f=y_fv/\sqrt2$ と与えられるので、
\begin{equation}
y_t=\frac{\sqrt2m_t}{v}\simeq0.99,
\qquad
y_e=\frac{\sqrt2m_e}{v}\simeq2.9\times10^{-6}
\end{equation}
となる。
\begin{note}
$y_t=\sqrt2\times173/246\simeq1.414\times0.703\simeq0.99$、
$y_e=\sqrt2\times0.000511/246\simeq2.94\times10^{-6}$。
$m_f=y_fv/\sqrt2$ の導出は \S\ref{sec:SM-lagrangian} の Yukawa 項の節と第7章で行う。
\end{note}
トップクォークだけが $O(1)$ の結合を持ち、
他はすべて小さい。この階層構造は標準模型の中では説明されない
（\term{フレーバー問題}{flavor problem}）。
ニュートリノの質量は $\lesssim0.1\ \mathrm{eV}$ とさらに小さく、
最小標準模型\en{minimal SM}では massless である。
\paragraph{クォーク}
クォークはさらに3つの\term{カラー}{color}自由度（red, green, blue）を持つ。
カラーの数 $N_c=3$ は、例えば
\begin{equation}
R\equiv\frac{\sigma(e^+e^-\to\text{ハドロン})}{\sigma(e^+e^-\to\mu^+\mu^-)}
\simeq N_c\sum_qQ_q^2
\end{equation}
の測定によって確かめられる。
$u,d,s,c,b$ が寄与するエネルギー領域では
$N_c\sum_qQ_q^2=3\times(4/9+1/9+1/9+4/9+1/9)=11/3$ となり、
実験値とよく一致する（$N_c=1$ では合わない）。

クォークは閉じ込めにより単独では現れず、
クォーク・反クォーク対からなる\term{中間子}{meson}（$\pi^+=u\bar d$ など）や、
3つのクォークからなる\term{バリオン}{baryon}（$p=uud$, $n=udd$ など）として観測される。
ただしトップクォークは寿命（$\sim5\times10^{-25}\ \mathrm{s}$）が
ハドロン化\en{hadronization}の時間スケール $1/\Lambda_{\mathrm{QCD}}$ より短いため、
ハドロンを作る前に崩壊する。
\paragraph{レプトン}
荷電レプトン $e,\mu,\tau$ は電磁相互作用と弱い相互作用を受け、
ニュートリノは弱い相互作用のみを受ける。
ニュートリノの世代数は、$Z$ ボソンの崩壊幅\en{decay width}のうち
観測されない部分（invisible width, $Z\to\nu\bar\nu$）から
\begin{equation}
N_\nu=2.984\pm0.008
\end{equation}
と決定されている（LEP）。
したがって、$m_\nu<M_Z/2$ の軽い\en{light}ニュートリノを持つ世代は
ちょうど3つである。
\paragraph{1世代の自由度}
各世代はクォーク $(u,d)\times3$ カラーとレプトン $(\nu,e)$ からなる。
後で見るように、これはちょうど15個の Weyl スピノルに相当する（\S\ref{sec:chirality}）。
\subsection{ボソン}
\begin{center}
\renewcommand{\arraystretch}{1.3}
\begin{tabular}{ccccc}
\toprule
 & spin & $Q$ & mass & 役割\\
\midrule
$g$ (gluon) & 1 & 0 & 0 & 強い相互作用（8種類）\\
$\gamma$ (photon) & 1 & 0 & 0 & 電磁相互作用\\
$W^\pm$ & 1 & $\pm1$ & $80.4$ GeV & 弱い相互作用（charged current）\\
$Z^0$ & 1 & 0 & $91.2$ GeV & 弱い相互作用（neutral current）\\
$h$ (Higgs) & 0 & 0 & $125$ GeV & 質量の生成\\
\bottomrule
\end{tabular}
\end{center}
スピン1のボソンは\term{ゲージボソン}{gauge boson}であり、
ゲージ対称性からその存在と相互作用の形が決まる。
対称性の破れの前のゲージ場
\begin{equation}
G_\mu^a\ (a=1,\dots,8),\qquad
W_\mu^i\ (i=1,2,3),\qquad
B_\mu
\end{equation}
のうち、$W_\mu^i$ と $B_\mu$ は電弱対称性の破れによって
物理的な場 $W_\mu^\pm$, $Z_\mu$, $A_\mu$（光子）に再編成される（第6章）：
\begin{equation}
W_\mu^\pm=\frac{1}{\sqrt2}\left(W_\mu^1\mp iW_\mu^2\right),
\qquad
\begin{pmatrix}Z_\mu\\ A_\mu\end{pmatrix}
=
\begin{pmatrix}
\cos\theta_W & -\sin\theta_W\\
\sin\theta_W & \cos\theta_W
\end{pmatrix}
\begin{pmatrix}W_\mu^3\\ B_\mu\end{pmatrix}.
\label{eq:ZA-preview}
\end{equation}
ここで $\theta_W$ は Weinberg 角\en{weak mixing angle}であり、
\begin{equation}
\cos\theta_W=\frac{g}{\sqrt{g^2+g'^2}},\qquad
\sin\theta_W=\frac{g'}{\sqrt{g^2+g'^2}},\qquad
\sin^2\theta_W\simeq0.23
\end{equation}
で与えられる。
\paragraph{ゲージボソンの偏極自由度（polarization degrees of freedom）}
massless なスピン1粒子は横波\en{transverse}の2つの偏極（ヘリシティ $\pm1$）しか持たないが、
massive なスピン1粒子は縦波\en{longitudinal}を含む3つの偏極を持つ。
対称性の破れによって $W^\pm$, $Z^0$ が massive になると、
各ボソンにつき1つの自由度が増える。
この増えた3つの自由度は、ヒッグス二重項の4つの実自由度のうちの3つ
（Nambu--Goldstone ボソン）が「食べられた」\en{eaten}ものである。
自由度の数え上げは
\begin{equation}
\underbrace{4\times2}_{W^i,B\ \text{(massless)}}
+\underbrace{4}_{H}
=12
=
\underbrace{3\times3}_{W^\pm,Z\ \text{(massive)}}
+\underbrace{2}_{\gamma}
+\underbrace{1}_{h}
\end{equation}
となり、対称性の破れの前後で一致する。
残った1つの実自由度が、物理的なスピン0粒子であるヒッグス粒子 $h$ である。
\subsection{相互作用の一覧（list of interactions）}
どの粒子がどのボソンと結合するかを、まとめておく。
\begin{center}
\renewcommand{\arraystretch}{1.3}
\begin{tabular}{lcccc}
\toprule
 & $g$ & $\gamma$ & $W^\pm$ & $Z^0$\\
\midrule
$u_L,\ d_L$ & ○ & ○ & ○ & ○\\
$u_R,\ d_R$ & ○ & ○ & × & ○\\
$e_L$ & × & ○ & ○ & ○\\
$e_R$ & × & ○ & × & ○\\
$\nu_L$ & × & × & ○ & ○\\
\bottomrule
\end{tabular}
\end{center}
$W^\pm$ は左巻きフェルミオンのみと結合し、
二重項の上成分と下成分を入れ替える（$e^-\leftrightarrow\nu_e$, $d\leftrightarrow u$ など）。
これがフレーバーを変える\en{flavor-changing}唯一の相互作用である。
$Z^0$ と光子はフレーバーを変えない（第10章）。
またヒッグス粒子は質量に比例する強さでフェルミオンや弱ボソンと結合する。

非可換ゲージ理論であるため、ゲージボソン同士も結合する。
グルーオンは $ggg$, $gggg$ の自己相互作用\en{self-interaction}を持ち、
電弱部分には $WW\gamma$, $WWZ$ および4点結合が存在する。
光子同士は直接結合しない。
%------------------------------------------------
\section{標準模型のゲージ群}
%------------------------------------------------
標準模型のゲージ群は
\begin{equation}
\boxed{
G_{\SM}=\SU(3)_C\times\SU(2)_L\times\U(1)_Y
}
\end{equation}
である。
各因子はそれぞれ独立な結合定数
\begin{equation}
g_s,\qquad g,\qquad g'
\end{equation}
を持つ。ゲージ場の数は各群の次元（生成子の数）に等しい。
\subsection{Lie 群の基礎事項（basics of Lie groups）}
\begin{mydef}{特殊ユニタリ群 $\SU(N)$}{SUN}
\begin{equation}
\SU(N)=\{U\in\mathrm{GL}(N,\mathbb{C})\mid U^\dagger U=1,\ \det U=1\}.
\end{equation}
単位元の近くの元は、エルミートかつトレースレスな生成子 $T^a$ を用いて
\begin{equation}
U=\exp\left(i\alpha^aT^a\right),\qquad \alpha^a\in\mathbb{R}
\end{equation}
と書ける。
\end{mydef}
\begin{note}
$U=e^{iX}$（$X=\alpha^aT^a$）とすると $U^\dagger=e^{-iX^\dagger}$ なので、
$U^\dagger U=1$ が任意の $\alpha^a$ で成り立つには $X^\dagger=X$、すなわち $T^a$ がエルミート。
また $\det e^{iX}=e^{i\tr X}$ なので、$\det U=1$ には $\tr T^a=0$ が必要。
\end{note}
$N\times N$ のエルミートかつトレースレスな行列の実自由度は $N^2-1$ なので
\begin{equation}
\dim\SU(N)=N^2-1
\end{equation}
である。生成子は Lie 代数
\begin{equation}
[T^a,T^b]=if^{abc}T^c
\end{equation}
を満たし、$f^{abc}$ を\term{構造定数}{structure constant}と呼ぶ。
規格化を $\tr(T^aT^b)=\frac12\delta^{ab}$（基本表現）とすると $f^{abc}$ は完全反対称になる。

標準模型で現れる表現は次の3種類である。
\begin{itemize}
    \item \term{自明表現}{trivial representation} $\mathbf{1}$：$T^a=0$。場は変換しない。
    \item \term{基本表現}{fundamental representation} $\mathbf{N}$：$N$ 成分の列ベクトルに $U$ が作用する。
    \item \term{随伴表現}{adjoint representation}：$(T^a_{\mathrm{adj}})^{bc}=-if^{abc}$。$N^2-1$ 次元で、ゲージ場が属する。
\end{itemize}
表現 $R$ の2次 Casimir 演算子\en{quadratic Casimir} $C_2(R)$ と Dynkin 指数\en{Dynkin index} $T(R)$ を
\begin{equation}
T^a_RT^a_R=C_2(R)\,\mathbf{1},\qquad
\tr(T^a_RT^b_R)=T(R)\,\delta^{ab}
\end{equation}
で定義すると、$\SU(N)$ では
\begin{equation}
C_2(\mathbf{N})=\frac{N^2-1}{2N},\qquad
C_2(\mathrm{adj})=N,\qquad
T(\mathbf{N})=\frac12,\qquad
T(\mathrm{adj})=N
\end{equation}
である。QCD では $C_F\equiv C_2(\mathbf{3})=4/3$, $C_A\equiv C_2(\mathbf{8})=3$ となる。
\begin{note}
$\sum_a\tr(T^a_RT^a_R)$ を2通りに計算する。
左の定義から $\tr(C_2(R)\mathbf 1)=C_2(R)\dim R$、
右の定義から $\sum_aT(R)\delta^{aa}=T(R)\dim G$。よって
\[
C_2(R)=\frac{T(R)\dim G}{\dim R}.
\]
基本表現では $T=1/2$, $\dim G=N^2-1$, $\dim R=N$ より $C_2(\mathbf N)=\frac{N^2-1}{2N}$。
$\SU(2)$ では $C_2(\mathbf2)=3/4=j(j+1)|_{j=1/2}$ となり、角運動量の結果と一致する。
随伴表現では $\dim R=\dim G$ なので $C_2(\mathrm{adj})=T(\mathrm{adj})=N$
（$T(\mathrm{adj})=N$ は $f^{acd}f^{bcd}=N\delta^{ab}$ による）。
\end{note}
これらはループ計算で頻繁に現れる（付録A）。
\subsection{$\SU(3)_C$：強い相互作用}
$\dim\SU(3)=8$ であり、生成子は Gell-Mann 行列 $\lambda^a$ を用いて $T^a=\lambda^a/2$ と書かれる：
\begin{gather}
\lambda^1=\begin{pmatrix}0&1&0\\1&0&0\\0&0&0\end{pmatrix},\quad
\lambda^2=\begin{pmatrix}0&-i&0\\i&0&0\\0&0&0\end{pmatrix},\quad
\lambda^3=\begin{pmatrix}1&0&0\\0&-1&0\\0&0&0\end{pmatrix},\quad
\lambda^4=\begin{pmatrix}0&0&1\\0&0&0\\1&0&0\end{pmatrix},\\
\lambda^5=\begin{pmatrix}0&0&-i\\0&0&0\\i&0&0\end{pmatrix},\quad
\lambda^6=\begin{pmatrix}0&0&0\\0&0&1\\0&1&0\end{pmatrix},\quad
\lambda^7=\begin{pmatrix}0&0&0\\0&0&-i\\0&i&0\end{pmatrix},\quad
\lambda^8=\frac{1}{\sqrt3}\begin{pmatrix}1&0&0\\0&1&0\\0&0&-2\end{pmatrix}.
\end{gather}
これらは $\tr(\lambda^a\lambda^b)=2\delta^{ab}$ を満たす。
$\lambda^3$ と $\lambda^8$ は対角的で互いに可換であり、
$\SU(3)$ の Cartan 部分代数\en{Cartan subalgebra}（階数\en{rank}2）を張る。
構造定数の非ゼロ成分は（添字の置換を除いて）
\begin{equation}
f^{123}=1,\quad
f^{147}=f^{246}=f^{257}=f^{345}=\frac12,\quad
f^{156}=f^{367}=-\frac12,\quad
f^{458}=f^{678}=\frac{\sqrt3}{2}
\end{equation}
である。

8個のゲージ場がグルーオンであり、
この部分の理論が量子色力学（Quantum Chromodynamics, QCD）である。
クォークは基本表現 $\mathbf{3}$（反クォークは $\bar{\mathbf{3}}$）に、
レプトンは自明表現 $\mathbf{1}$ に属する。
$\SU(3)_C$ は対称性の破れを起こさず\en{unbroken}、グルーオンは massless のままである。

カラーの観点では、ハドロンがカラー一重項であることは
\begin{equation}
\mathbf{3}\otimes\bar{\mathbf{3}}=\mathbf{1}\oplus\mathbf{8},
\qquad
\mathbf{3}\otimes\mathbf{3}\otimes\mathbf{3}=\mathbf{1}\oplus\mathbf{8}\oplus\mathbf{8}\oplus\mathbf{10}
\end{equation}
にそれぞれ一重項 $\mathbf{1}$ が含まれることに対応する。
中間子 $q\bar q$ とバリオン $qqq$ はこの一重項の状態である。
\subsection{$\SU(2)_L$：弱アイソスピン}
$\SU(2)_L$ の生成子は Pauli 行列
\begin{equation}
\sigma^1=\begin{pmatrix}0&1\\1&0\end{pmatrix},\qquad
\sigma^2=\begin{pmatrix}0&-i\\i&0\end{pmatrix},\qquad
\sigma^3=\begin{pmatrix}1&0\\0&-1\end{pmatrix}
\end{equation}
を用いて
\begin{equation}
T^i=\frac{\sigma^i}{2},
\qquad
[T^i,T^j]=i\epsilon^{ijk}T^k
\qquad(i,j,k=1,2,3)
\end{equation}
と表される。$\dim\SU(2)=3$ なので3個のゲージ場 $W_\mu^i$ が存在する。
この群の量子数を\term{弱アイソスピン}{weak isospin}と呼び、
二重項 $\mathbf{2}$ の上成分は $T_3=+1/2$、下成分は $T_3=-1/2$ を持つ。

昇降演算子\en{raising/lowering operator} $T^\pm=T^1\pm iT^2$ は
\begin{equation}
T^+=\begin{pmatrix}0&1\\0&0\end{pmatrix},\qquad
T^-=\begin{pmatrix}0&0\\1&0\end{pmatrix}
\end{equation}
であり、二重項の成分を入れ替える。
\begin{equation}
W_\mu^1T^1+W_\mu^2T^2
=\frac{1}{\sqrt2}\left(W_\mu^+T^++W_\mu^-T^-\right)
\end{equation}
が成り立つ。
\begin{note}
$W^\pm=(W^1\mp iW^2)/\sqrt2$ と $T^\pm=T^1\pm iT^2$ を代入すると
\[
W^+T^++W^-T^-
=\tfrac{1}{\sqrt2}\left[(W^1-iW^2)(T^1+iT^2)+(W^1+iW^2)(T^1-iT^2)\right]
=\sqrt2\left(W^1T^1+W^2T^2\right).
\]
交差項 $\pm i(W^1T^2-W^2T^1)$ は打ち消し合う。
\end{note}
したがって $W^\pm$ ボソンとの結合は $T^\pm$ を通じて
二重項の上下を入れ替える。これが荷電カレント\en{charged current, CC}相互作用である（第8章）。

添字 $L$ は、この群が\textbf{左巻き}フェルミオンにのみ作用することを表す。
左巻きフェルミオンは二重項 $\mathbf{2}$ に、
右巻きフェルミオンは一重項 $\mathbf{1}$ に属する。
\begin{myprop}{$\SU(2)$ の基本表現の擬実性}{pseudoreal}
Pauli 行列は
\begin{equation}
\sigma^2(\sigma^i)^*\sigma^2=-\sigma^i
\label{eq:pseudoreal}
\end{equation}
を満たす。したがって $\psi\to U\psi$（$U\in\SU(2)$）のとき
\begin{equation}
i\sigma^2\psi^*\to U\,(i\sigma^2\psi^*)
\end{equation}
となり、$\mathbf{2}$ とその複素共役表現 $\bar{\mathbf{2}}$ は同値である（$\mathbf{2}$ は\term{擬実}{pseudoreal}表現）。
\end{myprop}
\begin{note}
$\sigma^1,\sigma^3$ は実なので $\sigma^2\sigma^{1,3}\sigma^2=-\sigma^{1,3}$（$\sigma^2$ と反可換）。
$\sigma^2$ は純虚なので $\sigma^2(\sigma^2)^*\sigma^2=-\sigma^2\sigma^2\sigma^2=-\sigma^2$。

次に $U=e^{i\beta^i\sigma^i/2}$ に対して
\[
\sigma^2U^*\sigma^2
=\sigma^2e^{-i\beta^i(\sigma^i)^*/2}\sigma^2
=e^{-i\beta^i\sigma^2(\sigma^i)^*\sigma^2/2}
=e^{i\beta^i\sigma^i/2}=U .
\]
$(\sigma^2)^2=1$ を挟むと
$i\sigma^2(U\psi)^*=i\sigma^2U^*\sigma^2\sigma^2\psi^*=U(i\sigma^2\psi^*)$。
\end{note}
これは後でヒッグス場の共役 $\tilde H=i\sigma^2H^*$ を作るときに用いる。
\subsection{$\U(1)_Y$：弱超電荷}
$\U(1)_Y$ は可換群であり、ゲージ場は $B_\mu$ の1つである。
対応する量子数 $Y$ を\term{弱超電荷}{weak hypercharge}と呼ぶ。
$\U(1)$ の既約表現はすべて1次元であり、
超電荷 $Y$ の場は $\psi\to e^{i\theta Y/2}\psi$ と変換する。

電荷 $Q$ は弱アイソスピンの第3成分 $T_3$ と $Y$ によって
\begin{equation}
\boxed{
Q=T_3+\frac{Y}{2}
}
\label{eq:GNN}
\end{equation}
と表される（Gell-Mann--西島--中野の関係式）。
これは、電弱対称性の破れの後に残る $\U(1)_{\mathrm{EM}}$ の生成子が
$\SU(2)_L$ と $\U(1)_Y$ の生成子の線形結合であることを意味する。
実際、\eqref{eq:ZA-preview} を共変微分に代入すると、
光子 $A_\mu$ との結合は
\begin{equation}
g\sin\theta_W\,T_3+g'\cos\theta_W\frac{Y}{2}
=e\left(T_3+\frac Y2\right)=eQ
\end{equation}
となる。ここで電気素量\en{elementary charge}は
\begin{equation}
e=g\sin\theta_W=g'\cos\theta_W
=\frac{gg'}{\sqrt{g^2+g'^2}}
\label{eq:e-preview}
\end{equation}
で与えられる（第6章）。
\begin{note}
\eqref{eq:ZA-preview} を逆に解くと
$W^3_\mu=\cos\theta_WZ_\mu+\sin\theta_WA_\mu$,
$B_\mu=-\sin\theta_WZ_\mu+\cos\theta_WA_\mu$。
これを共変微分の中性部分 $gW^3_\mu T_3+g'B_\mu\frac Y2$ に代入して $A_\mu$ の係数を取り出すと
\[
g\sin\theta_W\,T_3+g'\cos\theta_W\,\frac Y2
=\frac{gg'}{\sqrt{g^2+g'^2}}\left(T_3+\frac Y2\right).
\]
$T_3$ と $Y/2$ の係数が等しくなるのは、
$\theta_W$ が $\tan\theta_W=g'/g$ で定まっているからである。
$Z_\mu$ の係数については第6章・第8章で扱う。
\end{note}
可換群の電荷の規格化は群論からは決まらないので、
$Y$ の値は理論の入力である。
ただし、後で見るゲージ異常の相殺条件（第3章）により、
1世代内の超電荷の比はほぼ一意に決まる。
その結果、電荷の量子化\en{charge quantization}（陽子と電子の電荷が厳密に逆符号で等しいこと）が説明される。
\subsection{ゲージ変換（gauge transformations）}
ゲージ群の元を
\begin{equation}
U(x)=
\exp\left(i\alpha^a(x)T^a\right)
\exp\left(i\beta^i(x)\frac{\sigma^i}{2}\right)
\exp\left(i\theta(x)\frac{Y}{2}\right)
\end{equation}
と書くと、物質場は $\psi\to U\psi$（ただし $\psi$ の表現に応じた部分だけが作用する）と変換する。
例えば
\begin{align}
Q_L&\to
e^{i\alpha^aT^a}\,e^{i\beta^i\sigma^i/2}\,e^{i\theta/6}\,Q_L,\\
u_R&\to e^{i\alpha^aT^a}\,e^{2i\theta/3}\,u_R,\\
L_L&\to e^{i\beta^i\sigma^i/2}\,e^{-i\theta/2}\,L_L,\\
e_R&\to e^{-i\theta}\,e_R
\end{align}
である。ゲージ場は、各群について
\begin{equation}
A_\mu\equiv A_\mu^aT^a\to
UA_\mu U^{-1}-\frac ig(\partial_\mu U)U^{-1}
\end{equation}
と変換する。$\U(1)_Y$ に対しては
$B_\mu\to B_\mu+\frac{1}{g'}\partial_\mu\theta$ である。
これにより $D_\mu\psi\to UD_\mu\psi$ が成り立つ（非可換の場合の証明は第2章）。
\begin{note}
例えば $L_L$ の $\U(1)_Y$ 部分を確かめる。$L_L\to e^{-i\theta/2}L_L$,
$B_\mu\to B_\mu+\frac{1}{g'}\partial_\mu\theta$ のとき
\[
\Bigl(\partial_\mu+\tfrac i2g'B_\mu\Bigr)L_L
\to
e^{-i\theta/2}\Bigl(
\partial_\mu L_L-\tfrac i2(\partial_\mu\theta)L_L
+\tfrac i2g'B_\mu L_L+\tfrac i2(\partial_\mu\theta)L_L
\Bigr)
=e^{-i\theta/2}\Bigl(\partial_\mu+\tfrac i2g'B_\mu\Bigr)L_L .
\]
超電荷 $Y$ の場に対しても、$\partial_\mu\theta$ の項は常に $\pm\frac{iY}{2}\partial_\mu\theta$ の形で打ち消す。
\end{note}
\subsection{ゲージ結合定数のランニング（running couplings）}
量子補正\en{quantum correction}により、結合定数はエネルギースケール $\mu$ に依存する（繰り込み群\en{renormalization group}）。
1ループでは
\begin{equation}
\mu\frac{dg_i}{d\mu}=\frac{b_i}{16\pi^2}g_i^3,
\qquad\text{すなわち}\qquad
\frac{1}{\alpha_i(\mu)}=\frac{1}{\alpha_i(\mu_0)}-\frac{b_i}{2\pi}\ln\frac{\mu}{\mu_0}
\end{equation}
であり、$\SU(N)$ ゲージ理論で $n_f$ 個の Dirac フェルミオン（基本表現）と
$n_s$ 個の複素スカラー（基本表現）がある場合
\begin{equation}
b=-\frac{11}{3}N+\frac23n_f+\frac16n_s
\end{equation}
となる（$n_f$ は Weyl フェルミオンでは $1/2$ ずつ数える）。
標準模型では
\begin{equation}
b_3=-11+\frac23\times6=-7,\qquad
b_2=-\frac{22}{3}+\frac43n_g+\frac16=-\frac{19}{6},\qquad
b_Y=+\frac{41}{6}
\end{equation}
である（$n_g=3$ は世代数、$b_Y$ は $Q=T_3+Y/2$ の $g'$ に対する値）。
$\beta$ 関数の導出は1ループの真空偏極の計算を要するので本ノートでは扱わず、結果のみ用いる。
\begin{note}
\textbf{$b_3$}：クォークは6フレーバーの Dirac 場（各々 $\mathbf 3$）なので $n_f=6$。

\textbf{$b_2$}：1世代あたりの $\SU(2)_L$ 二重項の Weyl 場は $Q_L$（3カラー）と $L_L$ の4個で、
Dirac 換算で $n_f=2$。3世代で $n_f=6$、すなわち $\frac23\times6=4=\frac43n_g$。
ヒッグス二重項は複素スカラー1個で $\frac16$。

\textbf{$b_Y$}：$\U(1)$ では $T(R)$ の代わりに電荷の2乗 $(Y/2)^2$ を用い、
Weyl 場1個あたり $\frac23(Y/2)^2$、複素スカラー1個あたり $\frac13(Y/2)^2$ が寄与する。
1世代の Weyl 場について
\[
\sum\left(\frac Y2\right)^2
=6\cdot\frac1{36}+3\cdot\frac49+3\cdot\frac19+2\cdot\frac14+1\cdot1
=\frac16+\frac43+\frac13+\frac12+1=\frac{10}{3}.
\]
したがって $b_Y=\frac23\cdot\frac{10}{3}\cdot3+\frac13\cdot2\cdot\frac14=\frac{20}{3}+\frac16=\frac{41}{6}$。
大統一理論の規格化 $g_1=\sqrt{5/3}\,g'$ では $b_1=\frac35b_Y=\frac{41}{10}$ となる。
\end{note}
$b_3<0$ は\term{漸近的自由性}{asymptotic freedom}を意味する。
非可換ゲージ場の自己相互作用に由来する $-\frac{11}{3}N$ が
フェルミオンの寄与に勝るためである。
一方 $\U(1)_Y$ の結合は高エネルギーで大きくなる。
3つの結合定数が高エネルギーで近づくことは
大統一理論\en{grand unified theory, GUT}の動機となる（第14章）。
\begin{remark}
ゲージ場の総数は $8+3+1=12$ である。
対称性の破れの後、そのうち3個（$W^\pm$, $Z^0$）が massive になり、
9個（グルーオン8個と光子）が massless のまま残る。
\end{remark}
\begin{remark}
厳密には、標準模型の全ての場に自明に作用する
離散部分群 $\mathbb{Z}_6$ が存在するため、
ゲージ群の大域的構造は
$(\SU(3)\times\SU(2)\times\U(1))/\mathbb{Z}_6$
とも書ける。摂動論\en{perturbation theory}の範囲ではこの違いは現れない。
\end{remark}
%------------------------------------------------
\section{粒子の表現と量子数}
%------------------------------------------------
\subsection{1世代の量子数（quantum numbers of one generation）}
標準模型の各場は、ゲージ群の各因子の表現の組
\begin{equation}
(R_3,\ R_2)_Y
\end{equation}
で指定される。
ここで $R_3$ は $\SU(3)_C$ の表現（$\mathbf{3}$ または $\mathbf{1}$）、
$R_2$ は $\SU(2)_L$ の表現（$\mathbf{2}$ または $\mathbf{1}$）、
$Y$ は弱超電荷である。

1世代分の場の量子数を以下にまとめる。
左巻き二重項を
\begin{equation}
Q_L=\begin{pmatrix}u_L\\ d_L\end{pmatrix},
\qquad
L_L=\begin{pmatrix}\nu_L\\ e_L\end{pmatrix}
\end{equation}
と書く。
\begin{center}
\renewcommand{\arraystretch}{1.35}
\begin{tabular}{cccccc}
\toprule
場 & $\SU(3)_C$ & $\SU(2)_L$ & $Y$ & $T_3$ & $Q=T_3+Y/2$\\
\midrule
$Q_L$ & $\mathbf{3}$ & $\mathbf{2}$ & $+1/3$ &
$\begin{matrix}+1/2\\ -1/2\end{matrix}$ &
$\begin{matrix}+2/3\ (u_L)\\ -1/3\ (d_L)\end{matrix}$\\
$u_R$ & $\mathbf{3}$ & $\mathbf{1}$ & $+4/3$ & $0$ & $+2/3$\\
$d_R$ & $\mathbf{3}$ & $\mathbf{1}$ & $-2/3$ & $0$ & $-1/3$\\
\midrule
$L_L$ & $\mathbf{1}$ & $\mathbf{2}$ & $-1$ &
$\begin{matrix}+1/2\\ -1/2\end{matrix}$ &
$\begin{matrix}0\ (\nu_L)\\ -1\ (e_L)\end{matrix}$\\
$e_R$ & $\mathbf{1}$ & $\mathbf{1}$ & $-2$ & $0$ & $-1$\\
\midrule
$H$ & $\mathbf{1}$ & $\mathbf{2}$ & $+1$ &
$\begin{matrix}+1/2\\ -1/2\end{matrix}$ &
$\begin{matrix}+1\ (H^+)\\ 0\ (H^0)\end{matrix}$\\
\bottomrule
\end{tabular}
\end{center}
第2・第3世代は、$(u,d,\nu_e,e)$ を
$(c,s,\nu_\mu,\mu)$, $(t,b,\nu_\tau,\tau)$ に置き換えたものである。
\begin{myex}{$u$ クォークの電荷}{u-charge}
$Q_L$ の上成分 $u_L$ については $T_3=+1/2$, $Y=1/3$ なので
\begin{equation}
Q=\frac12+\frac16=\frac23
\end{equation}
となり、$u$ クォークの電荷を正しく与える。
右巻きの $u_R$ は $T_3=0$ なので $Q=Y/2=2/3$ であり、
左巻き成分と右巻き成分の電荷は一致する。
これは電磁相互作用がパリティを保存することと整合している。
\end{myex}
\subsection{表から読み取れること}
\begin{itemize}
    \item 同じ粒子の左巻き成分と右巻き成分は、
    異なる $\SU(2)_L$ 表現と異なる $Y$ を持つが、
    同じ電荷 $Q$ を持つ。
    \item 二重項の上成分と下成分の電荷の差は常に $1$ であり、
    $W^\pm$ ボソンの電荷 $\pm1$ に対応する。
    \item 二重項の超電荷は、その平均電荷の2倍に等しい：
    $Y(Q_L)=2\times\frac12(\frac23-\frac13)=\frac13$,
    $Y(L_L)=2\times\frac12(0-1)=-1$。
    一重項の超電荷は電荷の2倍である。
    \item 最小標準模型には右巻きニュートリノ $\nu_R$ は含まれない。
    仮に導入すると $(\mathbf{1},\mathbf{1})_0$ であり、
    いかなるゲージ相互作用も受けない（\term{ステライル}{sterile}）。
    \item ヒッグス二重項の下成分 $H^0$ は電気的に中性であり、
    これが真空期待値を持っても $\U(1)_{\mathrm{EM}}$ は破れない。
\end{itemize}
\subsection{ヒッグス場の共役（conjugate Higgs field）}
$H\sim(\mathbf{1},\mathbf{2})_{+1}$ に対し、
複素共役 $H^*$ は $(\mathbf{1},\bar{\mathbf{2}})_{-1}$ に属する。
\eqref{eq:pseudoreal} を用いると
\begin{equation}
\tilde H\equiv i\sigma^2H^*
=\begin{pmatrix}0&1\\-1&0\end{pmatrix}
\begin{pmatrix}H^{+*}\\ H^{0*}\end{pmatrix}
=\begin{pmatrix}H^{0*}\\ -H^-\end{pmatrix}
\end{equation}
は $(\mathbf{1},\mathbf{2})_{-1}$ の二重項として変換する（$H^-\equiv H^{+*}$）。
$\tilde H$ の上成分は中性、下成分は電荷 $-1$ である。
$H$ と $\tilde H$ の両方を用いることで、
上型クォークと下型クォークの両方に質量を与えることができる。
\subsection{すべて左巻きで書く表示（all left-handed notation）}
右巻きの Weyl 場の荷電共役\en{charge conjugate}は左巻きの Weyl 場になる。
$(\psi_R)^c$ は $\psi_R$ の複素共役表現に属するので、
すべてのフェルミオンを左巻き場で表すと次のようになる。
\begin{center}
\renewcommand{\arraystretch}{1.3}
\begin{tabular}{ccccc}
\toprule
場 & $\SU(3)_C$ & $\SU(2)_L$ & $Y$ & Weyl 成分数\\
\midrule
$Q$ & $\mathbf{3}$ & $\mathbf{2}$ & $+1/3$ & 6\\
$u^c$ & $\bar{\mathbf{3}}$ & $\mathbf{1}$ & $-4/3$ & 3\\
$d^c$ & $\bar{\mathbf{3}}$ & $\mathbf{1}$ & $+2/3$ & 3\\
$L$ & $\mathbf{1}$ & $\mathbf{2}$ & $-1$ & 2\\
$e^c$ & $\mathbf{1}$ & $\mathbf{1}$ & $+2$ & 1\\
\bottomrule
\end{tabular}
\end{center}
この表示はゲージ異常の計算や大統一理論（$\SU(5)$ では
$\bar{\mathbf{5}}=d^c\oplus L$, $\mathbf{10}=Q\oplus u^c\oplus e^c$ にまとまる）で便利である。
\subsection{異常相殺の予告（preview of anomaly cancellation）}
カイラルなゲージ理論では、量子効果によってゲージ対称性が破れる
\term{ゲージ異常}{gauge anomaly}が生じ得る。異常が相殺するための条件の一部は、
左巻き表示の Weyl フェルミオンについての和
\begin{equation}
\sum Y=0,\qquad\sum Y^3=0
\end{equation}
である（なぜこの条件が必要かは第3章で説明する。ここでは数値だけ確かめる）。1世代について確かめると
\begin{align}
\sum Y&=
6\cdot\frac13+3\cdot\left(-\frac43\right)+3\cdot\frac23+2\cdot(-1)+1\cdot2
=2-4+2-2+2=0,\\
\sum Y^3&=
6\cdot\frac{1}{27}+3\cdot\left(-\frac{64}{27}\right)+3\cdot\frac{8}{27}+2\cdot(-1)+1\cdot8
=-6-2+8=0
\end{align}
となる。
\begin{note}
$\sum Y^3$ の各項：
$Q$：$6\times\frac1{27}=\frac29$、
$u^c$：$3\times(-\frac{64}{27})=-\frac{64}{9}$、
$d^c$：$3\times\frac8{27}=\frac89$。
クォークの和は $\frac{2-64+8}{9}=-6$。
レプトンは $L$：$2\times(-1)=-2$、$e^c$：$1\times8=8$ で和は $+6$。
クォーク部分とレプトン部分がそれぞれ $\mp6$ で打ち消している。
\end{note}
クォークとレプトンの寄与が絶妙に打ち消し合っており、
\textbf{1世代はクォークとレプトンが揃って初めて矛盾のない理論になる}。
これは「クォークとレプトンの間に何らかの関係がある」ことを強く示唆している。
\subsection{大域的対称性：バリオン数とレプトン数（global symmetries: $B$ and $L$）}
標準模型のラグランジアンは、ゲージ対称性の他に、
課さなくても自動的に現れる\term{偶然的対称性}{accidental symmetry}を持つ。
\begin{itemize}
    \item \term{バリオン数}{baryon number} $B$：すべてのクォークに $B=1/3$、レプトンに $B=0$。
    \item \term{レプトン数}{lepton number} $L$：すべてのレプトンに $L=1$、クォークに $L=0$。
    ニュートリノが massless なら、各世代のレプトン数 $L_e$, $L_\mu$, $L_\tau$ が別々に保存する。
\end{itemize}
これらは繰り込み可能な項だけを許したことの帰結であり、
陽子が（摂動論の範囲で）安定である\en{proton stability}ことを説明する。
ただし $B$ と $L$ は量子異常により非摂動的に\en{non-perturbatively}破れ、$B-L$ のみが厳密に保存する。
%------------------------------------------------
\section{左巻き・右巻きフェルミオン}
\label{sec:chirality}
%------------------------------------------------
\subsection{カイラル射影（chiral projection）}
\begin{mydef}{カイラル射影とカイラリティ}{chirality}
Dirac スピノル $\psi$ に対して、カイラル射影演算子を
\begin{equation}
\PL=\frac{1-\gamma^5}{2},
\qquad
\PR=\frac{1+\gamma^5}{2}
\end{equation}
で定義し、左巻き\en{left-handed}・右巻き\en{right-handed}成分を
\begin{equation}
\psi_L=\PL\psi,\qquad \psi_R=\PR\psi,
\qquad
\psi=\psi_L+\psi_R
\end{equation}
と定義する。$\gamma^5\psi_L=-\psi_L$, $\gamma^5\psi_R=+\psi_R$ であり、
$\gamma^5$ の固有値 $\mp1$ を\term{カイラリティ}{chirality}と呼ぶ。
\end{mydef}
$(\gamma^5)^2=1$ から
\begin{equation}
\PL^2=\PL,\qquad \PR^2=\PR,\qquad
\PL\PR=\PR\PL=0,\qquad
\PL+\PR=1
\end{equation}
が成り立ち、これらは射影演算子\en{projection operator}である。
\begin{note}
$\PL^2=\frac14(1-2\gamma^5+(\gamma^5)^2)=\frac14(2-2\gamma^5)=\PL$。
$\PL\PR=\frac14(1-(\gamma^5)^2)=0$。
$\gamma^5\PL=\frac12(\gamma^5-1)=-\PL$ より $\gamma^5\psi_L=-\psi_L$。
\end{note}
Dirac 共役 $\bar\psi=\psi^\dagger\gamma^0$ については、
$\gamma^5$ がエルミートで $\gamma^0$ と反可換であることから
\begin{equation}
\overline{\psi_L}
=(\PL\psi)^\dagger\gamma^0
=\psi^\dagger\PL\gamma^0
=\psi^\dagger\gamma^0\PR
=\bar\psi\PR
\end{equation}
となる。同様に $\overline{\psi_R}=\bar\psi\PL$ である。
また $\gamma^\mu$ と $\gamma^5$ の反可換性から
\begin{equation}
\PR\gamma^\mu=\gamma^\mu\PL,\qquad
\PL\gamma^\mu=\gamma^\mu\PR
\end{equation}
が成り立つ。
\subsection{Weyl 表示（Weyl / chiral representation）}
カイラリティの構造は、$\gamma^5$ を対角にする Weyl（カイラル）表示で最も見やすい：
\begin{equation}
\gamma^\mu=\begin{pmatrix}0&\sigma^\mu\\ \bar\sigma^\mu&0\end{pmatrix},
\qquad
\gamma^5=\begin{pmatrix}-1&0\\0&1\end{pmatrix},
\qquad
\sigma^\mu=(1,\bm\sigma),\quad
\bar\sigma^\mu=(1,-\bm\sigma).
\end{equation}
このとき Dirac スピノルは2成分スピノルの組
\begin{equation}
\psi=\begin{pmatrix}\psi_L\\ \psi_R\end{pmatrix}
\end{equation}
に分かれ、$\psi_L$, $\psi_R$ はそれぞれ Lorentz 群の
$(\frac12,0)$, $(0,\frac12)$ 表現に属する。
これらは Lorentz 変換で混ざらないので、
\textbf{別々の場として独立に扱うことができる}。

Dirac ラグランジアンは
\begin{equation}
\Lag
=\bar\psi(i\gamma^\mu\partial_\mu-m)\psi
=
\psi_L^\dagger i\bar\sigma^\mu\partial_\mu\psi_L
+\psi_R^\dagger i\sigma^\mu\partial_\mu\psi_R
-m\left(\psi_L^\dagger\psi_R+\psi_R^\dagger\psi_L\right)
\end{equation}
となる。質量項だけが $\psi_L$ と $\psi_R$ を結びつけていることが明らかである。
\begin{note}
$\gamma^0=\begin{pmatrix}0&1\\1&0\end{pmatrix}$ なので
$\bar\psi=\psi^\dagger\gamma^0=(\psi_R^\dagger,\ \psi_L^\dagger)$。また
$\gamma^\mu\partial_\mu\psi=\begin{pmatrix}\sigma^\mu\partial_\mu\psi_R\\ \bar\sigma^\mu\partial_\mu\psi_L\end{pmatrix}$。
よって
\[
\bar\psi\gamma^\mu\partial_\mu\psi
=\psi_R^\dagger\sigma^\mu\partial_\mu\psi_R+\psi_L^\dagger\bar\sigma^\mu\partial_\mu\psi_L,
\qquad
\bar\psi\psi=\psi_R^\dagger\psi_L+\psi_L^\dagger\psi_R .
\]
$\gamma^5=\diag(-1,1)$ が $\gamma^5=i\gamma^0\gamma^1\gamma^2\gamma^3$ と一致することは、
$\gamma^0\gamma^i=\diag(-\sigma^i,\sigma^i)$ を用いて直接確かめられる。
\end{note}
\subsection{運動項と質量項（kinetic and mass terms）}
一般の表示で同じことを確かめよう。運動項は
\begin{equation}
\bar\psi i\gamma^\mu\partial_\mu\psi
=
\overline{\psi_L}i\gamma^\mu\partial_\mu\psi_L
+
\overline{\psi_R}i\gamma^\mu\partial_\mu\psi_R
\end{equation}
と左右に\textbf{分離}\en{decouple}する。
実際、交差項は
$\overline{\psi_L}\gamma^\mu\psi_R=\bar\psi\PR\gamma^\mu\PR\psi=\bar\psi\gamma^\mu\PL\PR\psi=0$
となる。
同様にゲージ場との結合 $\bar\psi\gamma^\mu A_\mu\psi$ も左右に分離するので、
$\psi_L$ と $\psi_R$ に異なるゲージ変換性を与えても
運動項のゲージ不変性は損なわれない。

一方、質量項は
\begin{equation}
m\bar\psi\psi
=
m\left(\overline{\psi_L}\psi_R+\overline{\psi_R}\psi_L\right)
\end{equation}
となり、左右を\textbf{混合}\en{mix}する
（$\overline{\psi_L}\psi_L=\bar\psi\PR\PL\psi=0$ に注意）。
標準模型では $\psi_L$ と $\psi_R$ が異なるゲージ表現に属するため、
この項はゲージ不変ではない。例えば電子の場合
\begin{equation}
\overline{e_L}e_R:\quad
e_L\ \text{は}\ \SU(2)_L\ \text{二重項の成分},\quad
e_R\ \text{は一重項},\quad
Y:\ +1-2=-1\neq0
\end{equation}
であり、$\SU(2)_L$ と $\U(1)_Y$ の両方で不変ではない。
したがって、この項をラグランジアンに直接書き込むことはできない。
これがフェルミオンの質量をヒッグス機構によって生成する必要がある理由である（第3章、第7章）。
\subsection{双線形形式と V$-$A 構造（bilinears and V$-$A structure）}
Dirac 双線形形式\en{bilinear}を左右の成分で書くと
\begin{align}
\bar\psi\gamma^\mu\psi&=\overline{\psi_L}\gamma^\mu\psi_L+\overline{\psi_R}\gamma^\mu\psi_R,\\
\bar\psi\gamma^\mu\gamma^5\psi&=-\overline{\psi_L}\gamma^\mu\psi_L+\overline{\psi_R}\gamma^\mu\psi_R
\end{align}
となる。
\begin{note}
$\gamma^5\psi=-\psi_L+\psi_R$ なので
$\bar\psi\gamma^\mu\gamma^5\psi=\bar\psi\gamma^\mu(-\psi_L+\psi_R)$。
$\bar\psi\gamma^\mu\psi_L=\overline{\psi_L}\gamma^\mu\psi_L$
（$\overline{\psi_R}\gamma^\mu\psi_L=\bar\psi\PL\gamma^\mu\PL\psi=\bar\psi\gamma^\mu\PR\PL\psi=0$ による）を用いればよい。
また $1-\gamma^5=2\PL$ から直接
$\bar\psi\gamma^\mu(1-\gamma^5)\psi=2\bar\psi\gamma^\mu\PL\psi=2\bar\psi\PR\gamma^\mu\PL\psi=2\overline{\psi_L}\gamma^\mu\psi_L$
（$\PL^2=\PL$ を用いて $\PL$ を2つに分けた）。
\end{note}
したがって
\begin{equation}
\bar\psi\gamma^\mu(1-\gamma^5)\psi=2\,\overline{\psi_L}\gamma^\mu\psi_L
\end{equation}
であり、左巻き成分のみが現れる。
ベクトル（vector, V）カレントから軸性ベクトル（axial vector, A）カレントを引いた
この形を \term{V$-$A 構造}{V$-$A structure} と呼ぶ。
$W$ ボソンとの結合はすべてこの形をとる（導出は第8章。ここでは結果のみ示す）：
\begin{equation}
\Lag_{\mathrm{CC}}=\frac{g}{\sqrt2}\,
\overline{\nu_L}\gamma^\mu e_L\,W_\mu^++\hc
=\frac{g}{2\sqrt2}\,
\bar\nu\gamma^\mu(1-\gamma^5)e\,W_\mu^++\hc
\end{equation}
\subsection{カイラリティとヘリシティ（chirality vs.\ helicity）}
カイラリティ（$\gamma^5$ の固有値）は Lorentz 不変な概念であるが、
質量があると運動方程式で保存しない。
これに対しヘリシティ\en{helicity}
\begin{equation}
h=\frac{\bm{S}\cdot\bm{p}}{|\bm{p}|}
\end{equation}
は自由粒子で保存するが、massive な粒子では
粒子を追い越す Lorentz 系で符号が反転するため、Lorentz 不変ではない。

massless の場合、Weyl 表示で $\psi_L\propto u\,e^{-ip\cdot x}$ とおくと
運動方程式 $i\bar\sigma^\mu\partial_\mu\psi_L=0$ は
\begin{equation}
(E+\bm\sigma\cdot\bm p)u=0
\quad\Longrightarrow\quad
\frac{\bm\sigma\cdot\bm p}{|\bm p|}u=-u
\end{equation}
となり（$E=|\bm p|$）、左巻きの粒子はヘリシティ $-1/2$ を持つ。
\begin{note}
$i\bar\sigma^\mu\partial_\mu(u\,e^{-ip\cdot x})=\bar\sigma^\mu p_\mu u\,e^{-ip\cdot x}$。
$\bar\sigma^\mu=(1,-\sigma^i)$、共変成分 $p_\mu=(E,-p^i)$ なので
$\bar\sigma^\mu p_\mu=E+(-\sigma^i)(-p^i)=E+\bm\sigma\cdot\bm p$ となる。
スピン演算子は $\bm S=\bm\sigma/2$ なので $h=-1/2$。
同様に $\psi_R$ は $(E-\bm\sigma\cdot\bm p)u=0$ からヘリシティ $+1/2$ を持つ。
\end{note}
同様に、左巻きの場が生成する反粒子はヘリシティ $+1/2$ を持つ。
したがってニュートリノは常にヘリシティ $-1/2$、
反ニュートリノは常にヘリシティ $+1/2$ で観測される（Goldhaber らの実験, 1958）。

massive な粒子（質量 $m$、エネルギー $E$）では、
左巻きカイラリティの状態がヘリシティ $+1/2$ の成分を含む振幅は
$O(m/E)$ である。
\begin{myex}{荷電パイ中間子の崩壊とヘリシティ抑制（helicity suppression）}{pion}
$\pi^-\to\ell^-\bar\nu_\ell$（$\ell=e,\mu$）を考える。
$\pi^-$ はスピン0なので、静止系で背中合わせに出る
$\ell^-$ と $\bar\nu_\ell$ のスピンは逆向き、
すなわち両者のヘリシティは等しくなければならない。
$\bar\nu_\ell$ はヘリシティ $+1/2$ なので $\ell^-$ もヘリシティ $+1/2$ でなければならないが、
V$-$A 結合は左巻きの $\ell^-$ しか生成しないので、
この振幅は $m_\ell/E$ に比例して suppress される。
その結果、崩壊幅の比は
\begin{equation}
R_\pi=\frac{\Gamma(\pi\to e\bar\nu)}{\Gamma(\pi\to\mu\bar\nu)}
=\frac{m_e^2}{m_\mu^2}
\left(\frac{m_\pi^2-m_e^2}{m_\pi^2-m_\mu^2}\right)^2
\simeq1.28\times10^{-4}
\end{equation}
となる。位相空間\en{phase space}だけを考えると軽い電子への崩壊の方が有利なのにもかかわらず、
実際には電子への崩壊は1万分の1程度しか起こらない。
放射補正\en{radiative correction}を含めた理論値 $1.235\times10^{-4}$ は実験値とよく一致し、
V$-$A 構造の強力な証拠となっている。
（崩壊幅の具体的な計算は第12章で行う。）
\end{myex}
\begin{note}
崩壊幅は $\Gamma(\pi\to\ell\bar\nu)\propto m_\ell^2\,(m_\pi^2-m_\ell^2)^2/m_\pi^3$ の形をとる。
$m_\ell^2$ が helicity suppression、$(m_\pi^2-m_\ell^2)^2$ が位相空間と行列要素の運動学的因子である。
数値は $m_e=0.511$, $m_\mu=105.66$, $m_\pi=139.57$（MeV）を用いて
\[
\frac{m_e^2}{m_\mu^2}=2.34\times10^{-5},
\qquad
\frac{m_\pi^2-m_e^2}{m_\pi^2-m_\mu^2}=\frac{19480}{8316}=2.343,
\qquad
R_\pi=2.34\times10^{-5}\times(2.343)^2\simeq1.28\times10^{-4}.
\]
\end{note}
\subsection{離散対称性 $C$, $P$, $CP$（discrete symmetries）}
パリティ $P$ は空間反転\en{spatial inversion} $\bm x\to-\bm x$ であり、運動量を反転しスピンを変えないので
ヘリシティを反転する。Dirac 場に対しては $\psi\to\gamma^0\psi$ と作用し、
$\psi_L\leftrightarrow\psi_R$ を入れ替える。
荷電共役 $C$ は粒子と反粒子を入れ替え、ヘリシティは変えない。

左巻きのニュートリノ $\nu_L$ に対して
\begin{equation}
\nu_L\ (h=-\tfrac12)
\xrightarrow{\ P\ }
\nu\ (h=+\tfrac12)\ \text{（存在しない）},
\qquad
\nu_L\ (h=-\tfrac12)
\xrightarrow{\ C\ }
\bar\nu\ (h=-\tfrac12)\ \text{（存在しない）}
\end{equation}
となり、$P$ と $C$ はそれぞれ最大限に破れている。
一方、$CP$ は $\nu\ (h=-\frac12)$ を存在する状態 $\bar\nu\ (h=+\frac12)$ に移すので、
ゲージ相互作用は $CP$ を保存する。
標準模型で $CP$ を破るのは Yukawa 結合に含まれる複素位相
（CKM 位相、第9章）と QCD の $\theta$ 項のみである。
\subsection{1世代の自由度（degrees of freedom per generation）}
各 Dirac 場 $\psi$ は2つの Weyl スピノル $\psi_L$, $\psi_R$ からなる。
標準模型の1世代に含まれる Weyl スピノルの数は
\begin{equation}
\underbrace{3\times2}_{Q_L}
+\underbrace{3}_{u_R}
+\underbrace{3}_{d_R}
+\underbrace{2}_{L_L}
+\underbrace{1}_{e_R}
=15
\end{equation}
である。3世代で45個となる。
右巻きニュートリノを加えると1世代16個となり、
これは $\mathrm{SO}(10)$ 大統一理論のスピノル表現 $\mathbf{16}$ にちょうど収まる。
%------------------------------------------------
\section{標準模型のラグランジアンの全体構造}
\label{sec:SM-lagrangian}
%------------------------------------------------
標準模型のラグランジアンは
\begin{equation}
\boxed{
\Lag_{\SM}
=
\Lag_{\mathrm{gauge}}
+\Lag_{\mathrm{fermion}}
+\Lag_{\mathrm{Higgs}}
+\Lag_{\mathrm{Yukawa}}
\ (+\Lag_\theta)
}
\end{equation}
と分解できる（量子化の際にはこれにゲージ固定項\en{gauge-fixing term}とゴースト項\en{ghost term}が加わる。付録E）。
以下に各項の具体形を示す。導出と物理的意味は後の章で詳しく扱う。
\subsection{共変微分（covariant derivative）}
\begin{equation}
D_\mu
=
\partial_\mu
-ig_sG_\mu^aT^a
-igW_\mu^i\frac{\sigma^i}{2}
-ig'\frac{Y}{2}B_\mu .
\end{equation}
各項は、その場が非自明な表現に属する場合にのみ作用する。各場に対して具体的に書くと
\begin{align}
D_\mu Q_L&=\left(\partial_\mu-ig_sG_\mu^aT^a-igW_\mu^i\frac{\sigma^i}{2}-\frac{i}{6}g'B_\mu\right)Q_L,\\
D_\mu u_R&=\left(\partial_\mu-ig_sG_\mu^aT^a-\frac{2i}{3}g'B_\mu\right)u_R,\\
D_\mu d_R&=\left(\partial_\mu-ig_sG_\mu^aT^a+\frac{i}{3}g'B_\mu\right)d_R,\\
D_\mu L_L&=\left(\partial_\mu-igW_\mu^i\frac{\sigma^i}{2}+\frac{i}{2}g'B_\mu\right)L_L,\\
D_\mu e_R&=\left(\partial_\mu+ig'B_\mu\right)e_R,\\
D_\mu H&=\left(\partial_\mu-igW_\mu^i\frac{\sigma^i}{2}-\frac{i}{2}g'B_\mu\right)H
\end{align}
である。
\subsection{ゲージ場の運動項（gauge kinetic terms）}
\begin{equation}
\Lag_{\mathrm{gauge}}
=
-\frac14G_{\mu\nu}^aG^{a\mu\nu}
-\frac14W_{\mu\nu}^iW^{i\mu\nu}
-\frac14B_{\mu\nu}B^{\mu\nu},
\end{equation}
\begin{align}
G_{\mu\nu}^a&=\partial_\mu G_\nu^a-\partial_\nu G_\mu^a+g_sf^{abc}G_\mu^bG_\nu^c,\\
W_{\mu\nu}^i&=\partial_\mu W_\nu^i-\partial_\nu W_\mu^i+g\epsilon^{ijk}W_\mu^jW_\nu^k,\\
B_{\mu\nu}&=\partial_\mu B_\nu-\partial_\nu B_\mu .
\end{align}
場の強さ\en{field strength}は $[D_\mu,D_\nu]=-igF_{\mu\nu}^aT^a$ から得られる（詳しくは第2章）。
\begin{note}
$D_\mu=\partial_\mu-igA_\mu$（$A_\mu=A_\mu^aT^a$）とすると、試験関数に作用させて
\[
[D_\mu,D_\nu]
=-ig(\partial_\mu A_\nu-\partial_\nu A_\mu)+(-ig)^2[A_\mu,A_\nu].
\]
$[A_\mu,A_\nu]=A_\mu^bA_\nu^c[T^b,T^c]=if^{abc}A_\mu^bA_\nu^cT^a$ を代入すると
\[
[D_\mu,D_\nu]=-ig\left(\partial_\mu A_\nu^a-\partial_\nu A_\mu^a+gf^{abc}A_\mu^bA_\nu^c\right)T^a
=-igF^a_{\mu\nu}T^a .
\]
\end{note}
非可換ゲージ場の場の強さに含まれる2次の項により、
$-\frac14F^2$ を展開すると
\begin{equation}
-\frac14F^a_{\mu\nu}F^{a\mu\nu}
\supset
-gf^{abc}(\partial_\mu A_\nu^a)A^{b\mu}A^{c\nu}
-\frac{g^2}{4}f^{abe}f^{cde}A_\mu^aA_\nu^bA^{c\mu}A^{d\nu}
\end{equation}
が現れ、グルーオンや $W$ ボソンは3点・4点の自己相互作用を持つ。
\begin{note}
$F^a_{\mu\nu}=f^a_{\mu\nu}+gf^{abc}A^b_\mu A^c_\nu$（$f^a_{\mu\nu}=\partial_\mu A_\nu^a-\partial_\nu A_\mu^a$）とおくと
\[
-\tfrac14F^aF^a
=-\tfrac14f^af^a
-\tfrac12gf^{abc}f^a_{\mu\nu}A^{b\mu}A^{c\nu}
-\tfrac14g^2f^{abe}f^{cde}A^a_\mu A^b_\nu A^{c\mu}A^{d\nu}.
\]
2項目で $f^a_{\mu\nu}$ の2つの項は、$f^{abc}$ の $b\leftrightarrow c$ 反対称性と $\mu\leftrightarrow\nu$ の入れ替えにより等しい寄与を与えるので、
$-gf^{abc}(\partial_\mu A^a_\nu)A^{b\mu}A^{c\nu}$ となる。
\end{note}
ゲージ場の質量項 $\frac12M^2A_\mu A^\mu$ はゲージ不変\en{gauge invariant}でないため含まれない。
\subsection{フェルミオンの運動項（fermion kinetic terms）}
\begin{equation}
\Lag_{\mathrm{fermion}}
=
\sum_{\text{世代}}
\left(
\overline{Q_L}i\gamma^\mu D_\mu Q_L
+\overline{u_R}i\gamma^\mu D_\mu u_R
+\overline{d_R}i\gamma^\mu D_\mu d_R
+\overline{L_L}i\gamma^\mu D_\mu L_L
+\overline{e_R}i\gamma^\mu D_\mu e_R
\right).
\end{equation}
この項がフェルミオンとゲージボソンの相互作用をすべて含む。
共変微分を展開すると、相互作用は
\begin{equation}
\Lag_{\mathrm{fermion}}\supset
g_s\,\bar q\gamma^\mu T^aq\,G_\mu^a
+g\,\overline{\Psi_L}\gamma^\mu\frac{\sigma^i}{2}\Psi_L\,W_\mu^i
+g'\,\bar\psi\gamma^\mu\frac{Y}{2}\psi\,B_\mu
\end{equation}
の形をとる（$\Psi_L$ は左巻き二重項）。
すなわち各ゲージ場は、対応する対称性の\term{保存カレント}{conserved current}と結合する。
ゲージ対称性のため、1つのゲージ群に属するすべての場は
\textbf{同じ}結合定数で結合する（ゲージ結合の普遍性\en{universality}）。
例えば $W$ ボソンとレプトンの結合と、$W$ ボソンとクォークの結合は同じ $g$ で与えられる。
\subsection{ヒッグス部分（Higgs sector）}
\begin{equation}
\Lag_{\mathrm{Higgs}}
=
(D_\mu H)^\dagger(D^\mu H)-V(H),
\qquad
V(H)=-\mu^2H^\dagger H+\lambda(H^\dagger H)^2 .
\end{equation}
$\SU(2)_L\times\U(1)_Y$ 不変かつ次元4以下のスカラーポテンシャルは $H^\dagger H$ の2次式に限られる
（$H^3$ のような奇数次の項は $Y$ の保存則から禁止される）。
真空の安定性\en{vacuum stability}から $\lambda>0$ が必要である。

$\mu^2>0$ のとき $V$ の最小値は
\begin{equation}
H^\dagger H=\frac{v^2}{2},\qquad v=\frac{\mu}{\sqrt\lambda}
\end{equation}
にあり、電弱対称性が自発的に破れる。
ユニタリーゲージ\en{unitary gauge}で $H=\frac{1}{\sqrt2}(0,\ v+h)^T$ とおくと、
\begin{itemize}
    \item 運動項 $(D_\mu H)^\dagger(D^\mu H)$ から $W^\pm$, $Z^0$ の質量項
    \begin{equation}
    M_W=\frac{gv}{2},\qquad
    M_Z=\frac{\sqrt{g^2+g'^2}\,v}{2}=\frac{M_W}{\cos\theta_W}
    \end{equation}
    および $hWW$, $hZZ$ 結合が現れる。
    \item ポテンシャルから、ヒッグス粒子の質量と自己結合\en{self-coupling}
    \begin{equation}
    m_h=\sqrt{2\lambda}\,v=\sqrt2\mu,\qquad
    V\supset\lambda vh^3+\frac\lambda4h^4
    \end{equation}
    が得られる。
\end{itemize}
（ゲージボソンの質量の導出は第6章で行う。）
\begin{note}
ポテンシャル部分だけ確かめておく。$H^\dagger H=\frac12(v+h)^2$ を代入し $\mu^2=\lambda v^2$ を用いると
\[
V=-\frac{\lambda v^2}{2}(v+h)^2+\frac\lambda4(v+h)^4 .
\]
$h^2$ の係数：$-\frac{\lambda v^2}{2}+\frac\lambda4\cdot6v^2=\lambda v^2$。
これを $\frac12m_h^2h^2$ と比べて $m_h^2=2\lambda v^2$。
$h^3$ の係数：$\frac\lambda4\cdot4v=\lambda v$。$h^4$ の係数：$\frac\lambda4$。
$h$ の1次の係数は $-\lambda v^3+\lambda v^3=0$ となり、$h=0$ が極値であることと整合する。
\end{note}
観測値 $m_h\simeq125\ \mathrm{GeV}$ から
$\lambda=m_h^2/(2v^2)=125^2/(2\times246^2)\simeq0.13$,
$\mu=m_h/\sqrt2\simeq88\ \mathrm{GeV}$ と決まる。
$M_W/M_Z=\cos\theta_W$ の関係は $\rho$ パラメータ\en{$\rho$ parameter}
\begin{equation}
\rho\equiv\frac{M_W^2}{M_Z^2\cos^2\theta_W}=1
\end{equation}
として表され、ヒッグス場が二重項であることの帰結である（第6章および第12章の電弱精密測定）。
\subsection{Yukawa 相互作用（Yukawa interactions）}
\begin{equation}
\Lag_{\mathrm{Yukawa}}
=
-y^d_{ij}\,\overline{Q_{Li}}\,H\,d_{Rj}
-y^u_{ij}\,\overline{Q_{Li}}\,\tilde H\,u_{Rj}
-y^e_{ij}\,\overline{L_{Li}}\,H\,e_{Rj}
+\hc
\end{equation}
ここで $i,j=1,2,3$ は世代の添字、$y^{u,d,e}$ は一般の $3\times3$ 複素行列である。

各項のゲージ不変性を確かめる。$\SU(3)_C$ については $\bar q\,q$ の形で縮約され、
$\SU(2)_L$ については $\overline{\Psi_L}\,\Phi$ が二重項同士の内積をなす。
超電荷の和は（$\bar\psi$ は $-Y$ を持つことに注意して）
\begin{align}
\overline{Q_L}Hd_R&:\quad -\tfrac13+1-\tfrac23=0,\\
\overline{Q_L}\tilde Hu_R&:\quad -\tfrac13-1+\tfrac43=0,\\
\overline{L_L}He_R&:\quad +1+1-2=0
\end{align}
となり、すべてゲージ不変である。

ヒッグス場が真空期待値を持つと、例えば
\begin{equation}
-y_e\overline{L_L}He_R+\hc
\ \longrightarrow\
-\frac{y_ev}{\sqrt2}\,\overline{e_L}e_R-\frac{y_e}{\sqrt2}\,h\,\overline{e_L}e_R+\hc
=-m_e\bar ee\left(1+\frac hv\right)
\end{equation}
となり、フェルミオンの質量 $m_e=y_ev/\sqrt2$ と、
質量に比例するヒッグス結合 $m_e/v$ が同時に得られる。
世代間の混合を含む一般の場合、
Yukawa 行列の対角化から CKM 行列が現れる（第9章）。
\subsection{$\theta$ 項（$\theta$ term）}
上記の他に、ゲージ不変かつ繰り込み可能な項として
QCD の $\theta$ 項
\begin{equation}
\Lag_\theta=\theta\frac{g_s^2}{32\pi^2}G_{\mu\nu}^a\tilde G^{a\mu\nu},
\qquad
\tilde G^{a\mu\nu}=\frac12\epsilon^{\mu\nu\rho\sigma}G^a_{\rho\sigma}
\end{equation}
が存在する。これは全微分であり摂動論には寄与しないが、
インスタントン\en{instanton}などの非摂動的効果を通じて $P$ と $CP$ を破る。
中性子の電気双極子モーメント\en{neutron electric dipole moment, nEDM}の上限から $|\bar\theta|\lesssim10^{-10}$ であり、
なぜこれほど小さいかは\term{強い CP 問題}{strong CP problem}と呼ばれる（第14章）。
$\SU(2)_L$ の $\theta$ 項は、$B+L$ の異常を用いた場の再定義で消去できるため物理的でない。
\subsection{繰り込み可能性による制限（constraints from renormalizability）}
4次元時空で作用 $S=\int d^4x\,\Lag$ が無次元であるためには、
$\Lag$ の質量次元は4でなければならない。
運動項から場の質量次元は
\begin{equation}
[\phi]=[A_\mu]=1,\qquad [\psi]=\frac32
\end{equation}
と決まる。各項の演算子の次元は次のとおりである。
\begin{center}
\renewcommand{\arraystretch}{1.3}
\begin{tabular}{lcc}
\toprule
項 & 演算子の次元 & 係数の次元\\
\midrule
$H^\dagger H$ & 2 & $[\mu^2]=2$\\
$(H^\dagger H)^2$, $|D H|^2$ & 4 & $[\lambda]=0$\\
$\bar\psi\gamma^\mu D_\mu\psi$ & 4 & 0\\
$F_{\mu\nu}F^{\mu\nu}$ & 4 & 0\\
$\bar\psi H\psi$ & 4 & $[y]=0$\\
$(\overline{L}\tilde H)(\tilde H^TL^c)$ & 5 & $-1$\\
$\overline{Q}\,\overline{Q}\,Q\,L$ & 6 & $-2$\\
\bottomrule
\end{tabular}
\end{center}
結合定数の質量次元が負にならない（演算子の次元が4以下である）ことが
繰り込み可能性の条件となる。
$\Lag_{\SM}$ は、与えられた場の内容とゲージ対称性の下で、
この条件を満たす\textbf{すべての}項を含んでいる。
有効場の理論の観点から次元 $d>4$ の演算子を加えると
\begin{equation}
\Lag_{\mathrm{eff}}=\Lag_{\SM}
+\frac{c_5}{\Lambda}\mathcal{O}_5
+\sum\frac{c_6}{\Lambda^2}\mathcal{O}_6+\cdots
\end{equation}
となる。唯一の次元5演算子は Weinberg 演算子
$(\overline{L_L}\tilde H)(\tilde H^T L_L^c)$ であり、
対称性の破れの後ニュートリノの Majorana 質量 $m_\nu\sim v^2/\Lambda$ を与える（第11章）。
$m_\nu\sim0.1\ \mathrm{eV}$ とすると $\Lambda\sim10^{14}$--$10^{15}\ \mathrm{GeV}$ となる。
次元6には陽子崩壊\en{proton decay}を引き起こす4フェルミオン演算子\en{four-fermion operator}が含まれる。
\subsection{相互作用頂点の分類（classification of vertices）}
対称性の破れの後、ラグランジアンを展開すると次のような頂点が現れる。
\begin{center}
\renewcommand{\arraystretch}{1.3}
\begin{tabular}{lll}
\toprule
起源 & 頂点 & 結合の強さ\\
\midrule
$\Lag_{\mathrm{fermion}}$ & $\bar qqg$ & $g_s$\\
 & $\bar ff\gamma$ & $eQ_f$\\
 & $\bar f f'W$ & $g/\sqrt2$\\
 & $\bar ffZ$ & $\frac{g}{\cos\theta_W}(T_3-Q\sin^2\theta_W)$\\
$\Lag_{\mathrm{gauge}}$ & $ggg$,\ $gggg$ & $g_s$,\ $g_s^2$\\
 & $WW\gamma$,\ $WWZ$ & $e$,\ $g\cos\theta_W$\\
 & $WWWW$,\ $WW\gamma\gamma$ など & $g^2$ など\\
$\Lag_{\mathrm{Higgs}}$ & $hWW$,\ $hZZ$ & $gM_W$,\ $gM_Z/\cos\theta_W$\\
 & $hhh$,\ $hhhh$ & $3m_h^2/v$,\ $3m_h^2/v^2$\\
$\Lag_{\mathrm{Yukawa}}$ & $h\bar ff$ & $m_f/v$\\
\bottomrule
\end{tabular}
\end{center}
（各頂点は第6--8章で導出する。数値係数や $i$ の因子を含む正確な Feynman 則は付録Fにまとめる。）
%------------------------------------------------
\section{標準模型の自由パラメータ}
\label{sec:SM-parameters}
%------------------------------------------------
\subsection{19個のパラメータ（the 19 parameters）}
ニュートリノを massless とする最小標準模型は、
以下の19個の独立なパラメータを持つ。
\begin{center}
\renewcommand{\arraystretch}{1.3}
\begin{tabular}{lcc}
\toprule
sector & パラメータ & 個数\\
\midrule
gauge & $g_s,\ g,\ g'$ & 3\\
Higgs potential & $\mu^2,\ \lambda$ \quad（または $v,\ m_h$） & 2\\
quark Yukawa & $m_u,m_d,m_s,m_c,m_b,m_t$ & 6\\
 & $V_{\mathrm{CKM}}$：$\theta_{12},\theta_{23},\theta_{13},\delta$ & 4\\
lepton Yukawa & $m_e,m_\mu,m_\tau$ & 3\\
QCD vacuum angle & $\theta$ & 1\\
\midrule
合計 & & 19\\
\bottomrule
\end{tabular}
\end{center}
このうち13個（質量と混合）が Yukawa 結合に由来し、
ゲージ部分とヒッグス部分は合わせて5個にすぎない。
標準模型のパラメータの大部分はフレーバーの構造に関するものである。
\subsection{Yukawa 行列のパラメータの数え上げ（parameter counting）}
Yukawa 行列の独立なパラメータが質量と CKM 行列の要素に
まとめられることは自明ではない。
複素 $3\times3$ 行列 $y^u$, $y^d$, $y^e$ は合計 $3\times18=54$ 個の実パラメータを持つが、
その大部分はフェルミオン場の再定義\en{field redefinition}
\begin{equation}
Q_L\to U_QQ_L,\quad
u_R\to U_uu_R,\quad
d_R\to U_dd_R,\quad
L_L\to U_LL_L,\quad
e_R\to U_ee_R
\end{equation}
（$U\in\U(3)$）によって消去できる。
このとき運動項は不変で、Yukawa 行列は
\begin{equation}
y^u\to U_Q^\dagger y^uU_u,\qquad
y^d\to U_Q^\dagger y^dU_d,\qquad
y^e\to U_L^\dagger y^eU_e
\end{equation}
と変わる。
5つの $\U(3)$ は合計 $5\times9=45$ 個のパラメータ（実角度15個、位相30個）を持つ。
ただし、そのうち Yukawa 項を不変に保つ変換
（バリオン数 $\U(1)_B$ と、ニュートリノ質量がゼロの場合に各世代ごとに保存するレプトン数
$\U(1)_{L_e}\times\U(1)_{L_\mu}\times\U(1)_{L_\tau}$ の計4個）は
パラメータを消去しないので、実際に消去できるのは $45-4=41$ 個である。したがって
\begin{equation}
54-41=13=6\ (\text{クォーク質量})+4\ (\text{CKM})+3\ (\text{レプトン質量})
\end{equation}
となる。

クォーク部分だけを取り出すと、
\begin{equation}
\underbrace{2\times18}_{y^u,y^d}
-\left(\underbrace{3\times9}_{U_Q,U_u,U_d}-\underbrace{1}_{\U(1)_B}\right)
=36-26=10=6+4
\end{equation}
であり、CKM 行列には3つの混合角と1つの位相が残る。これは $3\times3$ のユニタリ行列の自由度の数え上げと一致する。
$N$ 世代の場合には混合角 $N(N-1)/2$ 個、位相 $(N-1)(N-2)/2$ 個が残り、
$N\ge3$ で初めて CP を破る位相が現れる。
これが小林・益川による第3世代の予言の根拠である（第9章、付録C）。
\subsection{パラメータの数値（numerical values）}
19個のパラメータを直接 $g_s,g,g',\mu^2,\lambda,\ldots$ として測定するのではなく、
精密に測定できる量を入力\en{input parameters}として用いるのが普通である。
電弱部分の3つのパラメータ $(g,g',v)$ に対しては、例えば
\begin{equation}
\alpha^{-1}=137.036,\qquad
G_F=1.1664\times10^{-5}\ \mathrm{GeV}^{-2},\qquad
M_Z=91.188\ \mathrm{GeV}
\end{equation}
を入力とし、他の量（$M_W$, $\sin^2\theta_W$ など）は予言\en{prediction}とする。
ツリーレベルで換算すると、おおよそ
\begin{equation}
g\simeq0.65,\qquad g'\simeq0.35,\qquad g_s\simeq1.2\ (\mu=M_Z),\qquad
v\simeq246\ \mathrm{GeV},\qquad \lambda\simeq0.13
\end{equation}
である。CKM 行列の要素の大きさはおおよそ
\begin{equation}
|V_{\mathrm{CKM}}|\simeq
\begin{pmatrix}
0.974&0.225&0.004\\
0.225&0.973&0.041\\
0.009&0.040&0.999
\end{pmatrix}
\end{equation}
であり、対角成分が1に近い階層的な構造を持つ（第9章）。
\begin{remark}
ニュートリノ質量を導入すると、パラメータ数は質量生成機構に依存する。
Dirac 質量の場合は、質量3個と PMNS 行列のパラメータ4個（混合角3、CP位相1）が加わり26個となる。
Majorana 質量の場合は、さらに2個の Majorana 位相が加わり28個となる（第11章）。
\end{remark}
%------------------------------------------------
\section*{第1章のまとめ}
%------------------------------------------------
\begin{itemize}
    \item 標準模型はゲージ群 $\SU(3)_C\times\SU(2)_L\times\U(1)_Y$ に基づく、
    繰り込み可能なカイラルゲージ理論である。
    \item 物質は3世代のクォークとレプトンからなり、
    1世代は15個の Weyl フェルミオン
    $Q_L,u_R,d_R,L_L,e_R$ で構成される。
    \item 左巻きと右巻きのフェルミオンは異なるゲージ表現に属するため、
    裸の質量項は禁止され、質量はヒッグス機構によって生成される。
    \item 電荷は $Q=T_3+Y/2$ で与えられ、
    超電荷の割り当ては異常相殺条件を満たす。
    \item ラグランジアンはゲージ・フェルミオン・ヒッグス・Yukawa の4部分（と $\theta$ 項）からなり、
    19個のパラメータで完全に指定される。
\end{itemize}
%================================================
\chapter{標準模型のゲージ構造}
%================================================
本章では、第1章で結果だけを示したゲージ部分の構造を、
非可換ゲージ理論の一般論から組み立て直す。
まず一般の Lie 群 $G$ に対するゲージ理論を復習し（\S\ref{sec:YM-review}）、
それを $\SU(2)_L$, $\U(1)_Y$ に適用する。
最後に、標準模型が\textbf{カイラル}なゲージ理論であることの意味を整理する。

対称性の破れ（ヒッグス機構）は本章では扱わず、
すべてのゲージボソンは massless のままである。
質量の生成は第II部で扱う。
%------------------------------------------------
\section{非可換ゲージ理論の復習}
\label{sec:YM-review}
%------------------------------------------------
\subsection{大域的対称性から局所対称性へ}
コンパクト Lie 群 $G$ の生成子を $T^a$（$a=1,\dots,\dim G$）とし、
\begin{equation}
[T^a,T^b]=if^{abc}T^c,
\qquad
\tr(T^aT^b)=\frac12\delta^{ab}
\end{equation}
とする（$T^a$ は表現 $R$ の行列であり、規格化は基本表現について述べている）。
表現 $R$ に属する $\dim R$ 成分の Dirac 場 $\psi$ の自由ラグランジアン
\begin{equation}
\Lag_0=\bar\psi\,i\gamma^\mu\partial_\mu\psi
\end{equation}
は、定数行列 $U=e^{i\alpha^aT^a}\in G$ による\term{大域的変換}{global transformation}
$\psi\to U\psi$ で不変である。

これを時空の各点で独立な変換
\begin{equation}
\psi(x)\longrightarrow\psi'(x)=U(x)\psi(x),
\qquad
U(x)=\exp\left(i\alpha^a(x)T^a\right)
\label{eq:local-transf}
\end{equation}
（\term{局所ゲージ変換}{local gauge transformation}）に拡張すると、
$\partial_\mu\psi\to U\partial_\mu\psi+(\partial_\mu U)\psi$ となり、
第2項のために $\Lag_0$ は不変でなくなる。
\subsection{共変微分}
\begin{mydef}{共変微分（covariant derivative）}{covariant-derivative}
Lie 代数に値をとるベクトル場 $A_\mu=A_\mu^aT^a$（\term{ゲージ場}{gauge field}）と
結合定数 $g$ を導入し、\term{共変微分}{covariant derivative}を
\begin{equation}
D_\mu=\partial_\mu-igA_\mu=\partial_\mu-igA_\mu^aT^a
\end{equation}
で定義する。ゲージ場は局所ゲージ変換の下で
\begin{equation}
A_\mu\longrightarrow A'_\mu
=UA_\mu U^{-1}-\frac ig(\partial_\mu U)U^{-1}
\label{eq:A-transf}
\end{equation}
と変換するものとする。
\end{mydef}
\begin{myprop}{共変微分の共変性}{Dcov}
\eqref{eq:A-transf} の下で
\begin{equation}
D_\mu\psi\longrightarrow (D_\mu\psi)'=U(x)\,D_\mu\psi
\end{equation}
が成り立つ。したがって $\bar\psi\,i\gamma^\mu D_\mu\psi$ は局所ゲージ不変である。
\end{myprop}
\begin{note}
$(D_\mu\psi)'=\partial_\mu(U\psi)-igA'_\mu U\psi$ を計算する：
\[
\partial_\mu(U\psi)=U\partial_\mu\psi+(\partial_\mu U)\psi,
\qquad
-igA'_\mu U\psi
=-igUA_\mu\psi-(\partial_\mu U)U^{-1}U\psi
=-igUA_\mu\psi-(\partial_\mu U)\psi .
\]
$(\partial_\mu U)\psi$ が打ち消し、$(D_\mu\psi)'=U(\partial_\mu-igA_\mu)\psi=UD_\mu\psi$。
また $\bar\psi\to\bar\psi U^\dagger=\bar\psi U^{-1}$ なので
$\bar\psi\gamma^\mu D_\mu\psi\to\bar\psi U^{-1}\gamma^\mu UD_\mu\psi=\bar\psi\gamma^\mu D_\mu\psi$
（$U$ は内部空間の行列で $\gamma^\mu$ と可換）。
\end{note}
\paragraph{無限小変換}
$\alpha^a$ が小さいとき $U\simeq1+i\alpha^aT^a$ であり、
\eqref{eq:A-transf} は成分で
\begin{equation}
\delta A_\mu^a
=\frac1g\partial_\mu\alpha^a+f^{abc}A_\mu^b\alpha^c
=\frac1g\left(\partial_\mu\alpha^a+gf^{abc}A_\mu^b\alpha^c\right)
\equiv\frac1g(D_\mu\alpha)^a
\label{eq:A-infinitesimal}
\end{equation}
となる。最後の $(D_\mu\alpha)^a$ は随伴表現での共変微分である。
\begin{note}
\eqref{eq:A-transf} を $\alpha$ の1次まで展開すると
\[
A'_\mu=(1+i\alpha)A_\mu(1-i\alpha)-\frac ig(i\partial_\mu\alpha)
=A_\mu+i[\alpha,A_\mu]+\frac1g\partial_\mu\alpha
\qquad(\alpha\equiv\alpha^aT^a).
\]
交換子は $i[\alpha,A_\mu]=i\alpha^cA_\mu^b[T^c,T^b]=i\cdot if^{cba}\alpha^cA_\mu^bT^a=f^{abc}A_\mu^b\alpha^cT^a$
（$-f^{cba}=f^{abc}$ を用いた）。
第1項 $f^{abc}A^b\alpha^c$ は、$A_\mu^a$ 自身が随伴表現の「荷」を持ち
ゲージ変換で回転することを表す。
可換群（$f^{abc}=0$）ではこの項が消え、$\delta A_\mu=\frac1g\partial_\mu\alpha$ となる。
\end{note}
随伴表現の生成子は $(T^a_{\mathrm{adj}})^{bc}=-if^{abc}$ なので、
$(D_\mu\alpha)^a=\partial_\mu\alpha^a-igA^b_\mu(T^b_{\mathrm{adj}})^{ac}\alpha^c$ と書けば、
基本表現の共変微分と同じ形をしていることがわかる。
一般に、表現 $R$ の場に対しては $D_\mu=\partial_\mu-igA_\mu^aT^a_R$ とすればよい。
\subsection{場の強さ}
\begin{mydef}{場の強さ（field strength）}{field-strength}
共変微分の交換子により、\term{場の強さ}{field strength}
$F_{\mu\nu}=F_{\mu\nu}^aT^a$ を
\begin{equation}
[D_\mu,D_\nu]=-igF_{\mu\nu}
\end{equation}
で定義する。具体的には
\begin{equation}
F_{\mu\nu}=\partial_\mu A_\nu-\partial_\nu A_\mu-ig[A_\mu,A_\nu],
\qquad
F_{\mu\nu}^a=\partial_\mu A_\nu^a-\partial_\nu A_\mu^a+gf^{abc}A_\mu^bA_\nu^c .
\label{eq:Fmunu}
\end{equation}
\end{mydef}
計算は第1章の注で行った。
$[D_\mu,D_\nu]$ が微分演算子でなく（$\psi$ の微分を含まない）
単なる掛け算の行列になることが重要である。
\begin{myprop}{場の強さの変換性}{F-transf}
局所ゲージ変換の下で
\begin{equation}
F_{\mu\nu}\longrightarrow UF_{\mu\nu}U^{-1},
\qquad\text{無限小では}\quad
\delta F^a_{\mu\nu}=f^{abc}F^b_{\mu\nu}\alpha^c .
\end{equation}
すなわち $F_{\mu\nu}$ は随伴表現として\textbf{共変的に}変換する（不変ではない）。
\end{myprop}
\begin{note}
命題~\ref{prop:Dcov} より $D'_\mu=UD_\mu U^{-1}$（演算子として）。よって
$[D'_\mu,D'_\nu]=U[D_\mu,D_\nu]U^{-1}$、すなわち $-igF'_{\mu\nu}=-igUF_{\mu\nu}U^{-1}$。
無限小形は $\delta A$ の第1項と同じ計算で $i[\alpha,F]$ から得られる。
可換群では $F'_{\mu\nu}=F_{\mu\nu}$ となり、電磁場と同様にゲージ不変である。
\end{note}
場の強さは \term{Bianchi 恒等式}{Bianchi identity}
\begin{equation}
D_\mu F_{\nu\rho}+D_\nu F_{\rho\mu}+D_\rho F_{\mu\nu}=0,
\qquad
D_\mu F_{\nu\rho}\equiv\partial_\mu F_{\nu\rho}-ig[A_\mu,F_{\nu\rho}]
\end{equation}
を満たす。これは Jacobi 恒等式 $[D_\mu,[D_\nu,D_\rho]]+\text{(cyclic)}=0$ から直ちに従い、
電磁気学の斉次 Maxwell 方程式 $\nabla\cdot\bm B=0$, $\nabla\times\bm E+\partial_t\bm B=0$ の一般化である。
\begin{remark}
幾何学的には、$A_\mu$ は主 $G$ 束の\term{接続}{connection}、
$F_{\mu\nu}$ はその\term{曲率}{curvature}であり、
物質場は同伴ベクトル束の切断である。
共変微分は「異なる点のファイバーを比較する」ための平行移動を与える。
有限区間の平行移動は \term{Wilson 線}{Wilson line}
$W[C]=\mathcal{P}\exp\left(ig\int_CA_\mu dx^\mu\right)$ で与えられ、
閉曲線についてのトレース $\tr W[C]$（Wilson ループ）はゲージ不変量となる。
\end{remark}
%------------------------------------------------
\section{$\SU(2)_L$ ゲージ対称性}
%------------------------------------------------
\subsection{$\SU(2)_L$ 二重項とゲージ変換}
$\SU(2)_L$ の生成子は $T^i=\sigma^i/2$、構造定数は $f^{ijk}=\epsilon^{ijk}$ である。
標準模型で $\SU(2)_L$ が非自明に作用するのは、
左巻き二重項
\begin{equation}
Q_{Li}=\begin{pmatrix}u_{Li}\\ d_{Li}\end{pmatrix},
\qquad
L_{Li}=\begin{pmatrix}\nu_{Li}\\ e_{Li}\end{pmatrix}
\qquad(i=1,2,3)
\end{equation}
とヒッグス二重項 $H$ のみである。これらは
\begin{equation}
\Psi_L\longrightarrow\exp\left(i\beta^i(x)\frac{\sigma^i}{2}\right)\Psi_L
\end{equation}
と変換し、右巻き場 $u_R$, $d_R$, $e_R$ は変換しない。
二重項の上成分と下成分が「回転」して混ざるので、
\textbf{$\SU(2)_L$ 対称性が破れていなければ $\nu_L$ と $e_L$、$u_L$ と $d_L$ は区別できない}。
実際に区別できるのは、電弱対称性が破れているからである。
\subsection{$W$ ボソンの場の強さ}
3つのゲージ場 $W_\mu^i$ は随伴表現 $\mathbf{3}$（三重項）に属し、
\eqref{eq:A-infinitesimal} より
\begin{equation}
\delta W_\mu^i=\frac1g\partial_\mu\beta^i+\epsilon^{ijk}W_\mu^j\beta^k
\end{equation}
と変換する。第2項は3次元ベクトルの無限小回転 $\delta\bm W=-\bm\beta\times\bm W$ である。
場の強さは
\begin{equation}
W_{\mu\nu}^i=\partial_\mu W_\nu^i-\partial_\nu W_\mu^i+g\epsilon^{ijk}W_\mu^jW_\nu^k,
\qquad
\text{ベクトル記法で}\quad
\bm W_{\mu\nu}=\partial_\mu\bm W_\nu-\partial_\nu\bm W_\mu+g\,\bm W_\mu\times\bm W_\nu .
\end{equation}
\subsection{荷電基底：$W^\pm$ と $W^3$}
$T^\pm=T^1\pm iT^2$ と
\begin{equation}
W_\mu^\pm=\frac{1}{\sqrt2}\left(W_\mu^1\mp iW_\mu^2\right)
\end{equation}
を用いると（第1章の注を参照）、
\begin{equation}
W_\mu^i\frac{\sigma^i}{2}
=\frac12\begin{pmatrix}
W_\mu^3 & W_\mu^1-iW_\mu^2\\
W_\mu^1+iW_\mu^2 & -W_\mu^3
\end{pmatrix}
=\frac12\begin{pmatrix}
W_\mu^3 & \sqrt2W_\mu^+\\
\sqrt2W_\mu^- & -W_\mu^3
\end{pmatrix}
\label{eq:Wmatrix}
\end{equation}
となる。非対角成分 $W^\pm$ は二重項の上下を入れ替え、
対角成分 $W^3$ は上下の成分に $\pm\frac12$ の重みで結合する。

$W^\pm$ は $T_3$ の固有値 $\pm1$ を持つ。実際
\begin{equation}
[T^3,T^\pm]=\pm T^\pm
\end{equation}
なので、$W^+$ は $T_3$ を $+1$ だけ上げる演算子 $T^+$ と組になる。
$W^\pm$ は $\U(1)_Y$ の荷を持たない（$Y=0$）ので、
電荷は $Q=T_3=\pm1$ となる。
\begin{note}
$W^3_{\mu\nu}$ を $W^\pm$ で書き直す。
$W^1=\frac{1}{\sqrt2}(W^++W^-)$, $W^2=\frac{i}{\sqrt2}(W^+-W^-)$ より
\[
\epsilon^{3jk}W^j_\mu W^k_\nu
=W^1_\mu W^2_\nu-W^2_\mu W^1_\nu
=-i\left(W^+_\mu W^-_\nu-W^-_\mu W^+_\nu\right).
\]
したがって
\[
W^3_{\mu\nu}=\partial_\mu W^3_\nu-\partial_\nu W^3_\mu-ig\left(W^+_\mu W^-_\nu-W^-_\mu W^+_\nu\right).
\]
$-\frac14W^3_{\mu\nu}W^{3\mu\nu}$ を展開すると $W^3W^+W^-$ の3点結合が現れる。
$W^3$ は対称性の破れの後に光子と $Z$ の線形結合になるので、
これが $WW\gamma$, $WWZ$ 結合の起源である（第8章）。
\end{note}
\subsection{二重項に対する共変微分}
\eqref{eq:Wmatrix} を用いると、超電荷 $Y$ の二重項 $\Psi_L=(\psi_1,\psi_2)^T$ に対して
\begin{equation}
D_\mu\Psi_L
=\begin{pmatrix}
\partial_\mu-\dfrac{i}{2}\left(gW_\mu^3+g'YB_\mu\right) & -\dfrac{ig}{\sqrt2}W_\mu^+\\[2mm]
-\dfrac{ig}{\sqrt2}W_\mu^- & \partial_\mu+\dfrac{i}{2}\left(gW_\mu^3-g'YB_\mu\right)
\end{pmatrix}
\begin{pmatrix}\psi_1\\ \psi_2\end{pmatrix}
\label{eq:D-doublet}
\end{equation}
となる（カラー部分は省略）。
\begin{note}
$-igW^i\frac{\sigma^i}{2}-ig'\frac Y2B$ の $(1,1)$ 成分は
$-\frac{ig}{2}W^3-\frac{ig'Y}{2}B$、$(2,2)$ 成分は $+\frac{ig}{2}W^3-\frac{ig'Y}{2}B$、
$(1,2)$ 成分は $-ig\cdot\frac{\sqrt2}{2}W^+=-\frac{ig}{\sqrt2}W^+$。
\end{note}
\subsection{$\SU(2)_L$ の相互作用の構造}
$\SU(2)_L$ は単純群なので結合定数は $g$ ただ1つである。
したがって、すべての二重項（3世代のクォーク・レプトン、ヒッグス）は
同じ強さ $g$ で $W$ ボソンと結合する。
これが弱い相互作用の\term{普遍性}{universality}の起源である。
例えばミューオン崩壊 $\mu\to e\bar\nu_e\nu_\mu$ と
原子核の $\beta$崩壊 $d\to ue\bar\nu_e$ の結合の強さがほぼ等しいことは、
この普遍性の帰結である（クォークの場合の小さなずれは CKM 混合による。第9章）。

非可換ゲージ理論では、ゲージ場の規格化（$-\frac14W^iW^i$ の係数）を固定すると、
生成子の交換関係 $[T^i,T^j]=i\epsilon^{ijk}T^k$ によって
荷（$T_3$ の固有値）の大きさまで決まってしまう。
二重項の $T_3=\pm\frac12$、三重項の $T_3=0,\pm1$ のように、
荷は量子化されている。
%------------------------------------------------
\section{$\U(1)_Y$ とハイパーチャージ}
%------------------------------------------------
\subsection{可換ゲージ理論}
$\U(1)_Y$ の変換は、超電荷 $Y$ の場に対して
\begin{equation}
\psi\longrightarrow e^{i\theta(x)Y/2}\psi,
\qquad
B_\mu\longrightarrow B_\mu+\frac{1}{g'}\partial_\mu\theta
\end{equation}
である。構造定数がゼロなので
\begin{equation}
B_{\mu\nu}=\partial_\mu B_\nu-\partial_\nu B_\mu
\end{equation}
はそれ自体ゲージ不変であり、$B_\mu$ は自己相互作用を持たない。
$B_\mu$ は電磁場 $A_\mu$ と同様の可換ゲージ場であるが、
光子そのものではない（光子は $W^3$ と $B$ の線形結合である）。
\subsection{超電荷は量子化されない}
非可換群と異なり、$\U(1)$ の生成子は交換関係による制約を受けない。
$\U(1)$ の既約表現は任意の実数 $q$ に対して $\psi\to e^{iq\theta}\psi$ で与えられ、
\textbf{原理的には任意の実数値の超電荷が許される}。
例えば $Y=\sqrt2$ の粒子を導入しても、ゲージ不変性とは矛盾しない。

それにもかかわらず、観測される電荷は
\begin{equation}
Q(e)=-1,\qquad Q(u)=+\frac23,\qquad Q(d)=-\frac13
\end{equation}
のように $1/3$ の整数倍に量子化されている。特に
\begin{equation}
Q(p)+Q(e)=2Q(u)+Q(d)+Q(e)=0
\end{equation}
は $10^{-21}$ 程度の精度で成り立っている。
標準模型の枠内では、この量子化はゲージ異常の相殺条件（第3章）と
Yukawa 結合の存在から説明される。
大統一理論では、$Y$ が非可換群の生成子の一部となることでより自然に説明される（第14章）。
\subsection{超電荷の決め方}
各場の超電荷は、Gell-Mann--西島--中野の関係式 $Q=T_3+Y/2$ を
観測された電荷に当てはめて決める：
\begin{itemize}
    \item 一重項（$T_3=0$）では $Y=2Q$。
    例：$Y(e_R)=-2$, $Y(u_R)=\frac43$, $Y(d_R)=-\frac23$。
    \item 二重項では、上下の成分で $Q$ の差が1でなければならず、$Y$ は平均電荷の2倍。
    例：$Y(L_L)=Q(\nu)+Q(e)=-1$, $Y(Q_L)=Q(u)+Q(d)=\frac13$。
\end{itemize}
二重項の上下で電荷の差が1であることは、
$W^\pm$ が電荷 $\pm1$ を持つことと整合している。
\begin{remark}
$\U(1)$ 因子が複数ある理論では、場の強さ同士の
\term{運動項混合}{kinetic mixing} $\epsilon\,B_{\mu\nu}X^{\mu\nu}$ がゲージ不変に書けるが、
標準模型の $\U(1)$ は $\U(1)_Y$ の1つだけなので、この問題は生じない。
\end{remark}
%------------------------------------------------
\section{共変微分とゲージ場}
%------------------------------------------------
\subsection{直積群の共変微分}
標準模型のゲージ群は3つの因子の直積なので、
異なる因子の生成子は互いに可換である：
\begin{equation}
[T^a_{(3)},T^i_{(2)}]=[T^a_{(3)},Y]=[T^i_{(2)},Y]=0 .
\end{equation}
このため、各因子のゲージ場をそれぞれ独立に導入し、共変微分を
\begin{equation}
\boxed{
D_\mu
=
\partial_\mu
-ig_sG_\mu^aT^a
-igW_\mu^i\frac{\sigma^i}{2}
-ig'\frac{Y}{2}B_\mu
}
\label{eq:SM-covariant-derivative}
\end{equation}
とすればよい。ゲージ変換はそれぞれの因子ごとに \eqref{eq:A-transf} に従う：
\begin{align}
G_\mu^aT^a&\longrightarrow U_3\,G_\mu^aT^a\,U_3^{-1}-\frac{i}{g_s}(\partial_\mu U_3)U_3^{-1},\\
W_\mu^i\frac{\sigma^i}{2}&\longrightarrow U_2\,W_\mu^i\frac{\sigma^i}{2}\,U_2^{-1}-\frac{i}{g}(\partial_\mu U_2)U_2^{-1},\\
B_\mu&\longrightarrow B_\mu+\frac{1}{g'}\partial_\mu\theta .
\end{align}
ゲージ場 $G_\mu^a$ は $\SU(2)_L\times\U(1)_Y$ の下で不変、
$W_\mu^i$ は $\SU(3)_C\times\U(1)_Y$ の下で不変である。
\begin{note}
$U=U_3U_2U_1$ と書く。$U_3,U_2,U_1$ は互いに可換であり、
例えば $W$ の項は $U_3$, $U_1$ と可換なので
$U(W^i\sigma^i/2)U^{-1}=U_2(W^i\sigma^i/2)U_2^{-1}$。
また $(\partial_\mu U)U^{-1}=(\partial_\mu U_3)U_3^{-1}+(\partial_\mu U_2)U_2^{-1}+(\partial_\mu U_1)U_1^{-1}$
となり、各因子の非斉次項が独立に現れる。
したがって3つのゲージ場はそれぞれ自分の因子の変換だけを補償すればよい。
\end{note}
各因子の結合定数 $g_s$, $g$, $g'$ は互いに独立である。
$G_{\SM}$ は単純群ではないので、ゲージ対称性だけからは
結合定数の比は決まらない。
これが Weinberg 角 $\tan\theta_W=g'/g$ が自由パラメータとなる理由である。
\subsection{各場への作用}
共変微分 \eqref{eq:SM-covariant-derivative} は、場の表現に応じて作用する
（すべての場について書き下した形は第1章 \S\ref{sec:SM-lagrangian} を参照）。
例えば $L_L$ に対しては \eqref{eq:D-doublet} で $Y=-1$ とおいて
\begin{equation}
D_\mu L_L=
\begin{pmatrix}
\partial_\mu-\dfrac{i}{2}\left(gW_\mu^3-g'B_\mu\right) & -\dfrac{ig}{\sqrt2}W_\mu^+\\[2mm]
-\dfrac{ig}{\sqrt2}W_\mu^- & \partial_\mu+\dfrac{i}{2}\left(gW_\mu^3+g'B_\mu\right)
\end{pmatrix}
\begin{pmatrix}\nu_L\\ e_L\end{pmatrix}
\end{equation}
となる。ニュートリノ成分には $gW^3-g'B$ の組み合わせが現れる。
後で見るように、これは $Z$ ボソンに比例する組み合わせ
（$gW^3-g'B=\sqrt{g^2+g'^2}\,Z$）であり、
\textbf{ニュートリノが光子と結合しない}ことが自動的に保証される（第6章）。
\begin{note}
第1章の \eqref{eq:ZA-preview} より
$Z=\cos\theta_WW^3-\sin\theta_WB=(gW^3-g'B)/\sqrt{g^2+g'^2}$。
ニュートリノは $Q=0$ なので光子と結合しないのは当然であるが、
これが共変微分の形から見えるのは、$\theta_W$ が $Q=T_3+Y/2$ と整合するように定義されているからである。
\end{note}
\subsection{共役な場の共変微分}
$\psi$ が $D_\mu\psi$ の形で変換するとき、エルミート共役 $\psi^\dagger$ に対しては
\begin{equation}
(D_\mu\psi)^\dagger=\partial_\mu\psi^\dagger+ig\psi^\dagger A_\mu
\end{equation}
となる（$A_\mu$ はエルミート）。ヒッグス場の運動項
\begin{equation}
(D_\mu H)^\dagger(D^\mu H)
\end{equation}
はこの組み合わせで書かれている。
共役場 $\tilde H=i\sigma^2H^*$ は $Y=-1$ の二重項なので、
$D_\mu\tilde H$ は \eqref{eq:D-doublet} で $Y=-1$ とおいたものになり、
$D_\mu\tilde H=i\sigma^2(D_\mu H)^*$ が成り立つ。
\begin{note}
$(D_\mu H)^*=\left(\partial_\mu+igW^i\frac{(\sigma^i)^*}{2}+\frac{ig'}{2}B_\mu\right)H^*$。
左から $i\sigma^2$ を掛け、$\sigma^2(\sigma^i)^*=-\sigma^i\sigma^2$（擬実性 \eqref{eq:pseudoreal}）を用いると
\[
i\sigma^2(D_\mu H)^*
=\left(\partial_\mu-igW^i\frac{\sigma^i}{2}+\frac{ig'}{2}B_\mu\right)(i\sigma^2H^*)
\]
となり、これは $Y=-1$ の二重項に対する共変微分である。
\end{note}
%------------------------------------------------
\section{Yang--Mills ラグランジアン}
%------------------------------------------------
\subsection{ゲージ不変な運動項}
\begin{myprop}{Yang--Mills ラグランジアン}{YM}
\begin{equation}
\boxed{
\Lag_{\mathrm{YM}}
=-\frac12\tr\left(F_{\mu\nu}F^{\mu\nu}\right)
=-\frac14F^a_{\mu\nu}F^{a\mu\nu}
}
\end{equation}
は局所ゲージ不変である。
\end{myprop}
\begin{note}
$F\to UFU^{-1}$ とトレースの巡回性より
$\tr(F'F')=\tr(UFU^{-1}UFU^{-1})=\tr(FF)$。
成分では $\tr(T^aT^b)=\frac12\delta^{ab}$ より
$\tr(F_{\mu\nu}F^{\mu\nu})=\frac12F^a_{\mu\nu}F^{a\mu\nu}$。
\end{note}
係数 $-\frac14$ は、$g\to0$ で各成分 $A^a_\mu$ が
Maxwell 理論と同じ規格化の運動項
\begin{equation}
-\frac14(\partial_\mu A_\nu^a-\partial_\nu A_\mu^a)^2
=\frac12\left(\bm E^a\cdot\bm E^a-\bm B^a\cdot\bm B^a\right)
\end{equation}
を持つように選んでいる。
この規格化（\term{正準規格化}{canonical normalization}）では、
結合定数 $g$ は共変微分の中に現れる。
$A_\mu\to A_\mu/g$ と再定義すると $\Lag_{\mathrm{YM}}=-\frac{1}{4g^2}F^aF^a$ となり、
$g$ が運動項の前に現れる流儀もある。

ゲージ場の質量項
\begin{equation}
\frac12M^2A^a_\mu A^{a\mu}
\end{equation}
は \eqref{eq:A-infinitesimal} の $\frac1g\partial_\mu\alpha^a$ のためにゲージ不変ではない。
したがって\textbf{ゲージ対称性が破れていない限りゲージボソンは massless である}。
$W^\pm$, $Z^0$ の質量を説明するには、ヒッグス機構が必要となる（第5章、第6章）。
\subsection{自己相互作用}
$F^a_{\mu\nu}$ が $A$ について非線形であるため、$\Lag_{\mathrm{YM}}$ は
\begin{equation}
\Lag_{\mathrm{YM}}
=\underbrace{-\frac14(\partial_\mu A_\nu^a-\partial_\nu A_\mu^a)^2}_{\text{運動項}}
\underbrace{-gf^{abc}(\partial_\mu A_\nu^a)A^{b\mu}A^{c\nu}}_{\text{3点結合}}
\underbrace{-\frac{g^2}{4}f^{abe}f^{cde}A_\mu^aA_\nu^bA^{c\mu}A^{d\nu}}_{\text{4点結合}}
\end{equation}
となる（展開は第1章の注を参照）。
3点結合は $g$、4点結合は $g^2$ に比例し、
その相対的な大きさはゲージ対称性によって固定されている。
Feynman 則は付録Fにまとめる。

これらの自己相互作用は、ゲージ場自身が荷を持つ
（随伴表現に属する）ことの帰結である。
光子は電荷を持たないので光子同士は直接結合しないが、
グルーオンはカラー荷を、$W$ ボソンは弱アイソスピンを持つ。
この自己相互作用が、QCD の漸近的自由性（$\beta$ 関数の $-\frac{11}{3}N$）の起源である。
\subsection{運動方程式}
\begin{myprop}{Yang--Mills 方程式}{YM-eq}
\begin{equation}
\Lag=-\frac14F^a_{\mu\nu}F^{a\mu\nu}+\bar\psi\,i\gamma^\mu D_\mu\psi
\end{equation}
の $A_\nu^a$ に関する運動方程式は
\begin{equation}
(D_\mu F^{\mu\nu})^a
\equiv\partial_\mu F^{a\mu\nu}+gf^{abc}A^b_\mu F^{c\mu\nu}
=-g\,j^{a\nu},
\qquad
j^{a\nu}=\bar\psi\gamma^\nu T^a\psi
\end{equation}
である。
\end{myprop}
\begin{note}
フェルミオン部分は $\bar\psi i\gamma^\mu(-igA^a_\mu T^a)\psi=g\,j^{a\mu}A^a_\mu$ なので、
$\partial\Lag/\partial A^a_\nu\supset g\,j^{a\nu}$。
Yang--Mills 部分は
$\partial\Lag/\partial(\partial_\mu A^a_\nu)=-F^{a\mu\nu}$、
および非線形項から
\[
\frac{\partial\Lag_{\mathrm{YM}}}{\partial A^a_\nu}
=-\frac12F^{c\rho\sigma}\frac{\partial F^c_{\rho\sigma}}{\partial A^a_\nu}
=-gf^{cae}F^{c\nu\sigma}A^e_\sigma .
\]
Euler--Lagrange 方程式 $\partial_\mu(-F^{a\mu\nu})=-gf^{cae}F^{c\nu\sigma}A^e_\sigma+g\,j^{a\nu}$ を整理し、
$f^{cae}F^{c\nu\sigma}A^e_\sigma=-f^{abc}A^b_\mu F^{c\mu\nu}$（添字の付け替えと反対称性）を用いれば得られる。
\end{note}
左辺を $\partial_\mu F^{a\mu\nu}=-g\left(j^{a\nu}+f^{abc}A^b_\mu F^{c\mu\nu}\right)$ と書き直すと、
右辺の括弧はゲージ場自身の寄与を含む全カレントであり、
$\partial_\mu F^{a\mu\nu}$ の反対称性からこれは通常の意味で保存する。
一方、物質のカレント $j^{a\nu}$ だけでは保存せず、
\textbf{共変的に}保存する（$(D_\nu j^\nu)^a=0$）。
これも「ゲージ場が荷を持つ」ことの表れである。
\subsection{標準模型のゲージ部分}
標準模型のゲージ部分は3つの Yang--Mills 項の和
\begin{equation}
\Lag_{\mathrm{gauge}}
=-\frac14G^a_{\mu\nu}G^{a\mu\nu}
-\frac14W^i_{\mu\nu}W^{i\mu\nu}
-\frac14B_{\mu\nu}B^{\mu\nu}
\end{equation}
である。これに加えて、ゲージ不変かつ繰り込み可能な項として
\begin{equation}
\theta_3\frac{g_s^2}{32\pi^2}G^a_{\mu\nu}\tilde G^{a\mu\nu},
\qquad
\theta_2\frac{g^2}{32\pi^2}W^i_{\mu\nu}\tilde W^{i\mu\nu},
\qquad
\theta_1\frac{g'^2}{32\pi^2}B_{\mu\nu}\tilde B^{\mu\nu}
\end{equation}
が書けるが、いずれも全微分である。
$\theta_1$ は物理的効果を持たず、$\theta_2$ は場の再定義で消去でき、
$\theta_3$ のみが物理的意味を持つ（第1章、第14章）。

量子化にはゲージ固定が必要であり、それに伴って Faddeev--Popov ゴーストが現れる。
これについては付録Eで扱う。
%------------------------------------------------
\section{電弱相互作用のゲージ構造}
%------------------------------------------------
\subsection{電弱カレント}
$\SU(2)_L\times\U(1)_Y$ 部分のフェルミオンの相互作用は、
共変微分を展開して
\begin{equation}
\Lag_{\mathrm{EW}}
=g\,W^i_\mu J^{i\mu}+g'B_\mu J^\mu_Y
\end{equation}
と書ける。ここで
\begin{equation}
J^{i\mu}=\sum_{\Psi_L}\overline{\Psi_L}\gamma^\mu\frac{\sigma^i}{2}\Psi_L,
\qquad
J^\mu_Y=\sum_\psi\bar\psi\gamma^\mu\frac{Y}{2}\psi
\end{equation}
はそれぞれ\term{弱アイソスピンカレント}{weak isospin current}、
\term{超電荷カレント}{hypercharge current}である
（和は左巻き二重項およびすべての Weyl 場についてとる）。

荷電基底に移ると、\eqref{eq:Wmatrix} より
\begin{equation}
gW^i_\mu J^{i\mu}
=\frac{g}{\sqrt2}\left(W^+_\mu J^{+\mu}+W^-_\mu J^{-\mu}\right)+gW^3_\mu J^{3\mu},
\qquad
J^{+\mu}=\sum\overline{\Psi_L}\gamma^\mu T^+\Psi_L
\end{equation}
となる。1世代分を具体的に書くと
\begin{align}
J^{+\mu}&=\overline{\nu_L}\gamma^\mu e_L+\overline{u_L}\gamma^\mu d_L,
\qquad J^{-\mu}=(J^{+\mu})^\dagger,\\
J^{3\mu}&=\frac12\left(\overline{\nu_L}\gamma^\mu\nu_L-\overline{e_L}\gamma^\mu e_L
+\overline{u_L}\gamma^\mu u_L-\overline{d_L}\gamma^\mu d_L\right)
\end{align}
である（クォークの世代間混合は第9章で扱う）。
\subsection{電磁カレントの埋め込み}
$Q=T_3+Y/2$ をカレントで書くと
\begin{equation}
J^\mu_{\mathrm{EM}}=J^{3\mu}+J^\mu_Y,
\qquad
J^\mu_{\mathrm{EM}}=\sum_\psi Q_\psi\,\bar\psi\gamma^\mu\psi .
\end{equation}
電磁カレントは $\SU(2)_L$ と $\U(1)_Y$ のカレントの和として埋め込まれている。
中性部分を
\begin{equation}
W^3_\mu=\cos\theta_WZ_\mu+\sin\theta_WA_\mu,
\qquad
B_\mu=-\sin\theta_WZ_\mu+\cos\theta_WA_\mu
\end{equation}
で書き直すと、次の結果を得る。
\begin{myprop}{中性カレント相互作用}{neutral-current}
$\tan\theta_W=g'/g$ とすると
\begin{equation}
gW^3_\mu J^{3\mu}+g'B_\mu J^\mu_Y
=e\,A_\mu J^\mu_{\mathrm{EM}}
+\frac{g}{\cos\theta_W}Z_\mu\left(J^{3\mu}-\sin^2\theta_W J^\mu_{\mathrm{EM}}\right),
\qquad
e=g\sin\theta_W .
\end{equation}
\end{myprop}
\begin{note}
$J_Y=J_{\mathrm{EM}}-J^3$ を代入する。$s\equiv\sin\theta_W$, $c\equiv\cos\theta_W$ とおく。

$A_\mu$ の係数：$gs\,J^3+g'c\,(J_{\mathrm{EM}}-J^3)$。
$gs=g'c=gg'/\sqrt{g^2+g'^2}=e$ なので $J^3$ が打ち消し、$e\,J_{\mathrm{EM}}$。

$Z_\mu$ の係数：$gc\,J^3-g's\,(J_{\mathrm{EM}}-J^3)=(gc+g's)J^3-g's\,J_{\mathrm{EM}}$。
$gc+g's=(g^2+g'^2)/\sqrt{g^2+g'^2}=g/c$、$g's=g\,s^2/c$ より
$\frac gc\left(J^3-s^2J_{\mathrm{EM}}\right)$。
\end{note}
この命題は純粋にゲージ場の基底の取り替えであり、
対称性の破れとは独立に成り立つ。
ただし $A_\mu$ と $Z_\mu$ が質量の固有状態（物理的な粒子）になることは、
ヒッグス機構によって初めて示される（第6章）。
$Z$ ボソンの結合 $T_3-Q\sin^2\theta_W$ の現象論は第8章で扱う。
\subsection{電弱「統一」の意味}
$\SU(2)_L\times\U(1)_Y$ は単純群ではないので、
電弱理論は2つの独立な結合定数 $g$, $g'$ を持ち、
厳密な意味での「統一」ではない。
この理論の本質は、
\begin{itemize}
    \item 電磁相互作用と弱い相互作用が、1つのゲージ群の異なる部分として記述されること
    \item 電荷 $Q$ が $\SU(2)_L$ と $\U(1)_Y$ の生成子の組み合わせとして埋め込まれていること
    \item その結果、$e$, $\sin\theta_W$, $M_W$, $M_Z$, $G_F$ などの間に関係が生じること
\end{itemize}
にある。例えば中性カレント（$Z$ 交換）の存在とその強さは、
荷電カレント（$W$ 交換）の強さと $\sin\theta_W$ から予言され、
1973年の発見と整合した。
%------------------------------------------------
\section{カイラルゲージ理論としての標準模型}
%------------------------------------------------
\subsection{ベクトル的理論とカイラル理論}
すべてのフェルミオンを左巻き Weyl 場で表したとき（第1章の左巻き表示）、
フェルミオン全体が属するゲージ群の表現を $R_{\mathrm{tot}}$ とする。
\begin{mydef}{ベクトル的理論とカイラル理論}{chiral-theory}
$R_{\mathrm{tot}}$ が実表現（または擬実表現）、すなわち $R_{\mathrm{tot}}\simeq\bar R_{\mathrm{tot}}$ のとき、
理論は\term{ベクトル的}{vector-like}であるという。
そうでないとき（$R_{\mathrm{tot}}$ が複素表現のとき）、理論は\term{カイラル}{chiral}であるという。
\end{mydef}
Dirac 場 $\psi$ の言葉では、$\psi_L$ と $\psi_R$ が同じ表現 $R$ に属するとき、
左巻き表示では $\psi_L\in R$, $(\psi_R)^c\in\bar R$ となり、$R\oplus\bar R$ は実なので
ベクトル的である。このとき質量項 $m\bar\psi\psi$ はゲージ不変に書ける。
\begin{myex}{QCD と QED はベクトル的}{vectorlike}
QCD では各クォークについて $q_L,q_R\in\mathbf 3$ なので、
左巻き表示で $\mathbf 3\oplus\bar{\mathbf 3}$ となりベクトル的である。
したがって $\SU(3)_C$ だけを考えれば、クォーク質量 $m_q\bar qq$ はゲージ不変である。
QED でも $e_L$, $e_R$ はともに $Q=-1$ なのでベクトル的である。
\end{myex}
\begin{myex}{標準模型はカイラル}{SM-chiral}
1世代の左巻き表示
\begin{equation}
R_{\mathrm{tot}}
=(\mathbf3,\mathbf2)_{1/3}
\oplus(\bar{\mathbf3},\mathbf1)_{-4/3}
\oplus(\bar{\mathbf3},\mathbf1)_{2/3}
\oplus(\mathbf1,\mathbf2)_{-1}
\oplus(\mathbf1,\mathbf1)_{2}
\end{equation}
の複素共役は
\begin{equation}
\bar R_{\mathrm{tot}}
=(\bar{\mathbf3},\mathbf2)_{-1/3}
\oplus(\mathbf3,\mathbf1)_{4/3}
\oplus(\mathbf3,\mathbf1)_{-2/3}
\oplus(\mathbf1,\mathbf2)_{1}
\oplus(\mathbf1,\mathbf1)_{-2}
\end{equation}
であり（$\bar{\mathbf2}\simeq\mathbf2$ を用いた）、
例えば $(\mathbf3,\mathbf2)_{1/3}$ は $\bar R_{\mathrm{tot}}$ に含まれない。
したがって $R_{\mathrm{tot}}\not\simeq\bar R_{\mathrm{tot}}$ であり、標準模型はカイラルである。
\end{myex}
\subsection{カイラル性の帰結}
標準模型がカイラルであることから、次の重要な帰結が得られる。
\paragraph{1. 裸の質量項の禁止}
ゲージ不変な質量項は、表現 $R$ と $\bar R$ の組からしか作れない。
カイラルな理論ではすべてのフェルミオンについてこの組が欠けているため、
\textbf{すべてのフェルミオンが massless でなければならない}。
質量を得るには、ヒッグス場との結合（Yukawa 結合）を通じて対称性を破る必要がある（第3章、第7章）。
\paragraph{2. 破れずに残る対称性はベクトル的}
対称性の破れの後に残る $\SU(3)_C\times\U(1)_{\mathrm{EM}}$ の下では、
標準模型のフェルミオンはベクトル的である（例~\ref{ex:vectorlike}）。
フェルミオンが質量を持てるのはこのためである。
逆に、もしフェルミオンが $\U(1)_{\mathrm{EM}}$ の下でもカイラルであったなら、
電磁対称性を破らずにそのフェルミオンに質量を与えることはできない。
\paragraph{3. フェルミオン質量の上限}
フェルミオン質量は $m_f=y_fv/\sqrt2$ と電弱スケールに比例する。
摂動論が有効な範囲（$y_f\lesssim\sqrt{4\pi}$）では
$m_f\lesssim\mathcal O(\mathrm{TeV})$ となる。
標準模型のフェルミオンが電弱スケールより重くなれないのは、カイラル性の帰結である。
\paragraph{4. 量子補正に対する保護}
$y_f\to0$ の極限では理論に
\term{カイラル対称性}{chiral symmetry} $\psi_L\to e^{i\alpha}\psi_L$ が現れる。
このためフェルミオン質量の量子補正は $m_f$ 自身に比例し
（$\delta m_f\propto m_f\ln\Lambda$）、
カットオフ $\Lambda$ に対して対数的にしか依存しない。
これを\term{技術的自然さ}{technical naturalness}と呼ぶ。
ヒッグス質量にはこのような保護がなく、$\delta m_h^2\propto\Lambda^2$ となることが
階層性問題の根源である（第14章）。
\paragraph{5. ゲージ異常}
カイラルなフェルミオンはゲージ異常を生じ得る。
ベクトル的な理論では異常は自動的に相殺するが、
カイラルな理論では表現の選び方に非自明な条件が課される。
標準模型ではこの条件が1世代ごとに満たされている（第3章）。
\paragraph{6. パリティの破れ}
$\psi_L$ と $\psi_R$ が異なる表現に属するので、
これらを入れ替えるパリティ変換はゲージ相互作用の対称性ではない。
これが弱い相互作用におけるパリティの最大の破れの起源である。
\begin{remark}
カイラルゲージ理論は、摂動論を超えた定義に技術的な困難を伴う。
例えば格子上にカイラルフェルミオンを置くと、
Nielsen--二宮の定理によって必ず逆のカイラリティを持つ
\term{ダブラー}{doubler}が現れる。
これは格子 QCD（ベクトル的理論）では回避できるが、
標準模型の電弱部分を格子上で非摂動的に定義することは現在でも未解決の問題である。
\end{remark}
%------------------------------------------------
\section*{第2章のまとめ}
%------------------------------------------------
\begin{itemize}
    \item 局所ゲージ不変性は、ゲージ場 $A_\mu$ と共変微分 $D_\mu=\partial_\mu-igA_\mu$ の導入を要求し、
    ゲージ場とフェルミオンの相互作用の形を一意に決める。
    \item 場の強さ $F_{\mu\nu}$ は随伴表現で共変的に変換し、
    $-\frac14F^aF^a$ が唯一のゲージ不変な（繰り込み可能な）運動項である。
    非可換ゲージ場は自己相互作用を持つ。
    \item ゲージボソンの質量項はゲージ不変ではない。
    \item $\SU(2)_L$ は左巻き二重項にのみ作用し、単一の結合定数 $g$ が弱い相互作用の普遍性を与える。
    $\U(1)_Y$ の超電荷は群論からは量子化されない。
    \item 電磁カレントは $J_{\mathrm{EM}}=J^3+J_Y$ として埋め込まれ、
    基底の回転により $eA_\mu J_{\mathrm{EM}}$ と $Z$ の中性カレント結合が得られる。
    \item 標準模型はカイラルなゲージ理論であり、これが
    裸の質量項の禁止、パリティの破れ、異常相殺条件の起源である。
\end{itemize}
%================================================
\chapter{フェルミオンの表現とゲージ不変性}
%================================================
第1章では標準模型のフェルミオンの量子数を表として与え、
第2章ではゲージ場の側の構造を調べた。
本章では物質場（フェルミオン）の側に立ち返り、
\begin{itemize}
    \item フェルミオン場がどのような表現に属し（\S\ref{sec:fermion-rep}--\S\ref{sec:quantum-numbers}）、
    \item その表現からどのような項がゲージ不変に書けるか（\S\ref{sec:invariant-terms}, \S\ref{sec:no-bare-mass}）、
    \item 表現の選び方が量子論的な整合性からどう制限されるか（\S\ref{sec:anomaly}）
\end{itemize}
を議論する。
最後の点は、超電荷の値が（規格化を除いて）\textbf{理論の整合性から決まってしまう}という
著しい結果につながる。
%------------------------------------------------
\section{フェルミオン場の表現}
\label{sec:fermion-rep}
%------------------------------------------------
\subsection{場の添字構造}
標準模型のフェルミオン場は、
Lorentz 群・ゲージ群・世代（フレーバー）の3種類の添字を持つ。
例えば左巻きクォーク二重項は
\begin{equation}
Q_{L\,\alpha a i}^{\ s}
\qquad
\begin{cases}
s=1,2 & \text{Weyl スピノルの添字（Lorentz 群 $(\tfrac12,0)$）}\\
\alpha=1,2,3 & \text{カラーの添字（$\SU(3)_C$ の $\mathbf3$）}\\
a=1,2 & \text{弱アイソスピンの添字（$\SU(2)_L$ の $\mathbf2$）}\\
i=1,2,3 & \text{世代の添字}
\end{cases}
\end{equation}
を持つ。すなわち $Q_L$ は
\begin{equation}
\text{Lorentz}\ (\tfrac12,0)\ \otimes\ \SU(3)_C\ \mathbf3\ \otimes\ \SU(2)_L\ \mathbf2\ \otimes\ \U(1)_Y\ (Y=\tfrac13)\ \otimes\ \text{世代}\ \mathbf3
\end{equation}
の表現空間の元である。
ゲージ変換は $\alpha,a$ にのみ作用し、Lorentz 変換は $s$ にのみ作用する。
世代の添字にはゲージ変換は作用しない（世代はゲージ量子数では区別されない）。
以下では通常、スピノルとカラーの添字は省略する。
\subsection{Weyl スピノルと荷電共役}
第1章で見たように、Dirac 場は2つの Weyl 場 $\psi_L\in(\frac12,0)$, $\psi_R\in(0,\frac12)$ からなる。
Lorentz 群の表現として $(\frac12,0)$ と $(0,\frac12)$ は互いに複素共役の関係にある。
したがって右巻き場の複素共役をとれば左巻き場が得られる。
これを Dirac 場の言葉で表すのが\term{荷電共役}{charge conjugation}である。
\begin{mydef}{荷電共役場}{charge-conj}
Dirac 場 $\psi$ の荷電共役場を
\begin{equation}
\psi^c\equiv C\bar\psi^T=i\gamma^2\psi^*
\end{equation}
で定義する。ここで $C=i\gamma^2\gamma^0$（Weyl 表示）である。
\end{mydef}
\begin{myprop}{荷電共役とカイラリティ}{chiral-conj}
荷電共役はカイラリティを反転する：
\begin{equation}
(\psi_R)^c=\PL\,\psi^c,
\qquad
(\psi_L)^c=\PR\,\psi^c .
\end{equation}
また、$\psi$ がゲージ群の表現 $R$ に属し $\psi\to U\psi$ と変換するとき、
$\psi^c\to U^*\psi^c$ であり、$\psi^c$ は複素共役表現 $\bar R$ に属する。
特に $\U(1)$ 電荷は符号が反転する。
\end{myprop}
\begin{note}
Weyl 表示では $\gamma^0$ は実対称なので $C\bar\psi^T=i\gamma^2\gamma^0(\gamma^0)^T\psi^*=i\gamma^2\psi^*$。
また $\gamma^5=\diag(-1,1)$ は実なので $(\PR\psi)^*=\PR\psi^*$。
$\gamma^2$ と $\gamma^5$ は反可換なので $\gamma^2\PR=\PL\gamma^2$。よって
\[
(\psi_R)^c=i\gamma^2(\PR\psi)^*=i\gamma^2\PR\psi^*=\PL\,i\gamma^2\psi^*=\PL\psi^c .
\]
ゲージ変換については、$U$ は内部空間に作用し $\gamma^2$ と可換なので
$(U\psi)^c=i\gamma^2U^*\psi^*=U^*\psi^c$。
$\U(1)$ では $U=e^{iq\theta}$ に対し $U^*=e^{-iq\theta}$ である。
\end{note}
この性質により、\textbf{すべてのフェルミオンを左巻き Weyl 場で表す}ことができる。
右巻き場 $\psi_R\in R$ の代わりに、左巻き場 $(\psi_R)^c\in\bar R$ を用いればよい。
第1章の左巻き表示
\begin{equation}
Q,\quad u^c\equiv(u_R)^c,\quad d^c\equiv(d_R)^c,\quad L,\quad e^c\equiv(e_R)^c
\end{equation}
はこの書き方である。
左巻き表示では、ゲージ不変性・異常相殺などの条件が
「左巻き Weyl 場の集合」についての条件として統一的に書ける。
\begin{remark}
Lorentz 不変な2次形式は、2成分記法で
$\chi\psi\equiv\chi^s\psi_s=\epsilon^{st}\chi_t\psi_s$（左巻き同士、Lorentz スカラー）と
$\chi^\dagger\bar\sigma^\mu\psi$（左巻きとその共役、Lorentz ベクトル）の2種類である。
Dirac 記法の $\bar\psi_L\psi_R$ は前者に、$\bar\psi_L\gamma^\mu\psi_L$ は後者に対応する。
本ノートでは主に Dirac 記法を用い、必要に応じて左巻き表示に言い換える。
\end{remark}
\subsection{不変テンソル}
ゲージ不変な項を作るには、表現の積の中から一重項 $\mathbf1$ を取り出す必要がある。
そのための道具が\term{不変テンソル}{invariant tensor}である。
\begin{center}
\renewcommand{\arraystretch}{1.3}
\begin{tabular}{lll}
\toprule
群 & 不変テンソル & 作る一重項\\
\midrule
$\SU(N)$ & $\delta^\alpha{}_\beta$ & $\bar{\mathbf N}\otimes\mathbf N\ni\mathbf1$：$\bar\psi_\alpha\chi^\alpha$\\
$\SU(2)$ & $\epsilon^{ab}$ \ ($\epsilon^{12}=+1$) & $\mathbf2\otimes\mathbf2\ni\mathbf1$：$\epsilon^{ab}\psi_a\chi_b$\\
$\SU(3)$ & $\epsilon^{\alpha\beta\gamma}$ & $\mathbf3\otimes\mathbf3\otimes\mathbf3\ni\mathbf1$：$\epsilon^{\alpha\beta\gamma}q_\alpha q_\beta q_\gamma$\\
$\U(1)$ & — & 電荷の和がゼロ\\
\bottomrule
\end{tabular}
\end{center}
\begin{note}
$\epsilon^{ab}$ が $\SU(2)$ 不変であることは
$\epsilon^{ab}U_a{}^cU_b{}^d=(\det U)\,\epsilon^{cd}=\epsilon^{cd}$ による。
$\SU(3)$ の $\epsilon^{\alpha\beta\gamma}$ も同様に $\det U=1$ から不変である。
$\U(N)$ では $\det U=e^{i\phi}\ne1$ なので $\epsilon$ は不変テンソルにならない。
$\epsilon^{ab}$ の存在が $\mathbf2\simeq\bar{\mathbf2}$（擬実性）の別の表現であり、
$\tilde H^a=\epsilon^{ab}H^*_b$ である（$i\sigma^2=\begin{pmatrix}0&1\\-1&0\end{pmatrix}=\epsilon^{ab}$）。
\end{note}
$\SU(3)$ の $\epsilon^{\alpha\beta\gamma}$ はバリオンの波動関数
（3つのクォークのカラーの反対称な組み合わせ）を与える。
後で見るように、標準模型では $\epsilon^{\alpha\beta\gamma}$ を用いた
繰り込み可能な項は存在しない。これがバリオン数保存の理由である。
%------------------------------------------------
\section{左巻きダブレットと右巻きシングレット}
\label{sec:doublet-singlet}
%------------------------------------------------
\subsection{表現の割り当て}
標準模型の本質的な仮定は、
\begin{equation}
\boxed{
\text{左巻き場は }\SU(2)_L\text{ 二重項、右巻き場は }\SU(2)_L\text{ 一重項}
}
\end{equation}
という割り当てである。
1世代あたり
\begin{equation}
Q_L=\begin{pmatrix}u_L\\ d_L\end{pmatrix},\quad
L_L=\begin{pmatrix}\nu_L\\ e_L\end{pmatrix}
\qquad\text{と}\qquad
u_R,\quad d_R,\quad e_R
\end{equation}
を持つ。二重項の上下の成分は電荷が1だけ異なる粒子の組であり、
$W^\pm$ の放出・吸収で互いに移り変わる。
\subsection{実験的根拠}
この割り当ては次のような実験事実から決まった。
\paragraph{荷電カレントは左巻きのみ}
$\beta$崩壊やミューオン崩壊の分析から、
荷電カレントは V$-$A 形 $\bar\psi\gamma^\mu(1-\gamma^5)\psi'$ であり、
右巻きのカレント（V$+$A）は観測されていない。
ミューオン崩壊の Michel パラメータは V$-$A の予言と $10^{-3}$ 程度の精度で一致する。
これは右巻き場が $\SU(2)_L$ に関して一重項であることを意味する。
\paragraph{ニュートリノのヘリシティ}
Goldhaber らの実験（1958）により、ニュートリノのヘリシティは $-1/2$ であることが示された。
最小標準模型では $\nu_R$ を含めない。
\paragraph{右巻き場も中性カレントを持つ}
右巻き場は $\SU(2)_L$ 一重項であるが、超電荷を持つので
光子と $Z$ ボソンに結合する。第2章の命題~\ref{prop:neutral-current} より、
$Z$ との結合は $T_3-Q\sin^2\theta_W$ に比例し、
\begin{equation}
g_L^e=-\frac12+\sin^2\theta_W\simeq-0.27,
\qquad
g_R^e=\sin^2\theta_W\simeq0.23
\end{equation}
となる。
$e_R$ の $Z$ 結合がゼロでないことは、LEP での $Z$ の崩壊の精密測定で確かめられている（第12章）。
\begin{note}
$e_L$ は $T_3=-\frac12$, $Q=-1$ なので $g_L^e=-\frac12-(-1)\sin^2\theta_W$。
$e_R$ は $T_3=0$, $Q=-1$ なので $g_R^e=\sin^2\theta_W$。
$\sin^2\theta_W\simeq\frac14$ に近いため、
荷電レプトンの $Z$ へのベクトル結合 $g_V=g_L+g_R=-\frac12+2\sin^2\theta_W$ は非常に小さい。
\end{note}
\subsection{右巻きニュートリノ}
右巻きニュートリノ $\nu_R$ を加えるとすれば、その量子数は
\begin{equation}
\nu_R\sim(\mathbf1,\mathbf1)_0
\end{equation}
でなければならない（$T_3=0$, $Q=0$ より $Y=0$）。
これはすべてのゲージ相互作用に対して一重項であり（\term{ステライル}{sterile}）、
ゲージ相互作用だけでは生成も検出もできない。
$\nu_R$ の有無はニュートリノ質量の起源と関係しており、第11章で詳しく扱う。
\begin{remark}
右巻き場も二重項にまとめる\term{左右対称模型}{left-right symmetric model}
$\SU(2)_L\times\SU(2)_R\times\U(1)_{B-L}$ も考えられている。
この模型では $\nu_R$ が自然に現れ、パリティの破れは $\SU(2)_R$ の自発的破れとして説明される。
\end{remark}
\subsection{フレーバー対称性}
ゲージ相互作用は世代の添字に作用しないので、
フェルミオンの運動項 $\Lag_{\mathrm{fermion}}$ は世代の回転
\begin{equation}
Q_L\to U_QQ_L,\quad
u_R\to U_uu_R,\quad
d_R\to U_dd_R,\quad
L_L\to U_LL_L,\quad
e_R\to U_ee_R
\qquad(U_X\in\U(3))
\end{equation}
の下で不変である。この大域的対称性
\begin{equation}
G_F=\U(3)^5=\U(3)_Q\times\U(3)_u\times\U(3)_d\times\U(3)_L\times\U(3)_e
\end{equation}
を\term{フレーバー対称性}{flavor symmetry}と呼ぶ。
$G_F$ を破るのは Yukawa 結合だけである。
したがって、フレーバーを変える過程はすべて Yukawa 行列（質量と CKM 行列）を通じてのみ起こる。
これは第1章のパラメータの数え上げの基礎であり、
第9・10章の CKM 行列と GIM 機構の議論の出発点になる。
%------------------------------------------------
\section{クォークとレプトンのゲージ量子数}
\label{sec:quantum-numbers}
%------------------------------------------------
\subsection{量子数の一覧}
1世代分のゲージ量子数を、Dirac 表示と左巻き表示で並べると次のようになる
（第1章の表を再掲）。
\begin{center}
\renewcommand{\arraystretch}{1.3}
\begin{tabular}{cccc|cccc}
\toprule
\multicolumn{4}{c|}{Dirac 表示} & \multicolumn{4}{c}{左巻き表示}\\
場 & $\SU(3)_C$ & $\SU(2)_L$ & $Y$ & 場 & $\SU(3)_C$ & $\SU(2)_L$ & $Y$\\
\midrule
$Q_L$ & $\mathbf3$ & $\mathbf2$ & $+\frac13$ & $Q$ & $\mathbf3$ & $\mathbf2$ & $+\frac13$\\
$u_R$ & $\mathbf3$ & $\mathbf1$ & $+\frac43$ & $u^c$ & $\bar{\mathbf3}$ & $\mathbf1$ & $-\frac43$\\
$d_R$ & $\mathbf3$ & $\mathbf1$ & $-\frac23$ & $d^c$ & $\bar{\mathbf3}$ & $\mathbf1$ & $+\frac23$\\
$L_L$ & $\mathbf1$ & $\mathbf2$ & $-1$ & $L$ & $\mathbf1$ & $\mathbf2$ & $-1$\\
$e_R$ & $\mathbf1$ & $\mathbf1$ & $-2$ & $e^c$ & $\mathbf1$ & $\mathbf1$ & $+2$\\
\midrule
$H$ & $\mathbf1$ & $\mathbf2$ & $+1$ & $H$ & $\mathbf1$ & $\mathbf2$ & $+1$\\
\bottomrule
\end{tabular}
\end{center}
\subsection{超電荷を未知数として扱う}
第2章で見たように、$\U(1)_Y$ の超電荷は群論からは決まらない。
そこで、超電荷を未知数
\begin{equation}
Y_Q,\quad Y_u,\quad Y_d,\quad Y_L,\quad Y_e,\quad Y_H
\end{equation}
とおき、どのような条件で現実の値が決まるかを調べよう。
これは本章を通じたテーマである。

第1の条件は Yukawa 結合の存在である。
3種類の Yukawa 項 $\overline{Q_L}Hd_R$, $\overline{Q_L}\tilde Hu_R$, $\overline{L_L}He_R$ が
$\U(1)_Y$ 不変であるためには
\begin{equation}
Y_d=Y_Q-Y_H,\qquad
Y_u=Y_Q+Y_H,\qquad
Y_e=Y_L-Y_H
\label{eq:Y-yukawa}
\end{equation}
でなければならない。
第2の条件は、ヒッグス場の真空期待値が電荷を破らないことである。
真空期待値を持つ成分（下成分、$T_3=-\frac12$）が中性であるためには
\begin{equation}
Q=-\frac12+\frac{Y_H}{2}=0
\quad\Longrightarrow\quad
Y_H=1
\label{eq:YH}
\end{equation}
でなければならない（これは $\U(1)_Y$ の規格化を固定する条件とみなせる）。

これだけでは $Y_Q$ と $Y_L$ が未定のまま残る。
実際、\eqref{eq:Y-yukawa} と \eqref{eq:YH} の下で電荷は
\begin{equation}
Q(u)=\frac12+\frac{Y_Q}{2},\quad
Q(d)=-\frac12+\frac{Y_Q}{2},\quad
Q(\nu)=\frac12+\frac{Y_L}{2},\quad
Q(e)=-\frac12+\frac{Y_L}{2}
\end{equation}
となり、$Y_Q$, $Y_L$ は任意の実数でよい。
残りの2つの自由度を固定するのが、\S\ref{sec:anomaly} の異常相殺条件である。
%------------------------------------------------
\section{ゲージ不変な相互作用項}
\label{sec:invariant-terms}
%------------------------------------------------
\subsection{繰り込み可能な項の分類}
質量次元4以下で Lorentz 不変な演算子は、
場の次元 $[\psi]=\frac32$, $[\phi]=[A_\mu]=1$, $[\partial]=1$ から
次の形に限られる。
\begin{center}
\renewcommand{\arraystretch}{1.3}
\begin{tabular}{lll}
\toprule
型 & 例 & 標準模型での項\\
\midrule
$\bar\psi\gamma^\mu D_\mu\psi$ & フェルミオン運動項 & $\Lag_{\mathrm{fermion}}$\\
$F_{\mu\nu}F^{\mu\nu}$, $F\tilde F$ & ゲージ場運動項、$\theta$ 項 & $\Lag_{\mathrm{gauge}}$, $\Lag_\theta$\\
$|D_\mu\phi|^2$, $\phi^2$, $\phi^3$, $\phi^4$ & スカラーの運動項・ポテンシャル & $\Lag_{\mathrm{Higgs}}$\\
$\psi\psi$（$\bar\psi\psi$） & フェルミオン質量項 & なし（\S\ref{sec:no-bare-mass}）\\
$\psi\psi\phi$ & Yukawa 結合 & $\Lag_{\mathrm{Yukawa}}$\\
\bottomrule
\end{tabular}
\end{center}
運動項はゲージ不変性から共変微分の形に決まり、
第2章で見たように任意の表現に対して書ける。
非自明なのは、フェルミオンとスカラーを含む項（質量項・Yukawa 項）と
スカラーポテンシャルである。
\subsection{スカラーポテンシャル}
$H\sim(\mathbf1,\mathbf2)_1$ だけから作れるゲージ不変量は $H^\dagger H$ のみである。
\begin{itemize}
    \item $H^3$, $H^2H^\dagger$ などの奇数次の項は $Y$ の和がゼロにならない。
    \item $\epsilon^{ab}H_aH_b=0$（反対称テンソルと対称な積）。
\end{itemize}
したがって
\begin{equation}
V(H)=-\mu^2H^\dagger H+\lambda(H^\dagger H)^2
\end{equation}
が最も一般的なポテンシャルである（定数項を除く）。
ヒッグス二重項が2つ以上ある模型では、ポテンシャルはずっと複雑になる。
\subsection{Yukawa 項の網羅的な探索}
左巻き表示では、Yukawa 項は2つの左巻きフェルミオン $\psi_1\psi_2$ と
スカラー $H$ または $H^\dagger$ の積である。
1世代の左巻き場 $\{Q,u^c,d^c,L,e^c\}$ から2つを選び、
$H$（$Y=+1$）または $H^\dagger$（$Y=-1$）を掛けて一重項になるものを探す。
\begin{center}
\renewcommand{\arraystretch}{1.3}
\begin{tabular}{cccccc}
\toprule
$\psi_1\psi_2$ & $\SU(3)_C$ & $\SU(2)_L$ & $Y_1+Y_2$ & 相手 & 結果\\
\midrule
$Q\,u^c$ & $\mathbf3\otimes\bar{\mathbf3}\ni\mathbf1$ & $\mathbf2$ & $-1$ & $H$ & ○ 上型 Yukawa\\
$Q\,d^c$ & $\mathbf3\otimes\bar{\mathbf3}\ni\mathbf1$ & $\mathbf2$ & $+1$ & $H^\dagger$ & ○ 下型 Yukawa\\
$L\,e^c$ & $\mathbf1$ & $\mathbf2$ & $+1$ & $H^\dagger$ & ○ 荷電レプトン Yukawa\\
$Q\,L$ & $\mathbf3$ & & & & × カラー\\
$Q\,Q$ & $\mathbf3\otimes\mathbf3\not\ni\mathbf1$ & & & & × カラー\\
$u^c\,d^c$ & $\bar{\mathbf3}\otimes\bar{\mathbf3}\not\ni\mathbf1$ & & & & × カラー\\
$L\,L$ & $\mathbf1$ & $\mathbf1\oplus\mathbf3$ & $-2$ & & × $Y$\\
$e^c\,e^c$ & $\mathbf1$ & $\mathbf1$ & $+4$ & & × $Y$\\
$L\,e^c$ & $\mathbf1$ & $\mathbf2$ & $+1$ & $H$ & × $Y$\\
\bottomrule
\end{tabular}
\end{center}
その他の組み合わせも同様に、カラーまたは $Y$ の保存則で禁止される。
結局、許される Yukawa 項は
\begin{equation}
Q\,u^c\,H,\qquad Q\,d^c\,H^\dagger,\qquad L\,e^c\,H^\dagger
\end{equation}
の3種類（とそのエルミート共役）だけであり、これが第1章で与えた
\begin{equation}
\Lag_{\mathrm{Yukawa}}
=-y^d_{ij}\,\overline{Q_{Li}}Hd_{Rj}
-y^u_{ij}\,\overline{Q_{Li}}\tilde Hu_{Rj}
-y^e_{ij}\,\overline{L_{Li}}He_{Rj}+\hc
\label{eq:Yukawa-ch3}
\end{equation}
に対応する。
\subsection{ゲージ不変性の確認}
Dirac 表示で、$\SU(2)_L$ の添字を明示すると
\begin{equation}
\overline{Q_L}Hd_R=\overline{Q_L}{}_a\,H^a\,d_R,
\qquad
\overline{Q_L}\tilde Hu_R=\overline{Q_L}{}_a\,\epsilon^{ab}H^*_b\,u_R,
\qquad
\overline{L_L}He_R=\overline{L_L}{}_a\,H^a\,e_R
\end{equation}
である。
\begin{myprop}{Yukawa 項のゲージ不変性}{yukawa-invariance}
\eqref{eq:Yukawa-ch3} の各項は $\SU(3)_C\times\SU(2)_L\times\U(1)_Y$ 不変である。
\end{myprop}
\begin{note}
\textbf{$\SU(3)_C$}：$\overline{Q_L}$ と $d_R$ のカラー添字が $\delta$ で縮約され、
$\bar q\,U_3^\dagger U_3\,q$ の形になる。

\textbf{$\SU(2)_L$}：$\overline{Q_L}\to\overline{Q_L}U_2^\dagger$, $H\to U_2H$, $d_R\to d_R$ なので
$\overline{Q_L}Hd_R\to\overline{Q_L}U_2^\dagger U_2Hd_R=\overline{Q_L}Hd_R$。
$\tilde H$ も $U_2\tilde H$ と変換する（命題~\ref{prop:pseudoreal}）ので同様。

\textbf{$\U(1)_Y$}：$\bar\psi$ は $-Y$ を持つので、超電荷の和は
$\overline{Q_L}Hd_R$：$-\frac13+1-\frac23=0$、
$\overline{Q_L}\tilde Hu_R$：$-\frac13-1+\frac43=0$、
$\overline{L_L}He_R$：$+1+1-2=0$。
\end{note}
上型クォークに $\tilde H$ が必要なのは、
$\overline{Q_L}Hu_R$ の超電荷が $-\frac13+1+\frac43=2\ne0$ となるからである。
ヒッグス二重項が1つしかなくても上型・下型の両方に質量を与えられるのは、
$\SU(2)$ の擬実性（$\mathbf2\simeq\bar{\mathbf2}$）のおかげである。
超対称模型ではこの手が使えず、2つのヒッグス二重項が必要になる。
\subsection{偶然的対称性}
上の分類から、標準模型の繰り込み可能なラグランジアンには、
\begin{itemize}
    \item $\epsilon^{\alpha\beta\gamma}$（3つのクォーク）を含む項がない
    $\Rightarrow$ バリオン数 $B$ が保存する。
    \item $LL$, $Le^c$（同符号のレプトン数）を含む項がない
    $\Rightarrow$ レプトン数 $L$ が保存する。
\end{itemize}
これらは課した対称性ではなく、ゲージ対称性と繰り込み可能性から
自動的に従う\term{偶然的対称性}{accidental symmetry}である。
$B$ を破る最低次の演算子は次元6（$QQQL$ など）であり、
陽子崩壊の寿命の下限 $\tau_p\gtrsim10^{34}$ 年から
そのスケールは $\Lambda\gtrsim10^{15}\ \mathrm{GeV}$ と制限される。
\begin{remark}
最小超対称標準模型では、スカラー場の超対称パートナーを用いて
$u^cd^cd^c$ などの $B$ を破る繰り込み可能な項が書けてしまうため、
$R$ パリティと呼ばれる対称性を別途課す必要がある。
標準模型で $B$ と $L$ が自動的に保存することは、その意味で非自明な性質である。
\end{remark}
%------------------------------------------------
\section{なぜ裸の質量項を導入できないのか}
\label{sec:no-bare-mass}
%------------------------------------------------
\subsection{Dirac 質量項}
Dirac 質量項は
\begin{equation}
-m\bar\psi\psi=-m\left(\overline{\psi_L}\psi_R+\overline{\psi_R}\psi_L\right)
\end{equation}
であり、左巻き場と右巻き場を結ぶ。
これがゲージ不変であるためには、$\psi_L$ と $\psi_R$ が同じ表現に属さなければならない。
標準模型のすべての候補について確かめると次のようになる。
\begin{center}
\renewcommand{\arraystretch}{1.3}
\begin{tabular}{lccc}
\toprule
質量項 & $\SU(2)_L$ & $Y$ の和 & 判定\\
\midrule
$\overline{e_L}e_R$ & $\mathbf2$ の成分 $\otimes\ \mathbf1$ & $+1-2=-1$ & ×\\
$\overline{u_L}u_R$ & $\mathbf2$ の成分 $\otimes\ \mathbf1$ & $-\frac13+\frac43=1$ & ×\\
$\overline{d_L}d_R$ & $\mathbf2$ の成分 $\otimes\ \mathbf1$ & $-\frac13-\frac23=-1$ & ×\\
$\overline{\nu_L}\nu_R$（$\nu_R$ を加えた場合） & $\mathbf2$ の成分 $\otimes\ \mathbf1$ & $+1+0=1$ & ×\\
\bottomrule
\end{tabular}
\end{center}
すべての項で $Y$ の和が $\pm1$ となっている。
これは偶然ではなく、Yukawa 項 $\overline{\psi_L}H\psi_R$ が不変であることから
$\overline{\psi_L}\psi_R$ は $H$ の電荷 $\mp1$ を補えばよいことを意味する。
つまり\textbf{質量項に足りないのはちょうどヒッグス場1つ分の量子数}である。
$\SU(2)_L$ についても、$\overline{\psi_L}\psi_R$ は二重項の成分であり、
二重項 $H$ と縮約して初めて一重項になる。
\begin{note}
$\overline{e_L}e_R$ は $\SU(2)_L$ 二重項 $\overline{L_L}$ の下成分だけを取り出したものなので、
$\SU(2)_L$ 変換で $\overline{\nu_L}e_R$ と混ざる。
仮に二重項全体で $\overline{L_L}e_R$ と書いても、これは $\SU(2)_L$ 二重項（の共役）であり一重項ではない。
\end{note}
\subsection{Majorana 質量項}
1つの左巻き Weyl 場 $\chi$ だけで作る質量項
\begin{equation}
-\frac{m}{2}\left(\chi\chi+\hc\right)
\qquad\left(\text{Dirac 記法では}\ -\frac m2\overline{(\chi_L)^c}\chi_L+\hc\right)
\end{equation}
を \term{Majorana 質量項}{Majorana mass term} と呼ぶ。
これは $\chi$ の持つすべての $\U(1)$ 電荷を2倍だけ破るので、
$\chi$ が電荷を持つ場合は禁止される。
\begin{itemize}
    \item $e_L$, $u_L$ などの荷電フェルミオン：電荷 $Q\ne0$ を持つので禁止。
    \item $\nu_L$：電荷は持たないが、$LL$ は $Y=-2$ で $\SU(2)_L$ の $\mathbf1\oplus\mathbf3$ なので禁止。
    \item $\nu_R$（加えた場合）：$(\mathbf1,\mathbf1)_0$ なので\textbf{許される}。
    $\nu_R$ の Majorana 質量 $M_R$ はゲージ対称性で制限されず、任意に大きくてよい。
\end{itemize}
最後の点はシーソー機構の基礎となる（第11章）。
\begin{remark}
$LL$ の $\SU(2)_L$ 三重項部分（$Y=-2$）と結合するには、
$Y=+2$ の三重項スカラー $\Delta\sim(\mathbf1,\mathbf3)_2$ を導入すればよい
（type-II シーソー）。また、繰り込み不可能な Weinberg 演算子
$(LH)(LH)/\Lambda$ は $H$ を2つ用いて $Y$ と $\SU(2)_L$ を補っている（第11章）。
\end{remark}
\subsection{結論：すべての質量は電弱対称性の破れに由来する}
以上より、標準模型では
\begin{itemize}
    \item すべてのフェルミオン（$\nu_R$ を除く）の質量項はゲージ不変でない、
    \item ゲージボソンの質量項もゲージ不変でない（第2章）、
\end{itemize}
したがって、\textbf{電弱対称性が破れていなければ、標準模型の全粒子は massless である}
（ヒッグス場自身の $-\mu^2H^\dagger H$ を除く）。
観測される質量はすべて、電弱対称性の破れの大きさ $v$ に比例する：
\begin{equation}
m_f=\frac{y_fv}{\sqrt2},\qquad
M_W=\frac{gv}{2},\qquad
M_Z=\frac{\sqrt{g^2+g'^2}\,v}{2},\qquad
m_h=\sqrt{2\lambda}\,v .
\end{equation}
これらの導出は第6章・第7章で行う。
\begin{remark}
ヒッグス場がなくても、QCD のクォーク凝縮 $\langle\bar q_Lq_R\rangle\ne0$ は
$\overline{Q_L}$ と $q_R$ の積が二重項であることから $\SU(2)_L\times\U(1)_Y$ を破る。
このとき $W$, $Z$ はパイ中間子を「食べて」質量
$M_W\sim gf_\pi/2\sim30\ \mathrm{MeV}$ 程度を得る。
これは観測値よりはるかに小さいが、
「強い相互作用による動的な電弱対称性の破れ」というテクニカラー模型の着想の元になった。
\end{remark}
%------------------------------------------------
\section{ゲージ異常と異常相殺}
\label{sec:anomaly}
%------------------------------------------------
\subsection{量子異常}
古典的なラグランジアンが持つ対称性が、量子化の後に破れる現象を
\term{量子異常}{quantum anomaly}あるいは単に\term{アノマリー}{anomaly}と呼ぶ。
経路積分の言葉では、作用は不変でも
フェルミオンの経路積分測度 $\mathcal D\psi\mathcal D\bar\psi$ が不変でないことに由来する（藤川の方法）。
摂動論では、カイラルフェルミオンのループを含む
\term{三角形ダイアグラム}{triangle diagram}（カレント3つの相関関数）に現れる。

最も有名な例は、質量ゼロの Dirac フェルミオン（電荷 $e$）の軸性カレント
$j^\mu_5=\bar\psi\gamma^\mu\gamma^5\psi$ に対する
\term{Adler--Bell--Jackiw 異常}{ABJ anomaly}
\begin{equation}
\partial_\mu j^\mu_5=\frac{e^2}{8\pi^2}F_{\mu\nu}\tilde F^{\mu\nu},
\qquad
\tilde F^{\mu\nu}=\frac12\epsilon^{\mu\nu\rho\sigma}F_{\rho\sigma}
\end{equation}
である（全体の符号は $\epsilon^{0123}$ や $\gamma^5$ の規約に依存する）。
古典的には $\partial_\mu j^\mu_5=0$ であるが、量子論では右辺が現れる。
三角形ダイアグラムの計算そのものは場の量子論の標準的な教科書に譲り、本ノートでは結果のみを用いる。
\begin{myex}{$\pi^0\to\gamma\gamma$ とカラー数}{pi0}
$\pi^0$ の2光子崩壊は、軸性カレントの ABJ 異常によって支配される。
クォークの寄与を足し合わせると
\begin{equation}
\Gamma(\pi^0\to\gamma\gamma)
=\frac{\alpha^2m_\pi^3}{64\pi^3f_\pi^2}
\left[N_c\left(Q_u^2-Q_d^2\right)\right]^2
\end{equation}
となる（$f_\pi\simeq92\ \mathrm{MeV}$）。
$N_c=3$ では $N_c(Q_u^2-Q_d^2)=3\left(\frac49-\frac19\right)=1$ となり、
$\Gamma\simeq7.8\ \mathrm{eV}$ を与えて実験値とよく一致する。
$N_c=1$ では $1/9$ 倍になってしまう。
\end{myex}
\begin{note}
数値：$\alpha^2=5.33\times10^{-5}$、$m_\pi^3=(0.135)^3=2.46\times10^{-3}\ \mathrm{GeV}^3$、
$64\pi^3f_\pi^2=1984\times(0.0921)^2=16.8\ \mathrm{GeV}^2$。
よって $\Gamma=5.33\times10^{-5}\times2.46\times10^{-3}/16.8\simeq7.8\times10^{-9}\ \mathrm{GeV}$。
\end{note}
\subsection{大域的異常とゲージ異常}
ABJ 異常のように、\textbf{大域的}対称性が異常で破れることは問題ではない。
むしろ $\pi^0\to\gamma\gamma$ のような物理的な予言を与える。

一方、\textbf{ゲージ}対称性が異常で破れると理論は破綻する：
\begin{itemize}
    \item ゲージカレントが保存しないと Ward 恒等式が成り立たず、
    ゲージボソンの非物理的な偏極（縦波・時間成分）が物理的な状態から分離しない。
    \item その結果、ユニタリ性（確率の保存）や繰り込み可能性が失われる。
\end{itemize}
したがって、\textbf{ゲージ異常は相殺していなければならない}。
これはカイラルなゲージ理論の表現に対する非自明な制約である。
\subsection{異常係数}
左巻き Weyl フェルミオンの集合が、ゲージ群の生成子 $T^a$ の表現を持つとする
（右巻き場は荷電共役をとって左巻き場に直す）。
3つのゲージカレント $j^a,j^b,j^c$ の三角形ダイアグラムによる異常は、
\begin{equation}
\mathcal A^{abc}\equiv\sum_{\text{左巻き Weyl}}\tr\left[T^a\{T^b,T^c\}\right]
\end{equation}
に比例する。異常がないための条件は、すべての $a,b,c$ について
\begin{equation}
\boxed{\mathcal A^{abc}=0}
\end{equation}
である。
\begin{note}
右巻き場 $\psi_R\in R$ は、三角形ダイアグラムに左巻き場と逆符号で寄与する。
一方、その荷電共役 $(\psi_R)^c\in\bar R$ の生成子は $-T^{a*}=-T^{aT}$ であり、
\[
\tr\left[(-T^{aT})\{-T^{bT},-T^{cT}\}\right]
=-\tr\left[T^{aT}\{T^{bT},T^{cT}\}\right]
=-\tr\left[T^a\{T^b,T^c\}\right]
\]
（転置のトレースは不変で、積の順序が逆転するが反交換子なので同じ）。
したがって「右巻き場の寄与に $-1$ を掛ける」ことと
「左巻き表示で $\bar R$ として数える」ことは同値である。
\end{note}
\subsection{標準模型で調べるべき異常}
$G_{\SM}=\SU(3)_C\times\SU(2)_L\times\U(1)_Y$ の場合、
3つの生成子をどの因子から選ぶかで次の組み合わせがある。
\begin{center}
\renewcommand{\arraystretch}{1.3}
\begin{tabular}{lll}
\toprule
型 & 異常係数 & 自明か\\
\midrule
$\SU(3)^3$ & $\sum\tr_3[T^a\{T^b,T^c\}]$ & 非自明（$\mathbf3$ と $\bar{\mathbf3}$ の数の差）\\
$\SU(2)^3$ & $\tr[\sigma^a\{\sigma^b,\sigma^c\}]$ & 自明に $0$\\
$\SU(3)^2\,\U(1)$ & $\sum Y\,\tr_3[T^aT^b]$ & 非自明\\
$\SU(2)^2\,\U(1)$ & $\sum Y\,\tr_2[T^iT^j]$ & 非自明\\
$\U(1)^3$ & $\sum Y^3$ & 非自明\\
重力$^2\,\U(1)$ & $\sum Y$ & 非自明\\
$\SU(N)\times(\text{その他})^2$ など & $\propto\tr T^a$ & 自明に $0$\\
\bottomrule
\end{tabular}
\end{center}
\begin{note}
\textbf{$\SU(2)^3$}：$\{\sigma^b,\sigma^c\}=2\delta^{bc}$ なので
$\tr[\sigma^a\{\sigma^b,\sigma^c\}]=2\delta^{bc}\tr\sigma^a=0$。
より一般に、$\SU(2)$ の表現はすべて実または擬実なので $\SU(2)^3$ 異常は存在しない。

\textbf{生成子1つだけが非可換因子から来る場合}：
例えば $\SU(3)\,\U(1)^2$ では $\tr_3[T^a]\times(\cdots)=0$。
$\SU(3)\,\SU(2)\,\U(1)$ も $\tr T^a=0$ を含むので消える。

\textbf{重力}：ゲージカレント1つと重力子（エネルギー運動量テンソル）2つの三角形ダイアグラムも、
理論を重力と結合させたときに整合的であるためには相殺しなければならない。
その係数は $\sum Y$（$\U(1)$ 電荷の和）である。
\end{note}
\subsection{1世代での異常相殺}
左巻き表示の1世代 $\{Q,u^c,d^c,L,e^c\}$ について、非自明な条件を順に確かめる。
$T(\mathbf3)=T(\bar{\mathbf3})=T(\mathbf2)=\frac12$ を用いる。
\paragraph{$\SU(3)^3$}
$\mathbf3$ に属する Weyl 場は $Q$ の2成分（$\SU(2)$ の上下）、
$\bar{\mathbf3}$ に属するものは $u^c$, $d^c$ の2つ。
$\mathbf3$ と $\bar{\mathbf3}$ の異常係数は逆符号なので
\begin{equation}
\mathcal A_{333}\propto2-2=0 .
\end{equation}
これは $\SU(3)_C$ がベクトル的（第2章）であることの反映である。
\paragraph{$\SU(3)^2\,\U(1)_Y$}
カラーを持つ場の $Y$ を、$\SU(2)$ の多重度を掛けて足す：
\begin{equation}
\mathcal A_{33Y}\propto
\underbrace{2\cdot\tfrac13}_{Q}
+\underbrace{\left(-\tfrac43\right)}_{u^c}
+\underbrace{\tfrac23}_{d^c}
=0 .
\end{equation}
\paragraph{$\SU(2)^2\,\U(1)_Y$}
二重項の $Y$ を、カラーの多重度を掛けて足す：
\begin{equation}
\mathcal A_{22Y}\propto
\underbrace{3\cdot\tfrac13}_{Q}
+\underbrace{(-1)}_{L}
=0 .
\end{equation}
\paragraph{$\U(1)_Y^3$}
すべての Weyl 成分について
\begin{equation}
\mathcal A_{YYY}\propto
6\left(\tfrac13\right)^3+3\left(-\tfrac43\right)^3+3\left(\tfrac23\right)^3+2(-1)^3+(2)^3
=\tfrac29-\tfrac{64}9+\tfrac89-2+8=0 .
\end{equation}
\paragraph{重力$^2\,\U(1)_Y$}
\begin{equation}
\mathcal A_{\mathrm{grav}}\propto
6\cdot\tfrac13+3\left(-\tfrac43\right)+3\cdot\tfrac23+2(-1)+2
=2-4+2-2+2=0 .
\end{equation}
\begin{mythm}{標準模型の異常相殺}{anomaly-free}
標準模型のゲージ異常は、各世代ごとに独立に相殺する。
ただし相殺はクォークとレプトンの間で起こり、
クォークだけ、あるいはレプトンだけでは相殺しない。
\end{mythm}
例えば $\SU(2)^2\U(1)$ ではクォーク二重項の $+1$ とレプトン二重項の $-1$ が、
$\U(1)^3$ ではクォーク部分 $-6$ とレプトン部分 $+6$ が打ち消し合う。
もしトップクォークが存在しなければ第3世代の異常は相殺しないので、
異常相殺の観点からも、$b$ クォークと $\tau$ の発見の後にトップクォークの存在が強く期待された。
\subsection{異常相殺による超電荷の決定}
\S\ref{sec:quantum-numbers} で未定のまま残した超電荷を、異常相殺条件から決定しよう。
これは本章の主結果である。
\begin{mythm}{異常相殺と電荷の量子化}{charge-quantization}
1世代のフェルミオンの表現 $(\mathbf3,\mathbf2)_{Y_Q}$, $(\mathbf3,\mathbf1)_{Y_u}$, $(\mathbf3,\mathbf1)_{Y_d}$,
$(\mathbf1,\mathbf2)_{Y_L}$, $(\mathbf1,\mathbf1)_{Y_e}$（Dirac 表示）を仮定する。
4つの非自明な異常相殺条件の解は、$u\leftrightarrow d$ の入れ替えを除いて
\begin{equation}
(Y_Q,\ Y_u,\ Y_d,\ Y_L,\ Y_e)=Y_Q\times(1,\ 4,\ -2,\ -3,\ -6)
\end{equation}
または $Y_Q=0$ の自明に近い解（$Y_L=Y_e=0$, $Y_u=-Y_d$）である。
さらに Yukawa 結合の存在 \eqref{eq:Y-yukawa} と $Y_H=1$ \eqref{eq:YH} を課すと $Y_Q=\frac13$ となり、
すべての超電荷と電荷が一意に決まる：
\begin{equation}
Q(u)=\frac23,\quad Q(d)=-\frac13,\quad Q(\nu)=0,\quad Q(e)=-1 .
\end{equation}
\end{mythm}
\begin{note}
左巻き表示では $u^c$, $d^c$, $e^c$ の超電荷は $-Y_u$, $-Y_d$, $-Y_e$ である。

$\SU(2)^2\U(1)$：$3Y_Q+Y_L=0$ より $Y_L=-3Y_Q$。

$\SU(3)^2\U(1)$：$2Y_Q-Y_u-Y_d=0$。そこで $Y_u=Y_Q+x$, $Y_d=Y_Q-x$ とおく。

重力：$6Y_Q-3Y_u-3Y_d+2Y_L-Y_e=6Y_Q-6Y_Q-6Y_Q-Y_e=0$ より $Y_e=-6Y_Q$。

$\U(1)^3$：
\[
6Y_Q^3-3(Y_Q+x)^3-3(Y_Q-x)^3+2(-3Y_Q)^3-(-6Y_Q)^3
=6Y_Q^3-6Y_Q^3-18Y_Qx^2-54Y_Q^3+216Y_Q^3
=18Y_Q\left(9Y_Q^2-x^2\right)=0 .
\]
よって $Y_Q=0$ または $x=\pm3Y_Q$。
$x=3Y_Q$ のとき $Y_u=4Y_Q$, $Y_d=-2Y_Q$（$x=-3Y_Q$ は $u\leftrightarrow d$ の入れ替え）。

Yukawa 条件：$Y_H=Y_Q-Y_d=3Y_Q$, $Y_H=Y_u-Y_Q=3Y_Q$, $Y_H=Y_L-Y_e=3Y_Q$ はすべて整合しており、
$Y_H=1$ から $Y_Q=\frac13$。電荷は $Q(u)=\frac12+\frac16=\frac23$ などとなる。
\end{note}
この定理の意味は大きい。
$\U(1)$ の電荷は群論的には任意の実数でよいにもかかわらず、
\begin{itemize}
    \item 異常相殺（量子論の整合性）、
    \item Yukawa 結合の存在（フェルミオンが質量を持つこと）、
    \item ヒッグス場の真空が電荷を保存すること
\end{itemize}
を要求するだけで、クォークの電荷が $\frac23$, $-\frac13$、電子の電荷が $-1$ と決まる。
特に陽子と電子の電荷が厳密に逆符号で等しいこと
\begin{equation}
Q(p)+Q(e)=2\cdot\frac23-\frac13-1=0
\end{equation}
が説明される。
カラー数 $N_c=3$ も、$\SU(2)^2\U(1)$ の条件 $N_cY_Q+Y_L=0$ を通じて電荷の値と結びついている。
\subsection{Witten 異常}
$\SU(2)$ には、摂動論の三角形ダイアグラムでは見えない
\term{大域的異常}{global anomaly}（\term{Witten 異常}{Witten anomaly}）が存在する。
$\SU(2)$ 二重項の左巻き Weyl 場の数が奇数であると、
理論は数学的に定義できない（分配関数がゼロになる）。

標準模型の1世代には、$\SU(2)_L$ 二重項が
\begin{equation}
\underbrace{3}_{Q\text{ のカラー}}+\underbrace{1}_{L}=4
\end{equation}
個あり、偶数なので Witten 異常は存在しない。
これもクォークとレプトンが揃うことで満たされる条件である。
\subsection{大域的対称性の異常}
偶然的対称性 $\U(1)_B$, $\U(1)_L$ は、$\SU(2)_L$ のゲージ場との間に異常を持つ：
\begin{equation}
\partial_\mu j^\mu_B=\partial_\mu j^\mu_L
=n_g\frac{g^2}{32\pi^2}W^i_{\mu\nu}\tilde W^{i\mu\nu}+\cdots
\end{equation}
（$n_g=3$ は世代数）。
これらは大域的対称性なので理論の整合性は損なわれないが、
$B$ と $L$ が非摂動的に破れることを意味する。
一方 $B-L$ の異常は相殺するので、$B-L$ は厳密に保存する。
この $B+L$ を破る過程（\term{スファレロン}{sphaleron}過程）は
現在の宇宙では極めて稀であるが、初期宇宙の高温ではよく起こり、
バリオン数非対称性の生成に関与する（第14章）。
\begin{note}
$\SU(2)^2\U(1)_B$：二重項のバリオン数は $Q$ のみで、$3\times\frac13=1$（1世代あたり）。
$\SU(2)^2\U(1)_L$：二重項のレプトン数は $L$ のみで $1$。
両者は等しいので、$B-L$ では打ち消し合う。
\end{note}
\begin{remark}
右巻きニュートリノ $\nu^c\sim(\mathbf1,\mathbf1)_0$（$B-L=+1$）を各世代に加えると、
$\U(1)_{B-L}$ はすべての異常（$\U(1)_{B-L}^3$、重力$^2\U(1)_{B-L}$ を含む）が相殺し、
ゲージ化することが可能になる。例えば
\[
\sum(B-L)^3=6\left(\tfrac13\right)^3+3\left(-\tfrac13\right)^3+3\left(-\tfrac13\right)^3+2(-1)^3+1+1
=\tfrac29-\tfrac19-\tfrac19-2+2=0 .
\]
$\nu^c$ がないと $\sum(B-L)^3=-1\ne0$ となる。
これは $\SO(10)$ 大統一理論や左右対称模型で $\nu_R$ が自然に現れることと対応している。
\end{remark}
%------------------------------------------------
\section*{第3章のまとめ}
%------------------------------------------------
\begin{itemize}
    \item フェルミオン場は Lorentz・ゲージ・世代の添字を持つ。
    荷電共役によりすべてを左巻き Weyl 場で表すことができ、
    このときゲージ量子数は複素共役になる。
    \item 左巻き場は $\SU(2)_L$ 二重項、右巻き場は一重項であり、
    これは V$-$A 構造などの実験事実に基づく。
    ゲージ部分は大きな大域的フレーバー対称性 $\U(3)^5$ を持ち、それを破るのは Yukawa 結合のみである。
    \item ゲージ不変かつ繰り込み可能なフェルミオン--スカラー項は3種類の Yukawa 項のみであり、
    $B$ と $L$ は偶然的対称性として保存する。
    \item 裸の Dirac 質量項・Majorana 質量項はすべて禁止され（$\nu_R$ の Majorana 質量を除く）、
    質量項に足りない量子数はちょうどヒッグス場1つ分である。
    \item ゲージ異常は各世代でクォークとレプトンの間で相殺する。
    異常相殺と Yukawa 結合の存在から、超電荷（すなわち電荷）は一意に決まる。
\end{itemize}
%------------------------------------------------
\section*{第I部のまとめ}
%------------------------------------------------
第I部では、標準模型の対称性と場の内容を導入した。
\begin{enumerate}
    \item 標準模型はゲージ群 $\SU(3)_C\times\SU(2)_L\times\U(1)_Y$ に基づく
    繰り込み可能なカイラルゲージ理論である。
    \item ゲージ対称性は、ゲージ場の自己相互作用、フェルミオンとゲージ場の相互作用、
    および許される Yukawa 項の形を決める。
    \item ゲージ対称性はゲージボソンとフェルミオンの質量項を禁止する。
    \item 表現の割り当ては異常相殺条件を満たし、電荷の量子化を説明する。
\end{enumerate}
ゲージ対称性が厳密なら、すべての粒子は massless のはずである。
第II部では、自発的対称性の破れとヒッグス機構を導入し、
この対称性が「隠される」ことで粒子が質量を獲得する仕組みを導出する。
%================================================
% Part II
%================================================
\part{自発的対称性の破れとヒッグス機構}
%================================================
\chapter{自発的対称性の破れ（SSB）}
\label{ch:SSB}
%================================================
第I部の最後で見たように、標準模型のゲージ対称性が「そのまま」実現されていれば、
ゲージボソンもフェルミオンもすべて massless でなければならない。
この矛盾を解決する鍵が\term{自発的対称性の破れ}{spontaneous symmetry breaking, SSB}である。
本章ではまずゲージ場を含まない\textbf{大域的}対称性の場合にSSBを調べ、
\begin{itemize}
    \item 対称性が破れるとはどういうことか（\S\ref{sec:SSB-global}）、
    \item スカラー場のポテンシャルと真空の構造（\S\ref{sec:scalar-potential}, \S\ref{sec:vev}）、
    \item 連続対称性の破れに必ず伴う massless 粒子（\S\ref{sec:goldstone-thm}, \S\ref{sec:NG-boson}）、
    \item 具体例としての複素スカラー場の模型（\S\ref{sec:complex-SSB}）と自由度の数え上げ（\S\ref{sec:SSB-counting}）
\end{itemize}
を順に議論する。
局所（ゲージ）対称性の場合、ここで現れる massless 粒子はゲージ場に「食べられて」消え、
代わりにゲージ場が質量を得る。これが第5章のヒッグス機構である。
%------------------------------------------------
\section{大域的対称性と自発的対称性の破れ}
\label{sec:SSB-global}
%------------------------------------------------
\subsection{大域的対称性と Noether の定理}
場 $\phi_i$ の変換 $\phi_i\to\phi_i+\delta\phi_i$ が時空の点によらないパラメータ $\alpha^a$ で
$\delta\phi_i=\alpha^a\Delta^a_i(\phi)$ と書かれ、ラグランジアンを（全微分を除いて）不変に保つとき、
これを\term{大域的対称性}{global symmetry}という。
\begin{myprop}{Noether の定理}{noether}
連続的な大域的対称性 $\delta\phi_i=\alpha^a\Delta^a_i$ のもとで $\delta\Lag=\alpha^a\partial_\mu K^{a\mu}$ ならば、
\begin{equation}
j^{a\mu}=\frac{\partial\Lag}{\partial(\partial_\mu\phi_i)}\Delta^a_i-K^{a\mu}
\label{eq:noether-current}
\end{equation}
は運動方程式を用いると $\partial_\mu j^{a\mu}=0$ を満たす（\term{Noether カレント}{Noether current}）。
対応する\term{保存荷}{conserved charge}
\begin{equation}
Q^a=\int d^3x\,j^{a0}(t,\bm x)
\end{equation}
は時間によらず、量子論では場の変換を生成する：$[\,i\alpha^aQ^a,\phi_i\,]=\delta\phi_i$
（符号は $\alpha^a$ の向きの規約による）。
\end{myprop}
\begin{note}
パラメータを $\alpha^a\to\alpha^a(x)$ と局所化すると、
$\delta\Lag=\alpha^a\partial_\mu K^{a\mu}+(\partial_\mu\alpha^a)\frac{\partial\Lag}{\partial(\partial_\mu\phi_i)}\Delta^a_i$。
作用の変化は部分積分して $\delta S=-\int d^4x\,\alpha^a\,\partial_\mu\Bigl(\frac{\partial\Lag}{\partial(\partial_\mu\phi_i)}\Delta^a_i-K^{a\mu}\Bigr)$。
運動方程式を満たす配位では任意の $\alpha^a(x)$ に対して $\delta S=0$ なので、$\partial_\mu j^{a\mu}=0$ が従う。
\end{note}
対称性が量子論でどのように「実現」されるかには、本質的に異なる2つの様式がある。
\begin{mydef}{対称性の実現様式}{realization}
対称性の生成子 $Q^a$ に対して、
\begin{enumerate}
    \item 真空が不変 $Q^a|0\rangle=0$（すべての $a$）のとき、対称性は
    \term{Wigner--Weyl 様式}{Wigner--Weyl mode}で実現されているという。
    このとき粒子は対称性群の既約表現に従って縮退した多重項をなす。
    \item ラグランジアン（ハミルトニアン）は不変だが、真空が不変でない $Q^a|0\rangle\ne0$ とき、
    対称性は\term{自発的に破れている}{spontaneously broken}、あるいは
    \term{Nambu--Goldstone 様式}{Nambu--Goldstone mode}で実現されているという。
\end{enumerate}
\end{mydef}
SSB では、対称性は「失われた」のではなく、真空を選ぶことで\textbf{隠されて}いる\en{hidden symmetry}。
実際、ラグランジアンの対称性は、粒子のスペクトルや結合定数の間の関係として痕跡を残す
（\S\ref{sec:vev} の3点・4点結合の関係や、Goldstone の定理）。
これと区別して、ラグランジアン自身に対称性を破る項がある場合を
\term{陽な対称性の破れ}{explicit symmetry breaking}という。
\begin{note}
無限体積では $Q^a|0\rangle$ は規格化可能な状態ではない。
並進不変性から $\langle0|Q^aQ^a|0\rangle=\int d^3x\,\langle0|j^{a0}(\bm x)Q^a|0\rangle=V\langle0|j^{a0}(0)Q^a|0\rangle$ となり、
右辺が 0 でなければ体積 $V\to\infty$ で発散する。
したがって SSB の正確な定義は「$Q^a|0\rangle\ne0$」ではなく、
ある演算子 $\mathcal O$ について $\langle0|[Q^a,\mathcal O]|0\rangle\ne0$（\term{秩序変数}{order parameter}が 0 でない）
とするのがよい。交換子は局所性により有限である。
\end{note}
\subsection{SSB の例}
\begin{itemize}
    \item \textbf{強磁性体}\en{ferromagnet}：
    Heisenberg 模型 $H=-J\sum_{\langle ij\rangle}\bm S_i\cdot\bm S_j$（$J>0$）は空間回転 $\SO(3)$ で不変だが、
    基底状態ではスピンがそろって磁化 $\bm M\ne0$ を持ち、$\SO(3)$ が $\SO(2)$ に破れる。
    磁化の向きは任意で、無限に縮退した基底状態の1つが選ばれる。
    \item \textbf{QCD のカイラル対称性}\en{chiral symmetry}：
    $u,d$ クォークの質量を無視すると QCD は $\SU(2)_L\times\SU(2)_R$ で不変だが、
    クォーク凝縮 $\langle\bar qq\rangle\ne0$ により $\SU(2)_V$ に破れる。
    その痕跡がパイ中間子の軽さである（例題\ref{ex:pion-PNG}）。
    \item \textbf{電弱対称性}：
    ヒッグス場の真空期待値により $\SU(2)_L\times\U(1)_Y\to\U(1)_{\mathrm{EM}}$ と破れる（第6章）。
    ただしこれは局所対称性であり、その意味は第5章で注意深く議論する。
\end{itemize}
\subsection{なぜ場の理論では対称性が破れうるのか}
量子力学（有限自由度）では、対称な2重井戸ポテンシャルの基底状態は2つの井戸に広がった
対称な状態であり、対称性は破れない。
2つの「古典的真空」の間を\term{トンネル効果}{tunneling}が混ぜ、
エネルギー準位を $\Delta E\propto e^{-S_{\mathrm{inst}}/\hbar}$ だけ分裂させるからである。
場の理論（無限自由度）ではこの事情が変わる。
空間全体の場の値を一斉にひっくり返すトンネル振幅は体積に比例する作用で抑制され、
$\langle+|H|-\rangle\propto e^{-c V}\to0$（$V\to\infty$）となる。
したがって異なる真空は混ざらず、そのうちの1つが実現される。
\begin{note}
簡単な見積もり：空間を $N$ 個の独立な小領域に分け、各領域での2つの真空状態の重なりを $|\langle+|-\rangle_x|=r<1$ とすると、
全体では $\langle+|-\rangle=\prod_x\langle+|-\rangle_x$、すなわち $|\langle+|-\rangle|=r^N=e^{-N\ln(1/r)}\to0$（$N\propto V\to\infty$）。
異なる真空は互いに直交し、さらに局所演算子で互いに移り合うこともない。
このように異なる真空の上に作られた状態空間は、互いに干渉しない\term{超選択セクター}{superselection sector}をなす。
\end{note}
\begin{remark}
時空2次元（空間1次元）では、連続的な大域的対称性は自発的に破れない
（\term{Coleman--Mermin--Wagner の定理}{Coleman--Mermin--Wagner theorem}）。
2次元の massless スカラーの伝播関数 $\int d^2k/k^2$ が赤外発散し、揺らぎが秩序を壊すためである。
4次元の標準模型ではこの制限はない。
\end{remark}
%------------------------------------------------
\section{スカラー場のポテンシャル}
\label{sec:scalar-potential}
%------------------------------------------------
SSB を最も簡単に実現するのは、スカラー場のポテンシャルが原点以外に最小値を持つ場合である。
ローレンツ不変性を保ったまま真空期待値を持てるのは\textbf{スカラー場}だけである
（ベクトル場やスピノル場が真空期待値を持つとローレンツ対称性が破れる）。
\subsection{実スカラー場と $\mathbb Z_2$ 対称性}
1成分の実スカラー場を考える。
\begin{equation}
\Lag=\frac12\partial_\mu\phi\,\partial^\mu\phi-V(\phi),
\qquad
V(\phi)=\frac12\mu^2\phi^2+\frac\lambda4\phi^4 .
\label{eq:Z2-potential}
\end{equation}
この理論は\term{離散対称性}{discrete symmetry} $\mathbb Z_2:\phi\to-\phi$ を持つ。
ポテンシャルの形は次の要請で決まる。
\begin{itemize}
    \item \textbf{繰り込み可能性}：4次元ではスカラー場の質量次元は $[\phi]=1$ なので、
    $[V]=4$ より $\phi^n$（$n\le4$）の項だけが許される。
    \item \textbf{$\mathbb Z_2$ 対称性}：奇数べき $\phi,\phi^3$ は禁止される。
    \item \textbf{安定性}：エネルギーが下に有界であるためには $\lambda>0$ が必要。
\end{itemize}
\subsection{古典的な真空}
ハミルトニアン密度は
\begin{equation}
\mathcal H=\frac12\dot\phi^2+\frac12(\nabla\phi)^2+V(\phi)
\end{equation}
であり、最初の2項は非負である。
したがって最低エネルギーの配位（古典的な\term{真空}{vacuum}）は、
$\dot\phi=0$、$\nabla\phi=0$ の一様な場で、かつ $V$ の最小値を与えるものである。
$V'(\phi)=\phi(\mu^2+\lambda\phi^2)=0$ より、$\mu^2$ の符号で2つの場合に分かれる。
\begin{enumerate}
    \item $\mu^2>0$：最小値は $\phi=0$ のみ。真空は一意で $\mathbb Z_2$ 不変。
    $\phi$ は質量 $\mu$ の粒子を記述し、$\phi^4$ 項は自己相互作用である。
    \item $\mu^2<0$：$\phi=0$ は極大になり、最小値は
    \begin{equation}
    \phi=\pm v,\qquad v\equiv\sqrt{-\frac{\mu^2}{\lambda}}
    \label{eq:Z2-vev}
    \end{equation}
    の2点である（\term{2重井戸ポテンシャル}{double-well potential}）。
    2つの真空は $\mathbb Z_2$ で互いに移り合い、どちらを選んでも $\mathbb Z_2$ は破れる。
\end{enumerate}
\begin{note}
$\mu^2<0$ の場合 $\mu$ は「虚数の質量」であり、$\phi=0$ のまわりで展開すると
$(\Box+\mu^2)\phi=0$ の解は $\omega^2=\bm k^2+\mu^2$ となり、$|\bm k|<|\mu|$ のモードは指数関数的に増大する
\en{tachyonic instability}。$\phi=0$ は不安定な平衡点であり、正しい真空のまわりで展開しなければならない。
\end{note}
\subsection{量子補正と有効ポテンシャル}
上の議論は古典的なポテンシャルに基づく。
量子補正を含めると、真空は\term{有効ポテンシャル}{effective potential} $V_{\mathrm{eff}}(\phi_{\mathrm{cl}})$ の最小値で決まる。
1ループの補正（\term{Coleman--Weinberg ポテンシャル}{Coleman--Weinberg potential}）は
$\phi^4\ln(\phi^2/M^2)$ の形の項を与え、古典的には $\mu^2=0$ でも量子効果で対称性が破れうる。
標準模型のヒッグスポテンシャルの量子補正と真空の安定性は第14章で触れる。
本ノートでは以後、断らない限り古典的（ツリーレベル）なポテンシャルで議論する。
%------------------------------------------------
\section{真空期待値と真空の縮退}
\label{sec:vev}
%------------------------------------------------
\subsection{真空期待値}
量子論では、古典的な真空の値は場の\term{真空期待値}{vacuum expectation value, VEV}
\begin{equation}
\langle0|\phi(x)|0\rangle=v
\end{equation}
に対応する（ツリーレベル）。並進不変性から VEV は $x$ によらない定数である。
$\langle0|\phi|0\rangle\ne0$ は、$\mathbb Z_2$ のもとで奇な演算子が VEV を持つことを意味し、
$\mathbb Z_2$ が破れていることの秩序変数である。
\subsection{真空のまわりの展開}
摂動論は真空のまわりの小さな揺らぎに対して行う。そこで
\begin{equation}
\phi(x)=v+h(x)
\label{eq:Z2-shift}
\end{equation}
と置き、$h$ を新しい場とする（$\langle0|h|0\rangle=0$）。
ポテンシャルを $h$ で書き直すと
\begin{equation}
V=-\frac{\lambda v^4}{4}+\lambda v^2h^2+\lambda v\,h^3+\frac\lambda4h^4 .
\label{eq:Z2-shifted}
\end{equation}
\begin{note}
$\mu^2=-\lambda v^2$ を用いる。
\begin{align*}
\frac12\mu^2(v+h)^2&=-\frac{\lambda v^2}{2}\left(v^2+2vh+h^2\right),\\
\frac\lambda4(v+h)^4&=\frac\lambda4\left(v^4+4v^3h+6v^2h^2+4vh^3+h^4\right).
\end{align*}
和をとると、定数項 $(-\frac12+\frac14)\lambda v^4=-\frac14\lambda v^4$、
1次の項 $(-\lambda v^3+\lambda v^3)h=0$（最小値のまわりなので当然）、
2次の項 $(-\frac12+\frac32)\lambda v^2h^2=\lambda v^2h^2$、
3次の項 $\lambda vh^3$、4次の項 $\frac\lambda4h^4$。
\end{note}
2次の項 $\frac12m_h^2h^2$ と比べると、$h$ は質量
\begin{equation}
m_h^2=2\lambda v^2=-2\mu^2>0
\label{eq:Z2-mass}
\end{equation}
を持つ正常な（タキオンでない）粒子である。
\subsection{隠れた対称性の痕跡}
展開後のラグランジアン
\begin{equation}
\Lag=\frac12(\partial h)^2-\frac12m_h^2h^2-\lambda v\,h^3-\frac\lambda4h^4
\end{equation}
には $h\to-h$ の対称性はなく、$h^3$ 項が現れている。
しかし元の理論のパラメータは $(\mu^2,\lambda)$ の2つだけなので、
展開後の3つの量（質量・3点結合・4点結合）は独立ではない。
$(m_h,v)$ で表すと
\begin{equation}
\lambda=\frac{m_h^2}{2v^2},\qquad
\lambda v=\frac{m_h^2}{2v},\qquad
\frac\lambda4=\frac{m_h^2}{8v^2}
\label{eq:Z2-couplings}
\end{equation}
となる。この関係こそが「隠れた $\mathbb Z_2$ 対称性」の痕跡であり、
一般の $\mathbb Z_2$ を持たない理論（$h^3$ の係数が独立）とは区別できる。
\begin{note}
Feynman 則では、$-\frac{c_n}{n!}h^n$ の項に頂点因子 $-ic_n$ が対応する。
3点結合：$\lambda v h^3=\frac{6\lambda v}{3!}h^3$ より $-6i\lambda v=-3i\,m_h^2/v$。
4点結合：$\frac\lambda4h^4=\frac{6\lambda}{4!}h^4$ より $-6i\lambda=-3i\,m_h^2/v^2$。
標準模型のヒッグス粒子の自己結合も全く同じ形 $-3im_h^2/v$, $-3im_h^2/v^2$ になる（第7章）。
これを測ることはヒッグスポテンシャルの形を直接検証することに相当する。
\end{note}
\subsection{真空の縮退と領域壁}
2つの真空 $\phi=\pm v$ はエネルギーが等しく、\term{縮退}{degenerate}している。
どちらが選ばれるかは理論からは決まらず、初期条件や微小な外場で決まる。
初期宇宙で空間の異なる領域が異なる真空を選ぶと、その境界には
\term{領域壁}{domain wall}と呼ばれる位相的に安定な配位ができる。
\begin{myex}{キンク解（領域壁）}{kink}
1次元的な静的配位 $\phi=\phi(z)$ で $z\to\mp\infty$ で $\phi\to\mp v$ となるものを求めよ。
\tcblower
静的なので運動方程式は $\phi''=V'(\phi)$。両辺に $\phi'$ をかけて積分すると
$\frac12\phi'^2=V(\phi)-V(v)=\frac\lambda4(\phi^2-v^2)^2$。
したがって $\phi'=\sqrt{\lambda/2}\,(v^2-\phi^2)$ で、解は
\[
\phi(z)=v\tanh\left(\sqrt{\frac\lambda2}\,v\,z\right)=v\tanh\left(\frac{m_hz}{2}\right).
\]
壁の厚さは $\sim m_h^{-1}$、単位面積あたりのエネルギー（張力）は
$\sigma=\int dz\,\phi'^2=\frac{2\sqrt2}{3}\sqrt\lambda\,v^3$ である。
\end{myex}
\begin{note}
$V(\phi)-V(v)=\frac12\mu^2\phi^2+\frac\lambda4\phi^4+\frac{\lambda v^4}{4}=\frac\lambda4(\phi^4-2v^2\phi^2+v^4)=\frac\lambda4(\phi^2-v^2)^2$。
張力は $\sigma=\int_{-v}^{v}d\phi\,\phi'=\sqrt{\lambda/2}\int_{-v}^{v}(v^2-\phi^2)d\phi=\sqrt{\lambda/2}\cdot\frac43v^3$。
\end{note}
離散対称性の破れでは massless 粒子は現れない（$h$ は massive）。
これに対し\textbf{連続}対称性の破れでは、必ず massless 粒子が現れる。これが次節の主題である。
%------------------------------------------------
\section{Goldstone の定理}
\label{sec:goldstone-thm}
%------------------------------------------------
\begin{mythm}{Goldstone の定理}{goldstone}
ローレンツ不変で局所的な場の理論において、連続的な大域的対称性が自発的に破れると、
破れた生成子1つごとに1つの massless スカラー粒子
（\term{Nambu--Goldstone ボソン}{Nambu--Goldstone boson, NG boson}）が存在する。
対称性群 $G$ が部分群 $H$ に破れるとき、NG ボソンの数は
\begin{equation}
n_{\mathrm{NG}}=\dim G-\dim H
\label{eq:NG-number}
\end{equation}
である。
\end{mythm}
\subsection{古典論での証明}
$n$ 個の実スカラー場 $\phi_i$ が群 $G$ の（実）表現に従い、
無限小変換が
\begin{equation}
\delta\phi_i=\alpha^a\,t^a_{ij}\,\phi_j
\qquad(t^a:\ \text{実反対称行列},\ a=1,\dots,\dim G)
\end{equation}
で与えられるとする（エルミートな生成子とは $T^a=it^a$ で結ばれる）。
ポテンシャルが不変であることは、すべての $\phi$ に対して
\begin{equation}
\delta V=\frac{\partial V}{\partial\phi_i}\,\alpha^a t^a_{ij}\phi_j=0
\quad\Longrightarrow\quad
\frac{\partial V}{\partial\phi_i}\,t^a_{ij}\phi_j=0
\label{eq:V-invariance}
\end{equation}
を意味する。これを $\phi_k$ で微分すると
\begin{equation}
\frac{\partial^2V}{\partial\phi_k\partial\phi_i}\,t^a_{ij}\phi_j+\frac{\partial V}{\partial\phi_i}\,t^a_{ik}=0 .
\end{equation}
真空 $\phi=v$ で評価すると、第2項は $\partial V/\partial\phi_i|_v=0$ により消える。
\term{質量行列}{mass matrix} $(M^2)_{ki}\equiv\partial^2V/\partial\phi_k\partial\phi_i|_v$ を用いると
\begin{equation}
(M^2)_{ki}\,(t^av)_i=0 .
\label{eq:goldstone-classical}
\end{equation}
すなわち、ベクトル $t^av$ は質量行列の固有値 0 の固有ベクトルである。
\begin{itemize}
    \item 真空を変えない生成子（$t^av=0$）は自明に \eqref{eq:goldstone-classical} を満たし、何も与えない。
    これらの生成子は部分群 $H$（\term{非破れ部分群}{unbroken subgroup}、真空の\term{安定化部分群}{stabilizer}）をなす。
    \item 真空を変える生成子（$t^av\ne0$、\term{破れた生成子}{broken generator}）ごとに、
    質量 0 の方向 $t^av$ が1つ得られる。
\end{itemize}
\begin{note}
破れた生成子から得られるベクトルが線形独立であること：
$G$ のリー代数を $\mathfrak g=\mathfrak h\oplus\mathfrak m$ と分け、$\mathfrak h=\{t\in\mathfrak g\,|\,tv=0\}$ とする。
$\mathfrak m$ の基底 $t^{\hat a}$（$\hat a=1,\dots,\dim G-\dim H$）について $\sum_{\hat a}c_{\hat a}t^{\hat a}v=0$ ならば、
$\sum c_{\hat a}t^{\hat a}\in\mathfrak h$ となるが、$\mathfrak h\cap\mathfrak m=\{0\}$ なので $c_{\hat a}=0$。
したがって質量 0 の固有ベクトルは少なくとも $\dim G-\dim H$ 個ある。
「少なくとも」であるのは、ポテンシャルが偶然にもっと大きな対称性を持つ場合などに追加の massless 粒子がありうるためである
（\term{擬 NG ボソン}{pseudo-NG boson}、\S\ref{sec:NG-boson}）。
\end{note}
幾何学的には、$t^av$ は真空 $v$ を対称性で動かす方向である。
ポテンシャルは対称性で不変なので、この方向に沿ってポテンシャルは平坦であり、
平坦な方向の揺らぎは質量（ポテンシャルの曲率）を持たない。
\subsection{量子論での証明（概略）}
古典論の議論はツリーレベルに限られるが、Goldstone の定理は摂動論によらず成り立つ。
保存カレント $j^\mu$（$\partial_\mu j^\mu=0$）と、
\begin{equation}
\langle0|[Q,\phi(0)]|0\rangle\equiv\eta\ne0
\label{eq:SSB-condition}
\end{equation}
を満たす演算子 $\phi$ があるとする（SSB の条件）。
\begin{enumerate}
    \item $F(t)\equiv\int d^3x\,\langle0|[j^0(t,\bm x),\phi(0)]|0\rangle$ とおくと、
    \[
    \frac{dF}{dt}=\int d^3x\,\langle0|[\partial_0j^0,\phi(0)]|0\rangle=-\int d^3x\,\nabla\cdot\langle0|[\bm j(t,\bm x),\phi(0)]|0\rangle=0 .
    \]
    最後に、局所性（空間的に離れた演算子は可換）により十分遠方で交換子が 0 になり、表面項が消えることを用いた。
    したがって $F(t)=\eta$ は時間によらない。
    \item 完全系 $\sum_n|n\rangle\langle n|$ を挿入し、並進 $j^0(t,\bm x)=e^{iP\cdot x}j^0(0)e^{-iP\cdot x}$ を用いると
    \[
    F(t)=\sum_n(2\pi)^3\delta^3(\bm p_n)\Bigl[\langle0|j^0(0)|n\rangle\langle n|\phi(0)|0\rangle e^{-iE_nt}
    -\langle0|\phi(0)|n\rangle\langle n|j^0(0)|0\rangle e^{iE_nt}\Bigr].
    \]
    \item 真空自身の寄与は $\langle0|j^0|0\rangle=0$（ローレンツ不変性）で消える。
    $F$ が 0 でなく時間に依存しないためには、
    $\bm p_n=0$ で $E_n=0$ となる状態、すなわち $\bm p\to0$ で $E\to0$ となる励起が存在し、
    $\langle0|j^0|n\rangle\ne0$ かつ $\langle n|\phi|0\rangle\ne0$ でなければならない。
    ローレンツ不変性から、これは質量 0 の1粒子状態である。
\end{enumerate}
この粒子は $j^\mu$ によって真空から生成されるので、
\begin{equation}
\langle0|j^\mu(x)|\pi(p)\rangle=-iF\,p^\mu e^{-ip\cdot x}
\label{eq:NG-decay-constant}
\end{equation}
と書ける（$F$ は\term{崩壊定数}{decay constant}）。
$\partial_\mu j^\mu=0$ とすると $F\,p^2=0$ となり、$F\ne0$ なら $p^2=0$、すなわち massless であることが再確認できる。
NG ボソンは破れた生成子と同じ量子数を持つ。内部対称性の生成子はスカラーなので、NG ボソンはスピン 0 である
（ただし、生成子がパリティ奇なら擬スカラー）。
\begin{remark}
定理の仮定は (i) ローレンツ不変性、(ii) 局所性、(iii) 正定値のヒルベルト空間 である。
\begin{itemize}
    \item 非相対論的な系では (i) が成り立たず、NG ボソンの数や分散関係が変わりうる。
    例えば強磁性体のマグノンは $\omega\propto k^2$ で、破れた生成子2つ（$S_x,S_y$）に対して1種類しかない。
    \item ゲージ理論では、共変なゲージでは (iii) が、物理的なゲージ（クーロンゲージなど）では (i) と (ii) が成り立たず、定理を回避する。
    この結果、NG ボソンは物理的な粒子として現れず、ゲージ場が質量を得る。これがヒッグス機構である（第5章）。
\end{itemize}
\end{remark}
%------------------------------------------------
\section{Nambu--Goldstone ボソン}
\label{sec:NG-boson}
%------------------------------------------------
\subsection{NG ボソンの性質}
NG ボソンは次のような普遍的な性質を持つ。
\begin{enumerate}
    \item \textbf{シフト対称性}\en{shift symmetry}：
    NG ボソン場 $\pi$ は、破れた対称性のもとで $\pi\to\pi+v\alpha+\cdots$ と非斉次に変換する
    （\S\ref{sec:complex-SSB} で具体的に見る）。したがって $\pi$ の質量項 $\pi^2$ や $\pi^4$ のような項は禁止される。
    \item \textbf{微分結合}\en{derivative coupling}：
    適切な場の変数を選ぶと、NG ボソンは微分 $\partial_\mu\pi$ を通じてのみ他の場と結合する。
    \item \textbf{低エネルギー定理}\en{low-energy theorem}：
    その結果、NG ボソンが関わる散乱振幅は、NG ボソンの運動量 $p\to0$ で 0 になる
    （\term{Adler ゼロ}{Adler zero}）。NG ボソンは低エネルギーで弱く相互作用する。
    \item \textbf{崩壊定数}：NG ボソンの相互作用の強さは崩壊定数 $F$（ツリーレベルでは $F\sim v$）で決まり、
    典型的に $\partial\pi\partial\pi/F^2$ の形の結合を持つ。
\end{enumerate}
これらの性質は、NG ボソンを記述する低エネルギー有効理論
（\term{非線形シグマ模型}{nonlinear sigma model}、\term{カイラル摂動論}{chiral perturbation theory}）の基礎になる。
\subsection{擬 NG ボソン}
現実には対称性が厳密でなく、小さな陽な破れがあることが多い。
このとき NG ボソンは小さな質量を得て、\term{擬 NG ボソン}{pseudo-Nambu--Goldstone boson, PNGB}となる。
質量の2乗は陽な破れのパラメータに比例する（$m^2\propto\epsilon$、\S\ref{sec:complex-SSB} の例題）。
\begin{myex}{パイ中間子：QCD の擬 NG ボソン}{pion-PNG}
$u,d$ クォークの質量を 0 とすると、QCD ラグランジアン
$\bar q\,i\gamma^\mu D_\mu q=\bar q_Li\gamma^\mu D_\mu q_L+\bar q_Ri\gamma^\mu D_\mu q_R$（$q=(u,d)^T$）は、
左巻きと右巻きを独立に回す $\SU(2)_L\times\SU(2)_R$ で不変である。
強い相互作用の非摂動効果でクォーク凝縮 $\langle\bar qq\rangle=\langle\bar q_Lq_R+\bar q_Rq_L\rangle\ne0$ が生じると、
これは $q_L\to Lq_L$, $q_R\to Rq_R$ のもとで $L=R$ のときだけ不変なので、
\[
\SU(2)_L\times\SU(2)_R\to\SU(2)_V
\]
と破れる。NG ボソンの数は $3+3-3=3$ で、これが $\pi^\pm,\pi^0$ である。
破れた生成子は軸性（パリティ奇）なので、パイ中間子は擬スカラーである。
\tcblower
クォーク質量 $m_q\bar qq$ は陽な破れであり、パイ中間子は小さな質量を得る。
\term{Gell-Mann--Oakes--Renner の関係式}{Gell-Mann--Oakes--Renner relation}
\[
m_\pi^2f_\pi^2=(m_u+m_d)\,|\langle\bar qq\rangle|
\]
により $m_\pi^2\propto m_q$ である（$f_\pi\simeq92\ \mathrm{MeV}$）。
実際 $m_\pi\simeq140\ \mathrm{MeV}$ は、他のハドロン（$m_\rho\simeq775\ \mathrm{MeV}$, $m_p\simeq938\ \mathrm{MeV}$）より著しく軽い。
$s$ クォークまで含めると $\SU(3)_L\times\SU(3)_R\to\SU(3)_V$ で $8$ 個の擬 NG ボソン（$\pi,K,\eta$）が現れる。
なお $f_\pi$ は荷電パイ中間子の崩壊 $\pi^+\to\mu^+\nu_\mu$ の振幅を決める量であり（例題\ref{ex:pion}）、
\eqref{eq:NG-decay-constant} の $F$ にあたる。
\end{myex}
\begin{remark}
歴史的には、南部が超伝導の BCS 理論との類推から、核子の質量とパイ中間子の軽さを
カイラル対称性の自発的破れで説明した（1960--61年、南部--Jona-Lasinio 模型）。
Goldstone はスカラー場の模型でこれを一般化し（1961年）、
Goldstone--Salam--Weinberg が一般的な証明を与えた（1962年）。
\end{remark}
\subsection{標準模型における NG ボソン}
標準模型では、ヒッグス二重項の SSB により3個の NG ボソンが現れる（\S\ref{sec:SSB-counting}）。
しかし破れる対称性は\textbf{ゲージ}対称性なので、これらは物理的な粒子としては現れず、
$W^\pm$, $Z$ の縦偏極成分になる（第5章、第6章）。
それでも高エネルギーでは、縦偏極の $W,Z$ の散乱振幅は対応する NG ボソンの振幅で近似できる
（\term{等価定理}{Goldstone boson equivalence theorem}、第12章）。
%------------------------------------------------
\section{複素スカラー場の SSB}
\label{sec:complex-SSB}
%------------------------------------------------
連続対称性の SSB の最も簡単な例として、$\U(1)$ 対称性を持つ複素スカラー場を考える。
\begin{equation}
\Lag=\partial_\mu\phi^*\partial^\mu\phi-V(\phi),\qquad
V(\phi)=\mu^2|\phi|^2+\lambda|\phi|^4 .
\label{eq:U1-model}
\end{equation}
これは大域的 $\U(1)$ 変換 $\phi\to e^{i\alpha}\phi$ で不変である。
\subsection{ポテンシャルと真空}
$\lambda>0$, $\mu^2<0$ とする。ポテンシャルは $|\phi|$ のみの関数で、
\begin{equation}
V=\lambda\left(|\phi|^2-\frac{v^2}{2}\right)^2-\frac{\lambda v^4}{4},
\qquad v\equiv\sqrt{-\frac{\mu^2}{\lambda}}
\label{eq:mexican-hat}
\end{equation}
と書ける。最小値は円周
\begin{equation}
|\phi|=\frac{v}{\sqrt2}
\qquad\text{すなわち}\qquad
\phi=\frac{v}{\sqrt2}e^{i\theta}\quad(0\le\theta<2\pi)
\end{equation}
上にあり、真空は連続無限に縮退している（\term{メキシカンハット型ポテンシャル}{Mexican hat potential}、
\term{ワインボトル型ポテンシャル}{wine-bottle potential}）。
$\U(1)$ 変換は $\theta\to\theta+\alpha$ と真空を互いに移す。
一般性を失わずに $\theta=0$、すなわち
\begin{equation}
\langle0|\phi|0\rangle=\frac{v}{\sqrt2}\quad(\text{実数})
\end{equation}
を選ぶ。$1/\sqrt2$ は、実部・虚部を正準規格化された実スカラー場にするための因子である。
\begin{note}
$\lambda(|\phi|^2-\frac{v^2}{2})^2=\lambda|\phi|^4-\lambda v^2|\phi|^2+\frac{\lambda v^4}{4}$ であり、
$-\lambda v^2=\mu^2$ なので \eqref{eq:U1-model} と一致する。
\end{note}
\subsection{直交座標による展開}
真空のまわりで
\begin{equation}
\phi(x)=\frac{1}{\sqrt2}\bigl(v+h(x)+i\chi(x)\bigr)
\label{eq:U1-cartesian}
\end{equation}
と展開する（$h,\chi$ は実スカラー場）。運動項は $\partial\phi^*\partial\phi=\frac12(\partial h)^2+\frac12(\partial\chi)^2$、ポテンシャルは
\begin{equation}
V=\lambda v^2h^2+\lambda v\,h\,(h^2+\chi^2)+\frac\lambda4(h^2+\chi^2)^2+\text{const.}
\label{eq:U1-cartesian-V}
\end{equation}
となる。したがって
\begin{equation}
m_h^2=2\lambda v^2,\qquad m_\chi=0 .
\end{equation}
\begin{note}
$|\phi|^2=\frac12\left[(v+h)^2+\chi^2\right]$ なので
$|\phi|^2-\frac{v^2}{2}=\frac12\left(2vh+h^2+\chi^2\right)$。
これを \eqref{eq:mexican-hat} に代入して
$V+\frac{\lambda v^4}{4}=\frac\lambda4(2vh+h^2+\chi^2)^2=\lambda v^2h^2+\lambda vh(h^2+\chi^2)+\frac\lambda4(h^2+\chi^2)^2$。
\end{note}
$h$ はポテンシャルの「谷を横切る」動径方向の揺らぎで massive、
$\chi$ は「谷に沿う」角度方向の揺らぎで massless の NG ボソンである。
破れた生成子は $\U(1)$ の1つなので、Goldstone の定理と一致する。
実際、$\delta\phi=i\alpha\phi$ より $\delta(h+i\chi)=i\alpha(v+h+i\chi)$、すなわち
\begin{equation}
\delta h=-\alpha\chi,\qquad \delta\chi=\alpha(v+h)
\end{equation}
であり、真空の近傍では $\chi\to\chi+\alpha v$ というシフトになっている。
\subsection{極座標による展開}
同じ理論を極座標
\begin{equation}
\phi(x)=\frac{1}{\sqrt2}\bigl(v+h(x)\bigr)e^{i\pi(x)/v}
\label{eq:U1-polar}
\end{equation}
で書くこともできる。このとき $\U(1)$ 変換は\textbf{厳密に}
\begin{equation}
h\to h,\qquad \pi\to\pi+v\alpha
\end{equation}
というシフトになる。ラグランジアンは
\begin{equation}
\Lag=\frac12(\partial_\mu h)^2+\frac12\left(1+\frac hv\right)^2(\partial_\mu\pi)^2
-\lambda v^2h^2-\lambda vh^3-\frac\lambda4h^4 .
\label{eq:U1-polar-L}
\end{equation}
\begin{note}
$\partial_\mu\phi=\frac{1}{\sqrt2}e^{i\pi/v}\left[\partial_\mu h+i\frac{v+h}{v}\partial_\mu\pi\right]$ なので
$\partial_\mu\phi^*\partial^\mu\phi=\frac12\left[(\partial h)^2+\left(1+\frac hv\right)^2(\partial\pi)^2\right]$
（交差項 $\pm i\,\partial h\,\partial\pi$ は打ち消す）。
ポテンシャルは $|\phi|=\frac{v+h}{\sqrt2}$ のみに依存するので、\eqref{eq:U1-cartesian-V} で $\chi=0$ とおいたものになる。
\end{note}
極座標では NG ボソン $\pi$ がポテンシャルから完全に消え、
\begin{equation}
\Lag\supset\frac{h}{v}(\partial\pi)^2+\frac{h^2}{2v^2}(\partial\pi)^2
\end{equation}
のように\textbf{微分を通じてのみ}結合する。これは $\pi$ のシフト対称性の直接の帰結である。
一方、直交座標では $\lambda vh\chi^2$ のような微分を含まない結合がある。
2つの記述は場の再定義で結ばれており、S 行列（物理的な散乱振幅）は一致する。
直交座標で計算すると、$\lambda vh\chi^2$ 項と $h$ の交換図などが打ち消し合って、
最終的に低エネルギーで振幅が 0 になる（Adler ゼロ）。
\subsection{Noether カレントと崩壊定数}
$\U(1)$ の Noether カレントは、\eqref{eq:noether-current} から（全体の符号を選んで）
\begin{equation}
j^\mu=i\left(\phi\,\partial^\mu\phi^*-\phi^*\partial^\mu\phi\right)
=\frac{(v+h)^2}{v}\partial^\mu\pi
=v\,\partial^\mu\pi+2h\,\partial^\mu\pi+\frac{h^2}{v}\partial^\mu\pi .
\label{eq:U1-current}
\end{equation}
\begin{note}
$\phi^*\partial^\mu\phi=\frac12(v+h)\left[\partial^\mu h+i\frac{v+h}{v}\partial^\mu\pi\right]$ とその複素共役から
$\phi^*\partial^\mu\phi-\phi\,\partial^\mu\phi^*=i\frac{(v+h)^2}{v}\partial^\mu\pi$。
これに $-i$ をかけて \eqref{eq:U1-current} を得る。
\end{note}
カレントの最低次は $v\,\partial^\mu\pi$ という\textbf{1次の項}であり、
カレントが真空から NG ボソンを1つ生成する：
\begin{equation}
\langle0|j^\mu(x)|\pi(p)\rangle=v\,\partial^\mu e^{-ip\cdot x}=-iv\,p^\mu e^{-ip\cdot x}.
\end{equation}
\eqref{eq:NG-decay-constant} と比べて、崩壊定数は（ツリーレベルで）$F=v$ である。
SSB のない場合（$v=0$）、カレントは場の2次式で、真空から1粒子を作ることはない。
\begin{myex}{陽な破れと擬 NG ボソンの質量}{U1-explicit}
\eqref{eq:U1-model} に小さな陽な破れ $V_{\mathrm{br}}=-\epsilon(\phi+\phi^*)$（$\epsilon>0$）を加える。
NG ボソンの質量を $\epsilon$ の1次で求めよ。
\tcblower
極座標 \eqref{eq:U1-polar} で $V_{\mathrm{br}}=-\sqrt2\,\epsilon\,(v+h)\cos(\pi/v)$。
$\pi$ について展開すると $\cos(\pi/v)=1-\frac{\pi^2}{2v^2}+\cdots$ なので、$h=0$ で
\[
V_{\mathrm{br}}=-\sqrt2\,\epsilon v+\frac{\sqrt2\,\epsilon}{2v}\pi^2+\cdots
\quad\Longrightarrow\quad
m_\pi^2=\frac{\sqrt2\,\epsilon}{v}.
\]
質量の2乗は陽な破れのパラメータ $\epsilon$ に比例する（GMOR の関係式 $m_\pi^2\propto m_q$ と同じ構造）。
また、陽な破れは縮退していた真空のうち $\theta=0$（$\pi=0$）を選び出す。
$\epsilon$ による $v$ 自身のずれは $O(\epsilon)$ であり、$m_\pi^2$ には $O(\epsilon^2)$ でしか効かない。
\end{myex}
\subsection{$\mathrm O(N)$ 模型}
$\U(1)\simeq\SO(2)$ の模型は、$N$ 成分の実スカラー場 $\bm\phi=(\phi_1,\dots,\phi_N)$ の
\term{線形シグマ模型}{linear sigma model}
\begin{equation}
\Lag=\frac12\partial_\mu\bm\phi\cdot\partial^\mu\bm\phi-\frac12\mu^2\bm\phi^2-\frac\lambda4(\bm\phi^2)^2
\label{eq:ON-model}
\end{equation}
に一般化される（$N=2$ が $\U(1)$ 模型で、$\phi=(\phi_1+i\phi_2)/\sqrt2$）。
これは $\mathrm O(N)$ 回転 $\bm\phi\to R\bm\phi$ で不変である。
$\mu^2<0$ のとき真空は $|\bm\phi|=v=\sqrt{-\mu^2/\lambda}$ の球面 $S^{N-1}$ 上にあり、
$\langle\bm\phi\rangle=(0,\dots,0,v)$ を選ぶと、
\begin{equation}
\bm\phi=(\pi_1,\dots,\pi_{N-1},\,v+\sigma)
\end{equation}
と展開して
\begin{equation}
V=\lambda v^2\sigma^2+\lambda v\,\sigma(\sigma^2+\bm\pi^2)+\frac\lambda4(\sigma^2+\bm\pi^2)^2+\text{const.}
\end{equation}
となる。$\sigma$ は質量 $m_\sigma^2=2\lambda v^2$ を持ち、$N-1$ 個の $\pi_i$ は massless である。
\begin{note}
$\bm\phi^2-v^2=2v\sigma+\sigma^2+\bm\pi^2$ であり、$V=\frac\lambda4(\bm\phi^2-v^2)^2+\text{const.}$ を展開すればよい。
真空 $(0,\dots,0,v)$ を変えない回転は、最初の $N-1$ 成分の回転 $\mathrm O(N-1)$ である。
破れた生成子は $(i,N)$ 平面の回転 $(i=1,\dots,N-1)$ の $N-1$ 個で、
\[
\dim\mathrm O(N)-\dim\mathrm O(N-1)=\frac{N(N-1)}{2}-\frac{(N-1)(N-2)}{2}=N-1
\]
は NG ボソンの数と一致する。
\end{note}
$N=4$ の場合が特に重要である。標準模型のヒッグス二重項は4つの実スカラー場からなり、
そのポテンシャルは $\mathrm O(4)$ 模型と同じ形をしている（\S\ref{sec:SSB-counting}、第6章）。
QCD の $\SU(2)_L\times\SU(2)_R$ も $\SO(4)$ と局所的に同型であり、
$N=4$ の線形シグマ模型（$\sigma$ と3つの $\pi$）はパイ中間子の古典的な模型でもある。
%------------------------------------------------
\section{対称性の破れと自由度の数え上げ}
\label{sec:SSB-counting}
%------------------------------------------------
\subsection{真空多様体}
SSB で縮退した真空全体の集合 $\mathcal M_{\mathrm{vac}}=\{g\cdot v\,|\,g\in G\}$ を
\term{真空多様体}{vacuum manifold}という。
$H$ の元は $v$ を動かさないので、$\mathcal M_{\mathrm{vac}}$ は\term{剰余空間}{coset space}
\begin{equation}
\mathcal M_{\mathrm{vac}}\simeq G/H
\end{equation}
と同一視できる。その次元 $\dim G-\dim H$ は、真空を動かす独立な方向の数、
すなわち NG ボソンの数に等しい。NG ボソンの場は、真空多様体上の座標（の揺らぎ）とみなせる。
\begin{note}
別の真空 $v'=gv$ を選ぶと、非破れ部分群は $H'=gHg^{-1}$（共役な部分群）になる。
どの真空を選んでも物理（スペクトルや相互作用）は同じであり、
破れのパターンは $G\to H$ の共役類で特徴づけられる。
\end{note}
\subsection{自由度の保存}
SSB は場の自由度の数を変えない。$n$ 個の実スカラー場は、破れた後も $n$ 個の実スカラー場であり、
\begin{equation}
n=\underbrace{(\dim G-\dim H)}_{\text{NG ボソン（massless）}}+\underbrace{(n-\dim G+\dim H)}_{\text{massive なスカラー}}
\end{equation}
と分かれるだけである。
\begin{table}[htbp]
\centering
\caption{代表的な対称性の破れのパターンと NG ボソンの数}
\label{tab:SSB-patterns}
\begin{tabular}{llccc}
\toprule
模型 & $G\to H$ & 実スカラー & $\dim G-\dim H$ & $\mathcal M_{\mathrm{vac}}$\\
\midrule
実スカラー & $\mathbb Z_2\to\{1\}$ & 1 & 0（離散） & 2点\\
複素スカラー & $\U(1)\to\{1\}$ & 2 & $1-0=1$ & $S^1$\\
$\mathrm O(N)$ 模型 & $\mathrm O(N)\to\mathrm O(N-1)$ & $N$ & $N-1$ & $S^{N-1}$\\
$\SU(3)$ 三重項 & $\SU(3)\to\SU(2)$ & 6 & $8-3=5$ & $S^5$\\
ヒッグス二重項 & $\SU(2)\times\U(1)\to\U(1)$ & 4 & $4-1=3$ & $S^3$\\
QCD（2フレーバー） & $\SU(2)_L\times\SU(2)_R\to\SU(2)_V$ & -- & $6-3=3$ & $\SU(2)\simeq S^3$\\
QCD（3フレーバー） & $\SU(3)_L\times\SU(3)_R\to\SU(3)_V$ & -- & $16-8=8$ & $\SU(3)$\\
\bottomrule
\end{tabular}
\end{table}
\begin{note}
$\SU(3)$ 三重項：$\phi\in\mathbb C^3$ は6個の実自由度。$\langle\phi\rangle=(0,0,v/\sqrt2)$ は最初の2成分に作用する $\SU(2)$ で不変
（$\U(1)$ をすべて考えると $\U(3)\to\U(2)$ で $9-4=5$ と同じ）。真空多様体は $|\phi|=v/\sqrt2$ の $S^5$。
6 = 5（NG）+ 1（動径）。
QCD の場合、NG ボソンはクォークの複合状態でありスカラー場の数は意味を持たないので「--」とした。
\end{note}
\subsection{ヒッグス二重項の場合（第6章の予告）}
標準模型のヒッグス二重項
\begin{equation}
\Phi=\begin{pmatrix}\phi^+\\ \phi^0\end{pmatrix}
=\frac{1}{\sqrt2}\begin{pmatrix}\phi_1+i\phi_2\\ \phi_3+i\phi_4\end{pmatrix}
\sim(\mathbf1,\mathbf2)_{1}
\end{equation}
を考える。$\Phi^\dagger\Phi=\frac12(\phi_1^2+\phi_2^2+\phi_3^2+\phi_4^2)$ なので、
ポテンシャル $V=\mu^2\Phi^\dagger\Phi+\lambda(\Phi^\dagger\Phi)^2$ は $N=4$ の $\mathrm O(4)$ 模型と同じ形である。
真空を
\begin{equation}
\langle\Phi\rangle=\frac{1}{\sqrt2}\begin{pmatrix}0\\ v\end{pmatrix}
\end{equation}
に選ぶ。$\SU(2)_L\times\U(1)_Y$ の生成子 $T^i=\sigma^i/2$ と $Y/2=1/2$ を真空に作用させると
\begin{equation}
T^1\langle\Phi\rangle=\frac{v}{2\sqrt2}\begin{pmatrix}1\\0\end{pmatrix},\quad
T^2\langle\Phi\rangle=\frac{v}{2\sqrt2}\begin{pmatrix}-i\\0\end{pmatrix},\quad
T^3\langle\Phi\rangle=\frac{v}{2\sqrt2}\begin{pmatrix}0\\-1\end{pmatrix},\quad
\frac Y2\langle\Phi\rangle=\frac{v}{2\sqrt2}\begin{pmatrix}0\\1\end{pmatrix}.
\end{equation}
したがって
\begin{equation}
Q\langle\Phi\rangle=\left(T^3+\frac Y2\right)\langle\Phi\rangle=0
\end{equation}
であり、電荷 $Q$ に対応する $\U(1)_{\mathrm{EM}}$ が破れずに残る。
残りの3つ（$T^1$, $T^2$, $T^3-Y/2$ など）が破れた生成子で、NG ボソンは
\begin{equation}
n_{\mathrm{NG}}=\dim(\SU(2)\times\U(1))-\dim\U(1)=4-1=3
\end{equation}
個である。4つの実スカラーのうち、3つが NG ボソン、1つが massive なヒッグス粒子 $h$ になる。
\begin{note}
実際の変換は $\delta\Phi=i\alpha^aT^a\Phi$ なので、$i$ をかけた方向を見る。
$\Phi=\frac{1}{\sqrt2}(\phi_1+i\phi_2,\ \phi_3+i\phi_4)^T$ と比べると
\[
iT^1\langle\Phi\rangle=\frac{v}{2\sqrt2}\begin{pmatrix}i\\0\end{pmatrix}\to\phi_2,\quad
iT^2\langle\Phi\rangle=\frac{v}{2\sqrt2}\begin{pmatrix}1\\0\end{pmatrix}\to\phi_1,\quad
i\Bigl(T^3-\frac Y2\Bigr)\langle\Phi\rangle=\frac{v}{\sqrt2}\begin{pmatrix}0\\-i\end{pmatrix}\to\phi_4 .
\]
したがって NG 方向は $\phi_1,\phi_2,\phi_4$ で、残る $\phi_3$（$\langle\Phi\rangle$ の方向）が動径方向の揺らぎ、すなわちヒッグス粒子 $h$ である。
\end{note}
\begin{remark}
ヒッグスポテンシャルは実は $\SU(2)_L\times\U(1)_Y$ より大きな $\mathrm O(4)\simeq\SU(2)_L\times\SU(2)_R/\mathbb Z_2$ の対称性を持ち、
真空により $\SU(2)_L\times\SU(2)_R\to\SU(2)_V$ と破れる。
残る $\SU(2)_V$ は\term{カストディアル対称性}{custodial symmetry}と呼ばれ、
$W$ と $Z$ の質量比 $\rho=m_W^2/(m_Z^2\cos^2\theta_W)=1$ を保証する（第6章）。
NG ボソンの数は $\mathrm O(4)\to\mathrm O(3)$ でも $6-3=3$ で変わらない。
\end{remark}
\subsection{ゲージ対称性の場合への見通し}
電弱対称性はゲージ対称性なので、上の3つの NG ボソンは物理的な粒子としては現れない。
第5章で示すように、ゲージ対称性の破れでは
\begin{itemize}
    \item NG ボソンはゲージ変換で取り除くことができ（ユニタリーゲージ）、
    \item 代わりに対応するゲージボソンが質量を得て、縦偏極の自由度を獲得する。
\end{itemize}
自由度で数えると、破れる前は
\[
\underbrace{4}_{\text{実スカラー}}+\underbrace{4\times2}_{\text{massless ゲージ場 }W^{1,2,3},B}=12,
\]
破れた後は
\[
\underbrace{1}_{h}+\underbrace{3\times3}_{\text{massive }W^\pm,Z}+\underbrace{2}_{\text{massless }\gamma}=12
\]
となり、全体の数は保存する。NG ボソンの3自由度が、$W^\pm$, $Z$ の縦偏極に移ったのである。
%------------------------------------------------
\section*{第4章のまとめ}
%------------------------------------------------
\begin{itemize}
    \item ラグランジアンが対称でも真空が対称でないとき、対称性は自発的に破れる（SSB）。
    無限自由度の場の理論では異なる真空が混ざらないため、SSB が可能になる。
    \item スカラー場のポテンシャル $V=\mu^2|\phi|^2+\lambda|\phi|^4$ は、$\mu^2<0$ のとき原点以外に最小値を持ち、
    場は真空期待値 $v=\sqrt{-\mu^2/\lambda}$ を持つ。真空のまわりの揺らぎ $h$ は質量 $m_h^2=2\lambda v^2$ を持つ。
    \item 隠れた対称性は、質量と自己結合の間の関係（$h^3$ 結合 $\propto m_h^2/v$ など）として痕跡を残す。
    \item 連続対称性が $G\to H$ と破れると、$\dim G-\dim H$ 個の massless な NG ボソンが現れる（Goldstone の定理）。
    NG ボソンはシフト対称性を持ち、微分結合のみで相互作用する。小さな陽な破れがあると擬 NG ボソンとなる（例：パイ中間子）。
    \item 標準模型のヒッグス二重項では $\SU(2)_L\times\U(1)_Y\to\U(1)_{\mathrm{EM}}$ と破れ、3個の NG ボソンが現れる。
    これらはゲージ対称性のため物理的粒子にはならず、$W^\pm$, $Z$ の縦偏極成分となる（第5章、第6章）。
\end{itemize}
%================================================
\chapter{ゲージ対称性の破れとヒッグス機構}
\label{ch:Higgs-mechanism}
%================================================
第4章では大域的対称性の SSB を調べ、連続対称性の破れには必ず massless の NG ボソンが伴うことを見た。
しかし標準模型で破れるのは\textbf{ゲージ}対称性であり、自然界には対応する massless スカラーは見つかっていない。
本章では、ゲージ対称性のもとでは NG ボソンがゲージ場に吸収され、
代わりにゲージ場が質量を得ること（\term{ヒッグス機構}{Higgs mechanism}）を示す。
\begin{itemize}
    \item 局所対称性の「破れ」の意味（\S\ref{sec:local-breaking}）、
    \item 最も簡単な模型である Abelian Higgs model と、そのゲージ場の質量生成（\S\ref{sec:abelian-higgs}, \S\ref{sec:gauge-mass}）、
    \item NG ボソンの行方とゲージの選び方（\S\ref{sec:NG-gauge-dof}--\S\ref{sec:Rxi}）、
    \item 自由度の数え上げ（\S\ref{sec:Higgs-counting}）と非可換ゲージ群への一般化（\S\ref{sec:nonabelian-higgs}）
\end{itemize}
を順に議論する。本章の結果は、第6章で標準模型の $\SU(2)_L\times\U(1)_Y$ にそのまま適用される。
%------------------------------------------------
\section{局所対称性の破れ}
\label{sec:local-breaking}
%------------------------------------------------
\subsection{ゲージ対称性は「対称性」か}
大域的対称性は、物理的に異なる状態を互いに移す本当の対称性である。
これに対しゲージ対称性（局所対称性）は、同じ物理状態を表す場の配位の間の
\term{冗長性}{redundancy}であり、物理的な状態・観測量はすべてゲージ不変である。
この違いから、「ゲージ対称性の自発的破れ」という言葉には注意が必要である。
\begin{mythm}{Elitzur の定理}{elitzur}
局所的なゲージ対称性は自発的に破れない。
すなわち、ゲージ不変でない局所演算子の期待値は、ゲージを固定しなければ 0 である。
\end{mythm}
\begin{note}
直観的な理由：大域的対称性では、真空を一斉に回すには無限体積にわたる変換が必要で、
そのエネルギー障壁（トンネル振幅の抑制）が無限大になる（\S\ref{sec:SSB-global}）。
ゲージ変換は一点ごとに独立に行えるので、1点の場の向きを変えてもエネルギーは変わらない。
したがって $\langle\phi(x)\rangle$ をゲージ変換で平均すると 0 になる。
\end{note}
したがって、ゲージ不変な形で $\langle\phi\rangle\ne0$ を言うことはできない。
実際に行っていることは次の通りである。
\begin{enumerate}
    \item ゲージを固定する（例えばユニタリーゲージや $R_\xi$ ゲージ）。
    ゲージ固定後は局所対称性は残っていないが、\textbf{大域的な部分}（定数パラメータの変換）は残りうる。
    \item その大域的な部分が、通常の SSB として破れる。$\langle\phi\rangle=v/\sqrt2$ はこの意味の VEV である。
\end{enumerate}
「ゲージ対称性の破れ」は、この手続きを簡潔に表す便利な言葉と理解すればよい。
物理的な帰結（ゲージボソンの質量、スペクトル）はゲージ不変であり、言葉の選び方によらない。
\begin{remark}
基本表現のヒッグス場を持つ理論では、ヒッグス相と閉じ込め相の間に相転移がなく、
両者は連続的につながっている（\term{Fradkin--Shenker の定理}{Fradkin--Shenker theorem}）。
例えば標準模型の $W$ ボソンは、ゲージ不変な複合演算子 $\Phi^\dagger D_\mu\Phi$ などの束縛状態と見ることもできる
（Fr\"ohlich--Morchio--Strocchi）。弱結合の摂動論では、ゲージ固定した VEV の描像で正しい結果が得られる。
\end{remark}
\subsection{歴史と超伝導}
ゲージ場が質量を得る機構は、まず超伝導で見出された。
超伝導体の内部では磁場が侵入長 $\lambda_L$ の距離で指数関数的に減衰する
（\term{Meissner 効果}{Meissner effect}）。これは、超伝導体中の光子が質量 $m_\gamma=1/\lambda_L$ を持つことと同じである。
Anderson（1963年）は、ゲージ場があると NG ボソンが現れずに光子が質量を得ることを指摘し、
相対論的な場の理論での機構は 1964 年に Englert--Brout、Higgs、Guralnik--Hagen--Kibble によって独立に示された。
このため\term{Brout--Englert--Higgs 機構}{BEH mechanism}とも呼ばれる。
%------------------------------------------------
\section{Abelian Higgs model}
\label{sec:abelian-higgs}
%------------------------------------------------
第4章の $\U(1)$ 模型 \eqref{eq:U1-model} の対称性を局所化する。
\begin{equation}
\Lag=-\frac14F_{\mu\nu}F^{\mu\nu}+(D_\mu\phi)^*(D^\mu\phi)-V(\phi),
\qquad
V(\phi)=\mu^2|\phi|^2+\lambda|\phi|^4,
\label{eq:abelian-higgs}
\end{equation}
\begin{equation}
D_\mu\phi=(\partial_\mu-igA_\mu)\phi,\qquad F_{\mu\nu}=\partial_\mu A_\nu-\partial_\nu A_\mu .
\end{equation}
これを\term{Abelian Higgs model}{Abelian Higgs model}（スカラー QED にポテンシャルを加えたもの）という。
ゲージ変換は
\begin{equation}
\phi(x)\to e^{i\alpha(x)}\phi(x),\qquad A_\mu(x)\to A_\mu(x)+\frac1g\partial_\mu\alpha(x)
\label{eq:U1-gauge-transf}
\end{equation}
である（第1章の規約）。
\begin{note}
$D_\mu\phi\to\partial_\mu(e^{i\alpha}\phi)-ig\bigl(A_\mu+\tfrac1g\partial_\mu\alpha\bigr)e^{i\alpha}\phi
=e^{i\alpha}\bigl(\partial_\mu\phi+i(\partial_\mu\alpha)\phi-igA_\mu\phi-i(\partial_\mu\alpha)\phi\bigr)=e^{i\alpha}D_\mu\phi$。
したがって $|D_\mu\phi|^2$ と $|\phi|^2$ は不変。$F_{\mu\nu}$ は $\partial_\mu\partial_\nu\alpha-\partial_\nu\partial_\mu\alpha=0$ より不変。
\end{note}
運動方程式は
\begin{equation}
\partial_\mu F^{\mu\nu}=-g\,j^\nu,\qquad
j^\nu=i\bigl[\phi^*D^\nu\phi-\phi\,(D^\nu\phi)^*\bigr]
\label{eq:abelian-higgs-eom}
\end{equation}
であり（カレントの全体の符号は規約による）、$j^\nu$ は共変微分を含むためゲージ場自身に依存する。
SSB が起こると、このカレントが $A_\mu$ に比例する項を持ち、それが質量項として働く（次節）。
\begin{note}
$\partial\Lag/\partial A_\nu$ を計算する：$(D_\mu\phi)^*D^\mu\phi$ のうち $A$ を含む部分は
$ig A_\mu\phi^*\partial^\mu\phi-igA_\mu\phi\,\partial^\mu\phi^*+g^2A^2|\phi|^2$。
これを $A_\nu$ で微分すると $ig(\phi^*\partial^\nu\phi-\phi\partial^\nu\phi^*)+2g^2A^\nu|\phi|^2
=ig\bigl[\phi^*D^\nu\phi-\phi(D^\nu\phi)^*\bigr]$。
$-\frac14F^2$ からは $\partial_\mu F^{\mu\nu}$ が出て、Euler--Lagrange 方程式 $\partial_\mu F^{\mu\nu}+ig[\cdots]=0$、すなわち \eqref{eq:abelian-higgs-eom} を得る。
\end{note}
%------------------------------------------------
\section{ゲージ場の質量生成}
\label{sec:gauge-mass}
%------------------------------------------------
\subsection{直交座標による展開}
$\mu^2<0$ のとき、第4章と同様に
\begin{equation}
\phi=\frac{1}{\sqrt2}(v+h+i\chi),\qquad v=\sqrt{-\mu^2/\lambda}
\end{equation}
と展開する。共変微分は
\begin{equation}
D_\mu\phi=\frac{1}{\sqrt2}\Bigl[(\partial_\mu h+gA_\mu\chi)+i(\partial_\mu\chi-gvA_\mu-gA_\mu h)\Bigr]
\end{equation}
なので、
\begin{align}
|D_\mu\phi|^2
&=\frac12(\partial_\mu h+gA_\mu\chi)^2+\frac12(\partial_\mu\chi-gvA_\mu-gA_\mu h)^2 \notag\\
&=\frac12(\partial_\mu h)^2+\frac12(\partial_\mu\chi)^2
-gv\,A^\mu\partial_\mu\chi+\frac12g^2v^2A_\mu A^\mu \notag\\
&\quad+gA^\mu(\chi\partial_\mu h-h\partial_\mu\chi)+g^2v\,hA_\mu A^\mu+\frac12g^2(h^2+\chi^2)A_\mu A^\mu .
\label{eq:abelian-higgs-cartesian}
\end{align}
最も重要なのは
\begin{equation}
\frac12g^2v^2A_\mu A^\mu\equiv\frac12m_A^2A_\mu A^\mu,\qquad m_A=gv
\label{eq:mA}
\end{equation}
という\textbf{ゲージ場の質量項}が現れたことである。
ゲージ場の質量は、結合定数と VEV の積で決まる。
ただし同時に、$A_\mu$ と $\chi$ が混ざる2次の項 $-gvA^\mu\partial_\mu\chi$ も現れている。
この項の意味は \S\ref{sec:NG-gauge-dof} で調べる。
\begin{note}
$D_\mu\phi=\frac{1}{\sqrt2}\bigl[\partial_\mu h+i\partial_\mu\chi-igA_\mu(v+h+i\chi)\bigr]$ で、
$-igA_\mu\cdot i\chi=gA_\mu\chi$ が実部、$-igA_\mu(v+h)$ が虚部に入る。
2乗の展開：$\frac12(\partial h+gA\chi)^2=\frac12(\partial h)^2+gA^\mu\chi\partial_\mu h+\frac12g^2A^2\chi^2$、
$\frac12(\partial\chi-gA(v+h))^2=\frac12(\partial\chi)^2-gA^\mu(v+h)\partial_\mu\chi+\frac12g^2(v+h)^2A^2$。
\end{note}
\subsection{極座標と Proca 場}
NG ボソンの扱いは極座標
\begin{equation}
\phi=\frac{1}{\sqrt2}(v+h)\,e^{i\pi/v}
\end{equation}
の方が見やすい。このとき
\begin{equation}
D_\mu\phi=\frac{1}{\sqrt2}e^{i\pi/v}\Bigl[\partial_\mu h-ig(v+h)\Bigl(A_\mu-\frac{1}{gv}\partial_\mu\pi\Bigr)\Bigr]
\end{equation}
となるので、新しいベクトル場
\begin{equation}
B_\mu\equiv A_\mu-\frac{1}{gv}\partial_\mu\pi
\label{eq:Bmu}
\end{equation}
を定義すると、$\pi$ は $B_\mu$ の中に完全に吸収される。
$B_\mu$ は $A_\mu$ とゲージ変換の形で違うだけなので、場の強さは $F_{\mu\nu}(B)=F_{\mu\nu}(A)$ である。結果は
\begin{equation}
\Lag=-\frac14F_{\mu\nu}F^{\mu\nu}+\frac12m_A^2B_\mu B^\mu
+\frac12(\partial_\mu h)^2-V(h)
+\frac{m_A^2}{v}hB_\mu B^\mu+\frac{m_A^2}{2v^2}h^2B_\mu B^\mu
\label{eq:abelian-higgs-unitary}
\end{equation}
となる（$V(h)=\lambda v^2h^2+\lambda vh^3+\frac\lambda4h^4$）。
$\pi$ はラグランジアンから完全に消えた。
\begin{note}
$|D_\mu\phi|^2=\frac12(\partial h)^2+\frac12g^2(v+h)^2B_\mu B^\mu$（交差項は $\pm i$ で打ち消す）。
$\frac12g^2(v+h)^2=\frac12m_A^2\left(1+\frac hv\right)^2=\frac12m_A^2+\frac{m_A^2}{v}h+\frac{m_A^2}{2v^2}h^2$。
\end{note}
\eqref{eq:abelian-higgs-unitary} の $B_\mu$ の部分
\begin{equation}
\Lag_{\mathrm{Proca}}=-\frac14F_{\mu\nu}F^{\mu\nu}+\frac12m^2B_\mu B^\mu
\end{equation}
は massive なベクトル場の\term{Proca ラグランジアン}{Proca Lagrangian}である。
\begin{myprop}{Proca 場の性質}{proca}
Proca 場の運動方程式 $\partial_\mu F^{\mu\nu}+m^2B^\nu=0$ から
\begin{equation}
\partial_\nu B^\nu=0,\qquad(\Box+m^2)B^\nu=0
\end{equation}
が従う。したがって $B_\mu$ は質量 $m$ の粒子を記述し、
4成分のうち $\partial\cdot B=0$ で1つが除かれて、物理的な自由度は $3$（スピン1の $2s+1=3$ 状態）である。
偏極ベクトル $\epsilon^\mu_\lambda(k)$（$k\cdot\epsilon=0$, $\epsilon^2=-1$）は、$k^\mu=(E,0,0,k)$ に対して
\begin{equation}
\epsilon^\mu_\pm=\frac{1}{\sqrt2}(0,1,\pm i,0),\qquad
\epsilon^\mu_L=\frac1m(k,0,0,E)
\label{eq:polarization}
\end{equation}
であり、偏極和と伝播関数は
\begin{equation}
\sum_\lambda\epsilon^\mu_\lambda\epsilon^{\nu*}_\lambda=-g^{\mu\nu}+\frac{k^\mu k^\nu}{m^2},
\qquad
\tilde D^{\mu\nu}(k)=\frac{-i}{k^2-m^2+i\epsilon}\left(g^{\mu\nu}-\frac{k^\mu k^\nu}{m^2}\right)
\label{eq:proca-propagator}
\end{equation}
で与えられる。
\end{myprop}
\begin{note}
運動方程式に $\partial_\nu$ を作用させると、$\partial_\nu\partial_\mu F^{\mu\nu}=0$（反対称×対称）より $m^2\partial_\nu B^\nu=0$。
$m\ne0$ なので $\partial\cdot B=0$。これを使うと $\partial_\mu F^{\mu\nu}=\Box B^\nu-\partial^\nu(\partial\cdot B)=\Box B^\nu$。
縦偏極：$\epsilon_L\cdot k=(kE-Ek)/m=0$、$\epsilon_L^2=(k^2-E^2)/m^2=-1$。
\end{note}
縦偏極ベクトルは高エネルギー $E\gg m$ で
\begin{equation}
\epsilon^\mu_L=\frac{k^\mu}{m}+O\!\left(\frac mE\right)
\label{eq:epsL-highE}
\end{equation}
と振る舞い、エネルギーとともに大きくなる。
このため縦偏極ベクトルボソンの散乱振幅は、個々のダイアグラムでは $E^2/m^2$ のように増大しうる。
その増大がヒッグス粒子の寄与で相殺されることは、標準模型の重要な整合性条件である（第12章）。
また、$k^\mu/m$ の部分は「NG ボソンの名残り」であり、
高エネルギーでは縦偏極の振幅が NG ボソンの振幅に一致する（等価定理）。
\begin{note}
$\epsilon_L-\frac km=\frac1m(k-E,0,0,E-k)$ で、$E-k=\frac{m^2}{E+k}\simeq\frac{m^2}{2E}$ なので差は $O(m/E)$。
\end{note}
\begin{myex}{超伝導と London 方程式}{london}
超伝導体の中では、電子の Cooper 対（電荷 $2e$）が凝縮して秩序変数 $\langle\phi\rangle\ne0$ を持つ。
Abelian Higgs model の非相対論版（Ginzburg--Landau 理論）で、静的な磁場の振る舞いを求めよ。
\tcblower
静的で $A_0=0$ のとき、$B_\mu$ の運動方程式 $(\Box+m^2)\bm A=0$ は $\nabla^2\bm A=m^2\bm A$ となる。
回転をとると磁場 $\bm B=\nabla\times\bm A$ も $\nabla^2\bm B=m^2\bm B$ を満たし、
超伝導体の表面（$z=0$）から内部へ
\[
B(z)=B(0)\,e^{-z/\lambda_L},\qquad \lambda_L=\frac1m
\]
と減衰する（Meissner 効果）。$\lambda_L$ は London の侵入長で、典型的には $10$--$100\ \mathrm{nm}$、
すなわち「光子の質量」は $m_\gamma\sim2$--$20\ \mathrm{eV}$ である。
カレント $\bm j\propto-m^2\bm A$ は London 方程式そのものである。
\end{myex}
%------------------------------------------------
\section{Goldstone ボソンとゲージ自由度}
\label{sec:NG-gauge-dof}
%------------------------------------------------
\subsection{NG ボソンはゲージ自由度である}
極座標で、ゲージ変換 \eqref{eq:U1-gauge-transf} は
\begin{equation}
h\to h,\qquad \pi\to\pi+v\,\alpha(x),\qquad A_\mu\to A_\mu+\frac1g\partial_\mu\alpha
\end{equation}
と作用する。大域的な場合（第4章）には定数シフトだった NG ボソンの変換が、局所的なシフトになっている。
したがって $\alpha(x)=-\pi(x)/v$ と選べば、各点で $\pi$ を 0 にできる。
NG ボソンは物理的な自由度ではなく、\textbf{ゲージの自由度}そのものである。
\eqref{eq:Bmu} の $B_\mu$ はこのゲージ変換を施した後の $A_\mu$ であり、ゲージ不変な組み合わせになっている：
\begin{equation}
B_\mu=A_\mu-\frac{1}{gv}\partial_\mu\pi\ \to\ A_\mu+\frac1g\partial_\mu\alpha-\frac{1}{gv}(\partial_\mu\pi+v\partial_\mu\alpha)=B_\mu .
\end{equation}
ヒッグス機構のよく使われる表現では、ゲージ場が NG ボソンを\term{食べて}{eat}質量を得る、と言う。
\subsection{混合項と横波の質量項}
直交座標の2次の項
\begin{equation}
\Lag^{(2)}=-\frac14F_{\mu\nu}F^{\mu\nu}+\frac12(\partial_\mu\chi-mA_\mu)^2,\qquad m=gv
\label{eq:quadratic-mixing}
\end{equation}
において、$A_\mu$ と $\chi$ は混合 $-mA^\mu\partial_\mu\chi$ で結ばれている。
$\chi$ を古典的に積分すると、ゲージ場に対する有効作用が得られる。
$\chi$ の運動方程式は
\begin{equation}
\partial_\mu(\partial^\mu\chi-mA^\mu)=0\quad\Longrightarrow\quad\chi=m\,\frac{1}{\Box}\,\partial_\mu A^\mu
\end{equation}
であり、これを代入すると
\begin{equation}
\Lag^{(2)}_{\mathrm{eff}}=-\frac14F_{\mu\nu}F^{\mu\nu}+\frac12m^2A_\mu\left(g^{\mu\nu}-\frac{\partial^\mu\partial^\nu}{\Box}\right)A_\nu
\end{equation}
となる。運動量空間では $g^{\mu\nu}-k^\mu k^\nu/k^2$ という\textbf{横波への射影演算子}がかかった質量項で、
これはゲージ変換 $A_\mu\to A_\mu+\partial_\mu\alpha/g$（$\to k_\mu$ の付加）で不変である。
\begin{note}
$(\partial\chi-mA)_\mu=m\Bigl(\partial_\mu\frac1\Box\partial\cdot A-A_\mu\Bigr)\equiv-mA^T_\mu$ と書ける。
$A^T_\mu=(g_\mu^{\ \nu}-\partial_\mu\partial^\nu/\Box)A_\nu$ は横波成分で、射影演算子 $P_T$ の性質 $P_T^2=P_T$ から
$\frac12m^2A^T_\mu A^{T\mu}=\frac12m^2A_\mu P_T^{\mu\nu}A_\nu$。
\end{note}
ゲージ不変性は通常ゲージ場の質量項を禁止するが（第1章）、
それは\textbf{局所的な}質量項 $\frac12m^2A_\mu A^\mu$ の話である。
上の有効作用は $1/\Box$ という非局所的な（massless な極を持つ）構造を通じてゲージ不変性を保っている。
その massless な極こそ NG ボソンの伝播関数であり、NG ボソンがゲージ不変な質量項を可能にしている。
これがヒッグス機構の本質である。
\begin{remark}
同じ結果は、ゲージ場の真空偏極（自己エネルギー）$\Pi^{\mu\nu}(k)=(k^2g^{\mu\nu}-k^\mu k^\nu)\Pi(k^2)$ が
$k^2=0$ に極 $\Pi(k^2)\simeq-m^2/k^2$ を持つ、と言い換えられる
（\term{Schwinger 機構}{Schwinger mechanism}）。このとき伝播関数の極は $k^2=0$ から $k^2=m^2$ に移る。
\end{remark}
%------------------------------------------------
\section{ユニタリーゲージ}
\label{sec:unitary-gauge}
%------------------------------------------------
\subsection{定義}
NG ボソンを0にするゲージ
\begin{equation}
\pi(x)=0\qquad(\text{直交座標では }\chi=0)
\end{equation}
を\term{ユニタリーゲージ}{unitary gauge}という。
このゲージではスカラー場は実で正
\begin{equation}
\phi(x)=\frac{1}{\sqrt2}\bigl(v+h(x)\bigr)
\label{eq:unitary-gauge-phi}
\end{equation}
となり、ラグランジアンは \eqref{eq:abelian-higgs-unitary} で $B_\mu\to A_\mu$ としたもの
\begin{equation}
\Lag_{\mathrm{U}}=-\frac14F_{\mu\nu}F^{\mu\nu}+\frac12m_A^2\left(1+\frac hv\right)^2A_\mu A^\mu
+\frac12(\partial_\mu h)^2-\frac12m_h^2h^2-\lambda vh^3-\frac\lambda4h^4
\label{eq:L-unitary}
\end{equation}
になる。
\subsection{長所と短所}
\textbf{長所}：物理的な粒子（massive ベクトル $A_\mu$ と実スカラー $h$）だけが現れ、スペクトルが一目で分かる。
相互作用から直接 Feynman 則が読み取れる：
\begin{equation}
hAA:\ \ 2i\frac{m_A^2}{v}g^{\mu\nu},\qquad
hhAA:\ \ 2i\frac{m_A^2}{v^2}g^{\mu\nu}.
\label{eq:hAA-vertex}
\end{equation}
ヒッグス粒子とゲージボソンの結合は\textbf{ゲージボソンの質量の2乗に比例する}。
これはヒッグス粒子が質量の起源であることの直接の表れであり、標準模型の $hWW$, $hZZ$ 結合も同じ形になる（第6章、第12章）。
\begin{note}
$\frac{m_A^2}{v}hA_\mu A^\mu$ の頂点：$A_\mu A_\nu$ の2通りの縮約で因子 $2$、$i$ をかけて $2i\frac{m_A^2}{v}g^{\mu\nu}$。
$\frac{m_A^2}{2v^2}h^2A_\mu A^\mu$ の頂点：$h$ の縮約 $2$ と $A$ の縮約 $2$ で $2i\frac{m_A^2}{v^2}g^{\mu\nu}$。
\end{note}
\textbf{短所}：伝播関数 \eqref{eq:proca-propagator} は $k\to\infty$ で
\begin{equation}
\tilde D^{\mu\nu}(k)\sim\frac{i\,k^\mu k^\nu}{m^2k^2}\sim O(1)
\end{equation}
となり減衰しない。このため冪勘定（power counting）では理論が繰り込み不可能に見え、
ループ計算では個々のダイアグラムの発散がひどくなる（発散は最終的に相殺する）。
ヒッグス機構を伴うゲージ理論が繰り込み可能であることは、次節の $R_\xi$ ゲージを用いて
't Hooft（1971年）によって証明された。
%------------------------------------------------
\section{$R_\xi$ ゲージ}
\label{sec:Rxi}
%------------------------------------------------
\subsection{ゲージ固定関数}
\eqref{eq:quadratic-mixing} の混合項 $-mA^\mu\partial_\mu\chi$ を打ち消すように、ゲージ固定項を
\begin{equation}
\Lag_{\mathrm{gf}}=-\frac{1}{2\xi}G^2,\qquad G=\partial_\mu A^\mu+\xi m\chi
\label{eq:Rxi-gf}
\end{equation}
と選ぶ（$\xi$ は任意の実パラメータ）。これを\term{$R_\xi$ ゲージ}{$R_\xi$ gauge}（'t Hooft ゲージ）という。
展開すると
\begin{equation}
\Lag_{\mathrm{gf}}=-\frac{1}{2\xi}(\partial_\mu A^\mu)^2-m\chi\,\partial_\mu A^\mu-\frac12\xi m^2\chi^2 ,
\end{equation}
中央の項は部分積分で $+mA^\mu\partial_\mu\chi$ となり、\eqref{eq:quadratic-mixing} の混合項とちょうど打ち消す。
\begin{note}
$-m\chi\partial_\mu A^\mu=-m\partial_\mu(\chi A^\mu)+mA^\mu\partial_\mu\chi$ で、全微分は作用に寄与しない。
混合項の符号は共変微分の符号の規約に依存する。$D_\mu=\partial_\mu+igA_\mu$ の規約（Peskin--Schroeder など）では
$G=\partial_\mu A^\mu-\xi m\chi$ となる。
\end{note}
\subsection{伝播関数}
混合がなくなった2次の部分は
\begin{equation}
\Lag^{(2)}=-\frac14F_{\mu\nu}F^{\mu\nu}-\frac{1}{2\xi}(\partial_\mu A^\mu)^2+\frac12m^2A_\mu A^\mu
+\frac12(\partial_\mu\chi)^2-\frac12\xi m^2\chi^2+\frac12(\partial_\mu h)^2-\frac12m_h^2h^2
\end{equation}
であり、伝播関数は
\begin{align}
\text{ゲージ場：}&\quad
\tilde D^{\mu\nu}(k)=\frac{-i}{k^2-m^2}\left[g^{\mu\nu}-(1-\xi)\frac{k^\mu k^\nu}{k^2-\xi m^2}\right],
\label{eq:Rxi-propagator}\\
\text{NG ボソン：}&\quad
\tilde D_\chi(k)=\frac{i}{k^2-\xi m^2}
\end{align}
となる。NG ボソン $\chi$ はゲージに依存する質量 $\sqrt\xi\,m$ を持つ。
\begin{note}
$A$ の2次の部分は部分積分して $\frac12A_\mu\bigl[(\Box+m^2)g^{\mu\nu}-(1-\frac1\xi)\partial^\mu\partial^\nu\bigr]A_\nu$。
運動量空間で、射影演算子 $P_T^{\mu\nu}=g^{\mu\nu}-k^\mu k^\nu/k^2$, $P_L^{\mu\nu}=k^\mu k^\nu/k^2$ を用いると
\[
O^{\mu\nu}=-(k^2-m^2)P_T^{\mu\nu}-\frac1\xi(k^2-\xi m^2)P_L^{\mu\nu}.
\]
$P_T,P_L$ は直交射影なので逆は各成分の逆数で、伝播関数 $i(O^{-1})$ は
\[
\tilde D^{\mu\nu}=-i\left[\frac{P_T^{\mu\nu}}{k^2-m^2}+\frac{\xi P_L^{\mu\nu}}{k^2-\xi m^2}\right].
\]
$P_L$ の係数を通分すると $\frac{1}{k^2-m^2}-\frac{(1-\xi)k^2}{(k^2-m^2)(k^2-\xi m^2)}=\frac{\xi}{k^2-\xi m^2}$ となり、
\eqref{eq:Rxi-propagator} と一致する。
\end{note}
\subsection{ゴースト}
非可換でなくても、ゲージ固定関数 $G$ がスカラー場を含むため Faddeev--Popov ゴーストが相互作用する。
ゲージ変換 $\delta A_\mu=\frac1g\partial_\mu\alpha$, $\delta\chi=\alpha(v+h)$ のもとで
\begin{equation}
\delta G=\frac1g\left[\Box+\xi m^2\left(1+\frac hv\right)\right]\alpha
\end{equation}
なので、ゴーストのラグランジアンは
\begin{equation}
\Lag_{\mathrm{ghost}}=\bar c\left[-\Box-\xi m^2\left(1+\frac hv\right)\right]c
\end{equation}
となる。ゴーストは質量 $\sqrt\xi\,m$ を持ち、$h$ と結合する。
Faddeev--Popov の手続きの導出は付録Fで扱う。
\begin{note}
$\delta(h+i\chi)=i\alpha(v+h+i\chi)$ より $\delta\chi=\alpha(v+h)$。
$\delta G=\partial^\mu\delta A_\mu+\xi m\delta\chi=\frac1g\Box\alpha+\xi m\alpha(v+h)$、$m=gv$ より $\xi m v=\xi m^2/g$。
\end{note}
\subsection{ゲージパラメータの特別な値}
\begin{itemize}
    \item $\xi=1$（\term{'t Hooft--Feynman ゲージ}{'t Hooft--Feynman gauge}）：
    $\tilde D^{\mu\nu}=-ig^{\mu\nu}/(k^2-m^2)$ と最も簡単で、$\chi$ とゴーストも質量 $m$ を持つ。実用計算でよく使う。
    \item $\xi=0$（\term{Landau ゲージ}{Landau gauge}）：伝播関数は横波で、$\chi$ とゴーストは massless。
    \item $\xi\to\infty$：$\chi$ とゴーストは無限に重くなって切り離され、伝播関数は \eqref{eq:proca-propagator} に戻る。
    これがユニタリーゲージである。
\end{itemize}
有限の $\xi$ ではすべての伝播関数が $k\to\infty$ で $1/k^2$ のように減衰し、冪勘定で繰り込み可能性が明らかになる。
一方、$k^2=\xi m^2$ にゲージ依存の非物理的な極が現れるが、
物理的な振幅ではゲージ場の縦成分・$\chi$・ゴーストの寄与が打ち消し合い、$\xi$ に依存しない。
\begin{remark}
$R_\xi$ ゲージは、$\xi$ を連続的に動かすことでユニタリーゲージ（物理的・ユニタリー性が明白）と
繰り込み可能なゲージ（繰り込み可能性が明白）を補間する。
物理量が $\xi$ によらないことから、両方の性質を同時に持つことが分かる。名前の $R$ は renormalizable に由来する。
\end{remark}
%------------------------------------------------
\section{ヒッグス機構における自由度の数え上げ}
\label{sec:Higgs-counting}
%------------------------------------------------
Abelian Higgs model の物理的な自由度を数える。
\begin{table}[htbp]
\centering
\caption{Abelian Higgs model の自由度}
\label{tab:abelian-dof}
\begin{tabular}{lll}
\toprule
 & 場 & 自由度\\
\midrule
対称相（$v=0$） & 複素スカラー $\phi$ & 2\\
 & massless ゲージ場 $A_\mu$（横波2偏極） & 2\\
 & 合計 & 4\\
\midrule
ヒッグス相（$v\ne0$） & 実スカラー $h$ & 1\\
 & massive ゲージ場 $A_\mu$（3偏極） & 3\\
 & 合計 & 4\\
\bottomrule
\end{tabular}
\end{table}
全体の自由度は保存し、NG ボソン1個分の自由度がゲージ場の縦偏極に移っている。
$R_\xi$ ゲージで数えると、$A_\mu$ の4成分、$\chi$ と $h$ の2つ、ゴースト $c,\bar c$（それぞれ $-1$）で
$4+2-2=4$ となり、やはり一致する。
\begin{note}
massless ゲージ場の4成分のうち、ゲージ固定で1つ、Gauss の法則（$A_0$ の運動方程式は時間微分を含まない拘束条件）で1つが除かれ、2つが残る。
massive の場合は $\partial\cdot B=0$ で1つが除かれ3つが残る。$R_\xi$ ゲージでは、ゴーストがゲージ場の非物理的な2成分を相殺し、
残った $A_\mu$ の時間成分的な $k^2=\xi m^2$ の極と $\chi$ が互いに相殺する。
\end{note}
一般に、ゲージ群 $G$ が $H$ に破れるとき、
\begin{equation}
\text{massive なゲージボソンの数}=\text{食べられる NG ボソンの数}=\dim G-\dim H
\label{eq:massive-gauge-count}
\end{equation}
であり、$H$ に対応する $\dim H$ 個のゲージボソンは massless のまま残る。
NG ボソン以外のスカラー（$n-\dim G+\dim H$ 個）は物理的なヒッグス粒子として残る。
\begin{remark}
大域的対称性の一部だけをゲージ化した場合、ゲージ化されていない破れた生成子に対応する NG ボソンは食べられずに物理的に残る。
例えば QCD のカイラル対称性の破れでは、$\SU(2)_L\times\U(1)_Y$ の部分がゲージ化されているので、
パイ中間子は（ごくわずかに）$W,Z$ の縦偏極に混ざる。もしヒッグス場がなければ、
QCD のクォーク凝縮だけで $W,Z$ が $\sim gf_\pi$ 程度の質量を得る（テクニカラー模型の発想）。
\end{remark}
%------------------------------------------------
\section{非 Abelian Higgs model}
\label{sec:nonabelian-higgs}
%------------------------------------------------
\subsection{一般論：ゲージボソンの質量行列}
ゲージ群 $G$（生成子 $T^a$、結合定数 $g$）と、その表現に属するスカラー場 $\phi$ を考える。
\begin{equation}
\Lag=-\frac14F^a_{\mu\nu}F^{a\mu\nu}+(D_\mu\phi)^\dagger(D^\mu\phi)-V(\phi),\qquad
D_\mu\phi=(\partial_\mu-igA^a_\mu T^a)\phi .
\end{equation}
$\phi$ が VEV $\langle\phi\rangle=\phi_0$ を持つと、運動項から
\begin{equation}
(D_\mu\phi)^\dagger(D^\mu\phi)\supset g^2A^a_\mu A^{b\mu}\,\phi_0^\dagger T^aT^b\phi_0
=\frac12(M^2)_{ab}A^a_\mu A^{b\mu},
\end{equation}
\begin{equation}
(M^2)_{ab}=g^2\,\phi_0^\dagger\{T^a,T^b\}\phi_0
\label{eq:gauge-mass-matrix}
\end{equation}
というゲージボソンの\term{質量行列}{mass matrix}が得られる。
\begin{note}
$A^a_\mu A^{b\mu}$ は $a\leftrightarrow b$ で対称なので、$T^aT^b$ の対称部分 $\frac12\{T^a,T^b\}$ だけが効く。
複数の結合定数がある場合（積群）は $g^2T^aT^b\to g_ag_bT^aT^b$ とする。
\end{note}
\begin{myprop}{質量行列の性質}{gauge-mass-matrix}
$M^2$ は実対称で半正定値であり、実ベクトル $c^a$ に対して
\begin{equation}
c^a(M^2)_{ab}c^b=2g^2\bigl|c^aT^a\phi_0\bigr|^2\ge0 .
\end{equation}
等号は $c^aT^a\phi_0=0$、すなわち $c^aT^a$ が非破れ部分群 $H$ の生成子のときに限る。
したがって、massless のゲージボソンは $H$ の生成子に対応し、
massive なゲージボソンの数は破れた生成子の数 $\dim G-\dim H$ に等しい。
\end{myprop}
\begin{note}
$c^ac^b\phi_0^\dagger\{T^a,T^b\}\phi_0=2\phi_0^\dagger(c^aT^a)^2\phi_0=2|(c^aT^a)\phi_0|^2$（$c^aT^a$ はエルミート）。
実対称性：$(\phi_0^\dagger\{T^a,T^b\}\phi_0)^*=\phi_0^\dagger\{T^a,T^b\}^\dagger\phi_0=\phi_0^\dagger\{T^a,T^b\}\phi_0$。
\end{note}
NG ボソンは $\delta\phi=i\alpha^aT^a\phi_0$ の方向であり、
極座標の一般化 $\phi(x)=\exp\bigl(i\pi^{\hat a}(x)T^{\hat a}/v\bigr)\,(\phi_0+\text{動径方向の揺らぎ})$（$\hat a$ は破れた生成子。$\pi^{\hat a}$ の規格化は運動項が正準になるよう適宜選ぶ）
を用いれば、ユニタリーゲージで $\pi^{\hat a}=0$ とできる。
$R_\xi$ ゲージは $G^a=\partial_\mu A^{a\mu}+\xi\,(\text{$\pi^a$ に比例する項})$ と一般化され、
各 massive ゲージボソンに対応する NG ボソンとゴーストが質量 $\sqrt\xi\,m_a$ を持つ。
\subsection{例1：$\SU(2)$ と二重項}
$G=\SU(2)$、$\phi$ を二重項、$\phi_0=\frac{1}{\sqrt2}(0,v)^T$ とする。
$\{\sigma^a/2,\sigma^b/2\}=\frac12\delta^{ab}$ より
\begin{equation}
(M^2)_{ab}=g^2\cdot\frac12\delta^{ab}\cdot\phi_0^\dagger\phi_0=\frac{g^2v^2}{4}\delta^{ab}
\quad\Longrightarrow\quad
m_1=m_2=m_3=\frac{gv}{2}.
\end{equation}
$\SU(2)$ は完全に破れ（$H=\{1\}$）、3つのゲージボソンがすべて同じ質量を得る。
二重項の4つの実スカラーのうち3つが食べられ、1つがヒッグス粒子として残る。
3つの質量が等しいのは、\S\ref{sec:SSB-counting} で触れたカストディアル対称性 $\SU(2)_V$ による。
これに $\U(1)_Y$ を加えたものが標準模型の電弱理論であり、$W^\pm$ は $m_W=gv/2$ を得る（第6章）。
\subsection{例2：$\SU(2)$ と随伴表現（Georgi--Glashow 模型）}
$\phi^a$（$a=1,2,3$）を $\SU(2)$ の随伴表現（実3重項）とする。
$(D_\mu\phi)^a=\partial_\mu\phi^a+g\epsilon^{abc}A^b_\mu\phi^c$ であり、$\langle\phi\rangle=(0,0,v)$ とすると
\begin{equation}
\frac12(D_\mu\phi)^a(D^\mu\phi)^a\supset\frac12g^2v^2\bigl[(A^1_\mu)^2+(A^2_\mu)^2\bigr].
\end{equation}
$A^{1,2}$（すなわち $W^\pm=(A^1\mp iA^2)/\sqrt2$）が質量 $gv$ を得て、
$\langle\phi\rangle$ の方向の回転 $A^3$ は massless の「光子」として残る：$\SU(2)\to\U(1)$。
3つの実スカラーのうち2つが食べられ、1つが残る。
\begin{note}
$\epsilon^{abc}A^b\phi_0^c=v\,\epsilon^{ab3}A^b$：$a=1$ で $vA^2$、$a=2$ で $-vA^1$、$a=3$ で $0$。
2乗の和は $v^2[(A^1)^2+(A^2)^2]$。
\end{note}
\begin{remark}
Georgi--Glashow 模型（1972年）は、中性カレントの発見前に提案された電弱理論の候補であった。
中性カレント（$Z$ ボソン）の発見により、$\SU(2)\times\U(1)$ の Glashow--Weinberg--Salam 模型が選ばれた。
この模型は $\SU(2)\to\U(1)$ の破れの位相的性質から \term{'t Hooft--Polyakov モノポール}{'t Hooft--Polyakov monopole}を持つことでも知られる。
\end{remark}
\subsection{例3：$\SU(3)$ と三重項}
$\SU(3)$ の基本表現 $\phi_0=\frac{1}{\sqrt2}(0,0,v)^T$ では $\SU(3)\to\SU(2)$ と破れ、
$8-3=5$ 個のゲージボソンが質量を得る。
Gell-Mann 行列 $\lambda^a$（$T^a=\lambda^a/2$）のうち、第3成分に作用する $\lambda^{4,5,6,7}$ の4つは質量 $gv/2$、
$\lambda^8$ に対応する1つは質量 $gv/\sqrt3$ を得る。$\lambda^{1,2,3}$ に対応する3つは massless である。
\begin{note}
$\lambda^8=\frac{1}{\sqrt3}\diag(1,1,-2)$ なので $T^8\phi_0=-\frac{1}{\sqrt3}\phi_0$、
$(M^2)_{88}=2g^2|T^8\phi_0|^2=2g^2\cdot\frac13\cdot\frac{v^2}{2}=\frac{g^2v^2}{3}$。
$\lambda^4$ は $(1,3)$ 成分の $\sigma^1$ 型で、$T^4\phi_0=\frac12\cdot\frac{v}{\sqrt2}(1,0,0)^T$ より
$(M^2)_{44}=2g^2\cdot\frac{v^2}{8}=\frac{g^2v^2}{4}$。$\lambda^{5,6,7}$ も同様。
\end{note}
%------------------------------------------------
\section*{第5章のまとめ}
%------------------------------------------------
\begin{itemize}
    \item ゲージ対称性は冗長性であり、厳密には自発的に破れない（Elitzur の定理）。
    「ゲージ対称性の破れ」とは、ゲージ固定後に残る大域的部分の SSB を指す。
    \item Abelian Higgs model では、スカラー場の VEV によりゲージ場が質量 $m_A=gv$ を得る。
    NG ボソンはゲージ変換で取り除くことができ、ゲージ場の縦偏極になる（ゲージ場が NG ボソンを「食べる」）。
    \item ユニタリーゲージでは物理的な粒子だけが現れ、ヒッグス粒子とゲージボソンの結合は $m_A^2/v$ に比例する。
    $R_\xi$ ゲージでは伝播関数が $1/k^2$ で減衰し、繰り込み可能性が明白になる。物理量は $\xi$ によらない。
    \item 自由度は保存する：複素スカラー（2）+ massless ベクトル（2）= 実スカラー（1）+ massive ベクトル（3）。
    \item 非可換群 $G\to H$ では、ゲージボソンの質量行列は $(M^2)_{ab}=g^2\phi_0^\dagger\{T^a,T^b\}\phi_0$ で与えられ、
    $\dim G-\dim H$ 個のゲージボソンが質量を得る。
\end{itemize}
%================================================
\chapter{標準模型における電弱対称性の破れ}
\label{ch:EWSB}
%================================================
第5章のヒッグス機構を、標準模型の電弱ゲージ群 $\SU(2)_L\times\U(1)_Y$ に適用する。
本章の目標は、第1章で結果だけを示した
\begin{equation*}
M_W=\frac{gv}{2},\qquad M_Z=\frac{\sqrt{g^2+g'^2}\,v}{2},\qquad M_\gamma=0,\qquad e=g\sin\theta_W=g'\cos\theta_W
\end{equation*}
を導出し、ゲージボソンの質量固有状態 $W^\pm$, $Z^0$, $\gamma$ と Weinberg 角の意味を明らかにすることである。
電弱対称性の破れ\en{electroweak symmetry breaking, EWSB}は、
\term{電弱理論}{electroweak theory}（Glashow 1961, Weinberg 1967, Salam 1968）の中心部分である。
%------------------------------------------------
\section{Higgs doublet の導入}
\label{sec:higgs-doublet}
%------------------------------------------------
\subsection{ヒッグス場の量子数}
標準模型のヒッグス場は
\begin{equation}
H=\begin{pmatrix}\phi^+\\ \phi^0\end{pmatrix}
=\frac{1}{\sqrt2}\begin{pmatrix}\phi_1+i\phi_2\\ \phi_3+i\phi_4\end{pmatrix}
\sim(\mathbf1,\mathbf2)_{1}
\label{eq:H-doublet}
\end{equation}
という、カラー一重項・$\SU(2)_L$ 二重項・超電荷 $Y=1$ の複素スカラー場である
（第4章の $\Phi$ と同じもの）。成分の電荷は $Q=T_3+Y/2$ より、上成分が $+1$、下成分が $0$ である。
この選択は次の要請からほぼ一意に決まる。
\begin{enumerate}
    \item \textbf{破れのパターン}：$\SU(2)_L\times\U(1)_Y$ を $\U(1)_{\mathrm{EM}}$ に破るには、
    $T_3\ne0$ かつ $Q=0$ の成分が VEV を持たなければならない。そのためには $\SU(2)_L$ の非自明な表現で $Y\ne0$ が必要である。
    \item \textbf{Yukawa 結合}：$\overline{Q_L}Hd_R$ などのゲージ不変な Yukawa 項（第3章 \eqref{eq:Y-yukawa}）が書けるには、
    ヒッグス場は $(\mathbf1,\mathbf2)_{\pm1}$ でなければならない。これは $Y_H=1$ \eqref{eq:YH} を意味する。
    \item \textbf{$\rho$ パラメータ}：$M_W$ と $M_Z$ の比が実験とよく合う（$\rho\simeq1$）ためには、
    二重項（あるいは特殊な高次元表現）が必要である（\S\ref{sec:weinberg-angle}）。
\end{enumerate}
共役場 $\tilde H=i\sigma^2H^*=(\phi^{0*},\,-\phi^-)^T$ は $Y=-1$ の二重項である（第1章、第3章）。
\begin{remark}
電弱理論の歴史を簡単にまとめておく。
\begin{itemize}
    \item 1961年：Glashow が $\SU(2)\times\U(1)$ のゲージ群と $Z$ ボソンを導入した。ただし $W,Z$ の質量は手で加えたため、繰り込み可能性は失われていた。
    \item 1964年：ヒッグス機構（Englert--Brout, Higgs, Guralnik--Hagen--Kibble）。
    \item 1967--68年：Weinberg と Salam が、ヒッグス二重項による SSB を Glashow の模型に組み込み、レプトンの理論を完成させた。
    \item 1971年：'t Hooft が（Veltman とともに）SSB を伴うゲージ理論の繰り込み可能性を証明した。
    \item 1973年：CERN の Gargamelle 泡箱で中性カレント（$\nu_\mu e\to\nu_\mu e$ など）が発見された。
    \item 1983年：CERN の UA1, UA2 実験で $W^\pm$, $Z^0$ が発見され、質量が予言とよく一致した。
    \item 1989--2000年：LEP と SLC による $Z$ 共鳴の精密測定で、量子補正のレベルまで標準模型が検証された。
    \item 2012年：LHC の ATLAS, CMS 実験でヒッグス粒子が発見された。
\end{itemize}
\end{remark}
\subsection{ヒッグス場のラグランジアン}
ヒッグス場のラグランジアンは
\begin{equation}
\Lag_H=(D_\mu H)^\dagger(D^\mu H)-V(H),
\qquad
D_\mu H=\left(\partial_\mu-igW^i_\mu\frac{\sigma^i}{2}-i\frac{g'}{2}B_\mu\right)H
\label{eq:L-Higgs}
\end{equation}
である。\eqref{eq:D-doublet} で $Y=1$ とおくと、成分では
\begin{equation}
D_\mu H=
\begin{pmatrix}
\partial_\mu-\dfrac{i}{2}\left(gW_\mu^3+g'B_\mu\right) & -\dfrac{ig}{\sqrt2}W_\mu^+\\[2mm]
-\dfrac{ig}{\sqrt2}W_\mu^- & \partial_\mu+\dfrac{i}{2}\left(gW_\mu^3-g'B_\mu\right)
\end{pmatrix}
\begin{pmatrix}\phi^+\\ \phi^0\end{pmatrix}
\label{eq:DH-components}
\end{equation}
となる。下成分 $\phi^0$ には $gW^3-g'B$ の組み合わせだけが結合していることに注意する。
これが後で $Z$ ボソンになる組み合わせである。
%------------------------------------------------
\section{ヒッグスポテンシャル}
\label{sec:higgs-potential}
%------------------------------------------------
\subsection{ポテンシャルの一般形}
\begin{myprop}{ヒッグスポテンシャル}{higgs-potential}
1つのヒッグス二重項について、$\SU(2)_L\times\U(1)_Y$ 不変で繰り込み可能なポテンシャルは
\begin{equation}
V(H)=-\mu^2H^\dagger H+\lambda\left(H^\dagger H\right)^2
\label{eq:higgs-potential}
\end{equation}
に限られる。
\end{myprop}
\begin{note}
$H$ の2次の不変量は $H^\dagger H$ だけである。
$H^Ti\sigma^2H=\phi^+\phi^0-\phi^0\phi^+=0$ は恒等的に 0（反対称な縮約）であり、しかも $Y=2$ を持つので不変でない。
3次の項は $Y$ の和が 0 にならず禁止される（$H$ が $+1$、$H^\dagger$ が $-1$ で、3つの和は奇数）。
4次の不変量は $(H^\dagger H)^2$ に帰着する（二重項の2つの $\SU(2)$ 不変な縮約 $|H^\dagger H|^2$ と $|H^Ti\sigma^2H|^2=0$ のうち非自明なのは前者だけ）。
\end{note}
本章では第1章に合わせ、$\mu^2>0$ のとき対称性が破れる符号で書く（第4章・第5章の $\mu^2|\phi|^2$ とは $\mu^2$ の符号が逆）。
\begin{itemize}
    \item $\lambda>0$：ポテンシャルが下に有界であるための条件（\term{真空の安定性}{vacuum stability}）。
    \item $\mu^2>0$：$H=0$ が極大になり、電弱対称性が自発的に破れる。
\end{itemize}
平方完成すると
\begin{equation}
V(H)=\lambda\left(H^\dagger H-\frac{v^2}{2}\right)^2-\frac{\lambda v^4}{4},
\qquad v^2\equiv\frac{\mu^2}{\lambda}
\label{eq:higgs-potential-square}
\end{equation}
となる。
\subsection{$\mathrm O(4)$ 対称性とカストディアル対称性}
$H^\dagger H=\frac12(\phi_1^2+\phi_2^2+\phi_3^2+\phi_4^2)$ なので、
$V$ は4つの実スカラー場の回転 $\mathrm O(4)$ で不変であり、ゲージ群 $\SU(2)_L\times\U(1)_Y$ より大きな対称性を持つ。
これを見やすくするため、$2\times2$ 行列
\begin{equation}
\Sigma\equiv\left(\tilde H,\ H\right)
=\begin{pmatrix}\phi^{0*}&\phi^+\\ -\phi^-&\phi^0\end{pmatrix}
\label{eq:bidoublet}
\end{equation}
を定義する。このとき
\begin{equation}
H^\dagger H=\frac12\tr\left(\Sigma^\dagger\Sigma\right),
\end{equation}
であり、$\Sigma$ は $\SU(2)_L\times\SU(2)_R$ のもとで
\begin{equation}
\Sigma\to L\,\Sigma\,R^\dagger,\qquad L\in\SU(2)_L,\ R\in\SU(2)_R
\end{equation}
と変換する。$V$ と運動項（$g'=0$ のとき）はこの $\SU(2)_L\times\SU(2)_R\simeq\SO(4)$ で不変である。
$\U(1)_Y$ は $\SU(2)_R$ の $\sigma^3$ 方向の部分群 $R=e^{i\theta\sigma^3/2}$ として埋め込まれている。
\begin{note}
$\tilde H$ と $H$ の超電荷はそれぞれ $-1$, $+1$ なので、$\U(1)_Y$ 変換 $H\to e^{i\theta/2}H$ は
$\Sigma\to\Sigma\,\diag(e^{-i\theta/2},e^{i\theta/2})=\Sigma\,e^{-i\theta\sigma^3/2}$ と右から作用する。
これは $R^\dagger=e^{-i\theta\sigma^3/2}$ に対応する。
$\tr(\Sigma^\dagger\Sigma)=|\phi^0|^2+|\phi^-|^2+|\phi^+|^2+|\phi^0|^2=2H^\dagger H$。
\end{note}
VEV $\langle\Sigma\rangle=\frac{v}{\sqrt2}\mathbf 1$ は $L=R$ のときだけ不変なので、
\begin{equation}
\SU(2)_L\times\SU(2)_R\to\SU(2)_V
\end{equation}
と破れる。残る $\SU(2)_V$ が\term{カストディアル対称性}{custodial symmetry}であり、
$W^1,W^2,W^3$ を3重項として回すため、$g'\to0$ の極限で3つの質量が等しくなる。
$g'\ne0$ と Yukawa 結合（$y_t\ne y_b$）は $\SU(2)_R$ を陽に破り、$\rho$ パラメータに量子補正を与える（第12章）。
\subsection{真空のエネルギー}
真空でのポテンシャルの値は \eqref{eq:higgs-potential-square} より
\begin{equation}
V(\langle H\rangle)=-\frac{\lambda v^4}{4}=-\frac{m_h^2v^2}{8}\simeq-1.2\times10^{8}\ \mathrm{GeV}^4
\end{equation}
である。場の理論ではエネルギーの原点は任意だが、重力は真空のエネルギー密度にも結合し、宇宙項として働く。
観測される宇宙項（暗黒エネルギー）のエネルギー密度は $\rho_\Lambda\simeq(2.3\times10^{-3}\ \mathrm{eV})^4\simeq3\times10^{-47}\ \mathrm{GeV}^4$ であり、
電弱対称性の破れに伴う真空のエネルギーの変化よりも $10^{55}$ 倍ほど小さい。
この2つが（定数項の微調整なしに）なぜ合わないのかは、\term{宇宙項問題}{cosmological constant problem}の一部である（第14章）。
\begin{note}
$\frac{m_h^2v^2}{8}=\frac{(125.25)^2\times(246.22)^2}{8}\simeq\frac{1.569\times10^4\times6.062\times10^4}{8}\simeq1.19\times10^8\ \mathrm{GeV}^4$。
$\lambda v^4/4$ で $\lambda=m_h^2/(2v^2)$ を代入すると同じ式になる。
\end{note}
\subsection{有限温度と電弱相転移}
初期宇宙の高温では、熱浴の粒子との相互作用によりヒッグス場の有効的な質量項が
\begin{equation}
V_T(H)\simeq\left(-\mu^2+cT^2\right)H^\dagger H+\lambda\left(H^\dagger H\right)^2,
\qquad
c=\frac{3g^2+g'^2}{16}+\frac{y_t^2}{4}+\frac\lambda2
\label{eq:thermal-mass}
\end{equation}
のように持ち上がる（高温展開の主要項）。$T$ が十分高いと $H^2$ の係数が正になり、
対称性が回復する（$\langle H\rangle=0$、$W,Z$ やフェルミオンは massless）。
宇宙が冷えて
\begin{equation}
T_c\simeq\frac{\mu}{\sqrt c}
\end{equation}
を下回ると、ヒッグス場が VEV を持ち電弱対称性が破れる。
数値を入れると $c\simeq0.09+0.25+0.06\simeq0.40$、$T_c\sim140\ \mathrm{GeV}$ となる。
格子計算によれば、$m_h=125\ \mathrm{GeV}$ では相転移ではなく滑らかな\term{クロスオーバー}{crossover}であり、
その温度は $T\simeq160\ \mathrm{GeV}$ である（宇宙年齢 $\sim10^{-11}$ 秒）。
\begin{note}
高温展開では、質量 $m_i(\phi)$ を持つボソン1自由度あたり $\frac{T^2}{24}m_i^2(\phi)$、
Weyl フェルミオン1自由度あたり $\frac{T^2}{48}m_i^2(\phi)$ が有効ポテンシャルに加わる。
$W$（6自由度）、$Z$（3自由度）、トップクォーク（12自由度）、ヒッグスと NG ボソンの寄与を足し合わせると、
$H^\dagger H$ の係数として \eqref{eq:thermal-mass} の $c$ が得られる。
トップクォークの寄与 $y_t^2/4\simeq0.25$ が最大である。
\end{note}
\begin{remark}
電弱相転移が強い1次相転移であれば、宇宙のバリオン数非対称性を電弱スケールで生成する可能性があった
（\term{電弱バリオン数生成}{electroweak baryogenesis}）。
標準模型ではクロスオーバーなのでこれは働かず、新しい物理が必要になる（第14章）。
\end{remark}
%------------------------------------------------
\section{真空期待値の決定}
\label{sec:vev-determination}
%------------------------------------------------
\subsection{真空の選択}
\eqref{eq:higgs-potential-square} の最小値は
\begin{equation}
H^\dagger H=\frac{v^2}{2},\qquad v=\frac{\mu}{\sqrt\lambda}
\end{equation}
すなわち $\mathbb R^4$ の中の半径 $v$ の3次元球面 $S^3$ 上にある。
$\SU(2)_L$ は $S^3$ に推移的に作用する（$\SU(2)\simeq S^3$）ので、
ゲージ変換によって任意の真空を
\begin{equation}
\boxed{\langle H\rangle=\frac{1}{\sqrt2}\begin{pmatrix}0\\ v\end{pmatrix},\qquad v\ \text{は実で正}}
\label{eq:H-vev}
\end{equation}
の形にできる。
\begin{note}
任意の $H_0=(a,b)^T$（$|a|^2+|b|^2=v^2/2$）に対して
\[
U=\frac{\sqrt2}{v}\begin{pmatrix}b^*&a\\ -a^*&b\end{pmatrix}\in\SU(2)\quad(\det U=\tfrac{2}{v^2}(|a|^2+|b|^2)=1)
\]
とおくと $U^\dagger=\frac{\sqrt2}{v}\begin{pmatrix}b&-a\\ a^*&b^*\end{pmatrix}$ で、$U^\dagger H_0=\frac{\sqrt2}{v}(ba-ab,\ |a|^2+|b|^2)^T=(0,\ v/\sqrt2)^T$。
したがって、どの方向に VEV が生じても、ゲージ変換により下成分の実数にそろえられる。
\end{note}
VEV が下成分（$Q=0$ の成分）にあることは、\textbf{選択}ではなく\textbf{電荷の定義}と考えるべきである。
二重項が1つしかない場合、どの真空を選んでも破れずに残る $\U(1)$ が1つあり、それを電磁気の $\U(1)$ と呼ぶ。
したがって標準模型では電荷の保存は自動的である。
（2つ以上の二重項があると、電荷を破る真空も可能になり、パラメータに条件がつく。）
\subsection{$v$ の数値}
$v$ の値は Fermi 定数から決まる。第8章で導出するように、低エネルギーの荷電カレント相互作用は
\begin{equation}
\frac{G_F}{\sqrt2}=\frac{g^2}{8M_W^2}=\frac{1}{2v^2}
\quad\Longrightarrow\quad
v=\left(\sqrt2\,G_F\right)^{-1/2}\simeq246.22\ \mathrm{GeV}
\label{eq:v-value}
\end{equation}
を与える（$G_F=1.16638\times10^{-5}\ \mathrm{GeV}^{-2}$、$\mu$ 粒子の寿命から決まる）。
これが\term{電弱スケール}{electroweak scale}である。
さらにヒッグス粒子の質量 $m_h\simeq125.25\ \mathrm{GeV}$ を用いると（\S\ref{sec:higgs-dof}）
\begin{equation}
\lambda=\frac{m_h^2}{2v^2}\simeq0.129,\qquad \mu=\frac{m_h}{\sqrt2}\simeq88.6\ \mathrm{GeV}
\end{equation}
となり、ヒッグスポテンシャルのパラメータが（ツリーレベルで）すべて決まる。
\subsection{真空のまわりの展開}
真空のまわりの揺らぎは、2通りのパラメータ表示で書ける。
\begin{align}
\text{直交座標：}&\quad
H=\begin{pmatrix}G^+\\ \dfrac{1}{\sqrt2}\left(v+h+iG^0\right)\end{pmatrix},
\label{eq:H-cartesian}\\
\text{指数表示：}&\quad
H=\exp\!\left(\frac{i\pi^a\sigma^a}{v}\right)\begin{pmatrix}0\\ \dfrac{1}{\sqrt2}(v+h)\end{pmatrix}.
\label{eq:H-exponential}
\end{align}
$G^\pm$（$G^-=(G^+)^*$）と $G^0$（実）は NG ボソン、$h$ は\term{ヒッグス粒子}{Higgs boson}の場である。
直交座標で
\begin{equation}
H^\dagger H-\frac{v^2}{2}=G^+G^-+vh+\frac12h^2+\frac12(G^0)^2
\end{equation}
なので、ポテンシャルは
\begin{equation}
V=\lambda v^2h^2+\lambda vh\left(h^2+(G^0)^2+2G^+G^-\right)+\frac\lambda4\left(h^2+(G^0)^2+2G^+G^-\right)^2+\text{const.}
\label{eq:V-cartesian}
\end{equation}
となる。$m_h^2=2\lambda v^2=2\mu^2$ で、$G^\pm$, $G^0$ はポテンシャルの中では massless である。
\begin{note}
$H^\dagger H=G^-G^++\frac12\left[(v+h)^2+(G^0)^2\right]$。
$V=\lambda X^2+\text{const.}$（$X=H^\dagger H-\frac{v^2}{2}$）に $X=\frac12\left(2vh+h^2+(G^0)^2+2G^+G^-\right)$ を代入する。
\end{note}
指数表示では、$\pi^a$ は $\SU(2)_L$ のゲージ変換のパラメータそのものなので、
ゲージ変換 $H\to\exp(-i\pi^a\sigma^a/v)H$ で取り除ける（$\pi^a$ の運動項が正準になるよう $\sigma^a/v$ と規格化した）。これが\term{ユニタリーゲージ}{unitary gauge}
\begin{equation}
H=\frac{1}{\sqrt2}\begin{pmatrix}0\\ v+h(x)\end{pmatrix}
\label{eq:H-unitary}
\end{equation}
である。以下、ゲージボソンの質量を求めるにはこのゲージを用いる。
%------------------------------------------------
\section{$\SU(2)_L\times\U(1)_Y\to\U(1)_{\mathrm{EM}}$}
\label{sec:EW-breaking-pattern}
%------------------------------------------------
\subsection{破れない生成子}
4つの生成子 $T^1,T^2,T^3,Y/2$ を \eqref{eq:H-vev} に作用させる（\S\ref{sec:SSB-counting}）。
\begin{equation}
T^{1,2}\langle H\rangle\ne0,\qquad
T^3\langle H\rangle=-\frac12\langle H\rangle,\qquad
\frac Y2\langle H\rangle=+\frac12\langle H\rangle .
\end{equation}
$T^3$ と $Y/2$ の固有値が打ち消し合うので、
\begin{equation}
Q\langle H\rangle=\left(T^3+\frac Y2\right)\langle H\rangle=0 .
\end{equation}
したがって $e^{i\alpha Q}$ は真空を変えず、\textbf{電磁気の $\U(1)_{\mathrm{EM}}$ が破れずに残る}：
\begin{equation}
\boxed{\SU(2)_L\times\U(1)_Y\ \longrightarrow\ \U(1)_{\mathrm{EM}},\qquad Q=T_3+\frac Y2}
\label{eq:EW-breaking}
\end{equation}
破れた生成子は3つ（例えば $T^1$, $T^2$, $T^3-Y/2$）で、3つの NG ボソン $G^\pm$, $G^0$ が現れる。
これらはそれぞれ $W^\pm$, $Z$ に食べられる。$G^\pm$ の電荷 $\pm1$ が $W^\pm$ の電荷と一致することに注意。
\begin{note}
$G^+=\phi^+$ は上成分なので電荷 $+1$、$G^0=\sqrt2\,\Im\phi^0$ は中性。
破れた生成子のうち $T^\pm=T^1\pm iT^2$ は電荷 $\pm1$ を運ぶ（$[Q,T^\pm]=\pm T^\pm$）ので、対応するゲージボソンも電荷 $\pm1$ を持つ。
\end{note}
\subsection{$\U(1)_{\mathrm{EM}}$ は $\U(1)_Y$ ではない}
破れずに残る $\U(1)$ は、元のゲージ群の $\U(1)_Y$ そのものではなく、
$\SU(2)_L$ の $T_3$ と $\U(1)_Y$ を混ぜた「斜め」の部分群である。
このため、光子は $W^3$ と $B$ の線形結合となり（\S\ref{sec:W-Z-photon}）、
電荷 $Q$ は $T_3$ と $Y$ の両方に依存する。
\begin{remark}
仮に VEV が上成分（$Q=+1$）にあったとすると、$Q$ ではなく $T_3-Y/2$ が保存し、
それを改めて「電荷」と呼べば物理は同じである。
すなわち「どちらの成分が中性か」は、VEV の方向で決まる。
電荷の割り当て $Q=T_3+Y/2$ は、VEV が下成分にあるという規約と表裏一体である。
\end{remark}
%------------------------------------------------
\section{ゲージ場の質量行列}
\label{sec:EW-mass-matrix}
%------------------------------------------------
\subsection{運動項への代入}
ユニタリーゲージ \eqref{eq:H-unitary} を \eqref{eq:L-Higgs} の運動項に代入する。
$\partial_\mu$ は $h$ にのみ作用し、
\begin{align}
\left(gW^i_\mu\sigma^i+g'B_\mu\right)\begin{pmatrix}0\\1\end{pmatrix}
&=\begin{pmatrix}gW^3_\mu+g'B_\mu & g(W^1_\mu-iW^2_\mu)\\ g(W^1_\mu+iW^2_\mu) & -gW^3_\mu+g'B_\mu\end{pmatrix}
\begin{pmatrix}0\\1\end{pmatrix}
=\begin{pmatrix}g(W^1_\mu-iW^2_\mu)\\ -gW^3_\mu+g'B_\mu\end{pmatrix}
\end{align}
なので、
\begin{equation}
D_\mu H=\frac{1}{\sqrt2}\begin{pmatrix}-\dfrac{ig}{2}(W^1_\mu-iW^2_\mu)(v+h)\\[2mm] \partial_\mu h+\dfrac{i}{2}(gW^3_\mu-g'B_\mu)(v+h)\end{pmatrix}.
\end{equation}
したがって
\begin{equation}
(D_\mu H)^\dagger(D^\mu H)
=\frac12(\partial_\mu h)^2
+\frac{(v+h)^2}{8}\Bigl[g^2\bigl((W^1_\mu)^2+(W^2_\mu)^2\bigr)+\bigl(gW^3_\mu-g'B_\mu\bigr)^2\Bigr]
\label{eq:DH-squared}
\end{equation}
となる（$(W_\mu)^2\equiv W_\mu W^\mu$）。
\begin{note}
上成分の絶対値の2乗：$\frac12\cdot\frac{g^2}{4}|W^1-iW^2|^2(v+h)^2=\frac{g^2(v+h)^2}{8}\left[(W^1)^2+(W^2)^2\right]$。
下成分：
\[
\frac12\left|\partial h+\frac i2(gW^3-g'B)(v+h)\right|^2=\frac12(\partial h)^2+\frac{(v+h)^2}{8}(gW^3-g'B)^2
\]
（$\partial h$ は実、第2項は純虚数なので交差項は消える）。
\end{note}
\subsection{質量行列}
$h=0$ とおくと、ゲージ場の質量項は $\frac12\mathcal V_\mu^T M^2\mathcal V^\mu$（$\mathcal V=(W^1,W^2,W^3,B)^T$）の形で、
\begin{equation}
M^2=\frac{v^2}{4}
\begin{pmatrix}
g^2&0&0&0\\
0&g^2&0&0\\
0&0&g^2&-gg'\\
0&0&-gg'&g'^2
\end{pmatrix}
\label{eq:EW-mass-matrix}
\end{equation}
となる。これは第5章の一般公式 $(M^2)_{ab}=g_ag_b\langle H\rangle^\dagger\{T^a,T^b\}\langle H\rangle$ からも得られる。
\begin{note}
$\langle H\rangle^\dagger\{T^a,T^b\}\langle H\rangle$ を計算する。$\langle H\rangle^\dagger\langle H\rangle=\frac{v^2}{2}$ で、
$a,b\in\{1,2,3\}$ では $\{T^a,T^b\}=\frac12\delta^{ab}$ より $\frac{v^2}{4}\delta^{ab}$。
$a=3$, $b=Y$（$T^Y=Y/2=\frac12$）では $\{T^3,\frac12\}=T^3$ で、$\langle H\rangle^\dagger T^3\langle H\rangle=-\frac12\cdot\frac{v^2}{2}=-\frac{v^2}{4}$。
$a=b=Y$ では $2\cdot\frac14\cdot\frac{v^2}{2}=\frac{v^2}{4}$。$a=1,2$, $b=Y$ では $\langle H\rangle^\dagger T^{1,2}\langle H\rangle=0$。
結合定数をかけると \eqref{eq:EW-mass-matrix} を得る。
\end{note}
\subsection{固有値}
$W^1$, $W^2$ はすでに対角で、固有値 $g^2v^2/4$。
$(W^3,B)$ の $2\times2$ ブロック
\begin{equation}
\frac{v^2}{4}\begin{pmatrix}g^2&-gg'\\ -gg'&g'^2\end{pmatrix}
\end{equation}
は、行列式 $g^2g'^2-g^2g'^2=0$、トレース $g^2+g'^2$ なので、固有値は
\begin{equation}
0\qquad\text{と}\qquad\frac{(g^2+g'^2)v^2}{4}
\end{equation}
である。固有値 0 の存在は、$\U(1)_{\mathrm{EM}}$ が破れていないことの直接の帰結である
（第5章の命題 \ref{prop:gauge-mass-matrix}）。
固有ベクトルはそれぞれ $(g',g)$ と $(g,-g')$ に比例する。
\subsection{行列表示による導出とカストディアル対称性}
\S\ref{sec:higgs-potential} の行列 $\Sigma=(\tilde H,H)$ を用いると、質量項の構造がより明確になる。
指数表示 \eqref{eq:H-exponential} に対応して
\begin{equation}
\Sigma=\frac{v+h}{\sqrt2}\,U,\qquad U=\exp\!\left(\frac{i\pi^a\sigma^a}{v}\right)\in\SU(2)
\end{equation}
と書け、$\tr(\Sigma^\dagger\Sigma)=(v+h)^2$, $(D_\mu H)^\dagger(D^\mu H)=\frac12\tr\bigl[(D_\mu\Sigma)^\dagger D^\mu\Sigma\bigr]$ より
\begin{equation}
\Lag_H=\frac12(\partial_\mu h)^2+\frac{(v+h)^2}{4}\tr\bigl[(D_\mu U)^\dagger D^\mu U\bigr]-V(h),
\label{eq:EW-chiral-L}
\end{equation}
\begin{equation}
D_\mu U=\partial_\mu U-igW^a_\mu\frac{\sigma^a}{2}U+ig'B_\mu U\frac{\sigma^3}{2}
\end{equation}
となる。$\SU(2)_L$ は $U$ に左から、$\U(1)_Y$ は $\SU(2)_R$ の部分群として右から作用する。
\begin{note}
$\U(1)_Y$ のもとで $U\to Ue^{-i\theta\sigma^3/2}$, $B_\mu\to B_\mu+\frac{1}{g'}\partial_\mu\theta$ なので、
$\partial_\mu U$ から出る $-iU\frac{\sigma^3}{2}\partial_\mu\theta$ を $+ig'U\frac{\sigma^3}{2}\cdot\frac{1}{g'}\partial_\mu\theta$ が打ち消す。
$\tr[(D\Sigma)^\dagger D\Sigma]$ は $\Sigma$ の2つの列 $\tilde H$, $H$ の寄与の和で、両者は等しいので $2(DH)^\dagger DH$ になる。
\end{note}
ユニタリーゲージ $U=1$ では $D_\mu U=-\frac i2\left(gW^a_\mu\sigma^a-g'B_\mu\sigma^3\right)$ なので、
$\tr(\sigma^a\sigma^b)=2\delta^{ab}$ を用いて
\begin{equation}
\frac{(v+h)^2}{4}\tr\bigl[(D_\mu U)^\dagger D^\mu U\bigr]
=\frac{(v+h)^2}{8}\Bigl[g^2\bigl((W^1)^2+(W^2)^2\bigr)+\bigl(gW^3-g'B\bigr)^2\Bigr]
\end{equation}
となり、\eqref{eq:DH-squared} を再現する。
この形から次のことが読み取れる。
\begin{itemize}
    \item $g'=0$ のとき質量項は $\frac{g^2v^2}{8}W^a_\mu W^{a\mu}$ となり、$W^a$ の3成分が同じ質量を持つ。
    これは $W^a$ が $\SU(2)_V$（カストディアル対称性）の3重項であることの反映である。
    \item $g'\ne0$ は $\sigma^3$ を通じて $\SU(2)_R$ を破り、$W^3$ だけが $B$ と混合する。
    \item $v^2\tr[(D_\mu U)^\dagger D^\mu U]/4$ はヒッグス粒子がなくても書けるゲージ不変な質量項であり、
    \term{電弱カイラルラグランジアン}{electroweak chiral Lagrangian}の最低次の項である。
    QCD のパイ中間子の有効理論（$f_\pi^2\tr[\partial U^\dagger\partial U]/4$）と全く同じ形をしており、$v$ が $f_\pi$ の役割を果たす。
\end{itemize}
\begin{myex}{$W^3$--$B$ 混合の2準位系としての見方}{W3B-mixing}
$(W^3,B)$ の質量行列を、「対角成分 $\frac{g^2v^2}{4}$, $\frac{g'^2v^2}{4}$ と非対角成分 $-\frac{gg'v^2}{4}$ を持つ2準位系」と見て、
混合角が $\tan\theta_W=g'/g$ で与えられることを確かめよ。
\tcblower
実対称な $2\times2$ 行列 $\begin{pmatrix}a&c\\ c&b\end{pmatrix}$ を対角化する回転角 $\theta$ は $\tan2\theta=\frac{2c}{a-b}$ を満たす
（$\begin{pmatrix}Z\\A\end{pmatrix}=R(\theta)\begin{pmatrix}W^3\\B\end{pmatrix}$ の向きでは $\tan2\theta=-\frac{2c}{a-b}$）。
$a=g^2$, $b=g'^2$, $c=-gg'$（共通因子 $v^2/4$ を除く）より
\[
\tan2\theta=\frac{2gg'}{g^2-g'^2}=\frac{2(g'/g)}{1-(g'/g)^2},
\]
これは $\tan\theta=g'/g$ のときの倍角公式そのものである。
行列式が 0 なので一方の固有値は 0 となり、これが「混合しても光子が massless のまま」である理由である。
混合角が質量の値によらず結合定数の比だけで決まるのは、質量行列が $(g,-g')$ というベクトルの直積
$\frac{v^2}{4}\begin{pmatrix}g\\-g'\end{pmatrix}\begin{pmatrix}g&-g'\end{pmatrix}$ の形（ランク1）をしているからである。
\end{myex}
%------------------------------------------------
\section{$W^\pm$、$Z^0$、光子の導出}
\label{sec:W-Z-photon}
%------------------------------------------------
\subsection{荷電ゲージボソン}
$W^1$, $W^2$ は質量が縮退しているので、その任意の線形結合が質量固有状態である。
電荷の固有状態として
\begin{equation}
W^\pm_\mu=\frac{1}{\sqrt2}\left(W^1_\mu\mp iW^2_\mu\right)
\label{eq:Wpm}
\end{equation}
を選ぶ（第2章 \eqref{eq:Wmatrix}）。$(W^1)^2+(W^2)^2=2W^+_\mu W^{-\mu}$ なので、
\begin{equation}
\frac{g^2v^2}{8}\left[(W^1)^2+(W^2)^2\right]=\frac{g^2v^2}{4}W^+_\mu W^{-\mu}\equiv M_W^2W^+_\mu W^{-\mu},
\end{equation}
\begin{equation}
\boxed{M_W=\frac{gv}{2}}
\label{eq:MW}
\end{equation}
を得る。複素ベクトル場の質量項は $M^2W^+W^-$ と規格化される（実ベクトル場の $\frac12M^2V^2$ に対応）。
\begin{note}
$W^+W^-=\frac12(W^1-iW^2)(W^1+iW^2)=\frac12\left[(W^1)^2+(W^2)^2\right]$（$W^1W^2$ の交差項は打ち消す）。
\end{note}
\subsection{中性ゲージボソン}
$(W^3,B)$ ブロックの固有ベクトルを規格化して
\begin{equation}
Z_\mu=\frac{gW^3_\mu-g'B_\mu}{\sqrt{g^2+g'^2}},\qquad
A_\mu=\frac{g'W^3_\mu+gB_\mu}{\sqrt{g^2+g'^2}}
\label{eq:Z-A-def}
\end{equation}
と定義する。これは直交変換なので、運動項 $-\frac14W^3_{\mu\nu}W^{3\mu\nu}-\frac14B_{\mu\nu}B^{\mu\nu}$ の
可換な部分は $-\frac14Z_{\mu\nu}Z^{\mu\nu}-\frac14F_{\mu\nu}F^{\mu\nu}$ となり、正準規格化が保たれる。
質量項は
\begin{equation}
\frac{v^2}{8}\left(gW^3_\mu-g'B_\mu\right)^2=\frac{(g^2+g'^2)v^2}{8}Z_\mu Z^\mu\equiv\frac12M_Z^2Z_\mu Z^\mu,
\end{equation}
\begin{equation}
\boxed{M_Z=\frac{\sqrt{g^2+g'^2}\,v}{2},\qquad M_\gamma=0}
\label{eq:MZ}
\end{equation}
となる。$A_\mu$ は質量項に現れず、massless の光子である。
\subsection{ユニタリーゲージでのラグランジアン}
\eqref{eq:DH-squared} を質量固有状態で書くと
\begin{equation}
(D_\mu H)^\dagger(D^\mu H)=\frac12(\partial_\mu h)^2
+\left(M_W^2W^+_\mu W^{-\mu}+\frac12M_Z^2Z_\mu Z^\mu\right)\left(1+\frac hv\right)^2
\label{eq:higgs-gauge-unitary}
\end{equation}
となる。$h$ についての展開から、ヒッグス粒子とゲージボソンの結合
\begin{equation}
\begin{aligned}
hW^+W^-&:\ \ i\frac{2M_W^2}{v}g^{\mu\nu}=igM_Wg^{\mu\nu}, &\qquad
hZZ&:\ \ i\frac{2M_Z^2}{v}g^{\mu\nu}=i\frac{g}{\cos\theta_W}M_Zg^{\mu\nu},\\
hhW^+W^-&:\ \ i\frac{2M_W^2}{v^2}g^{\mu\nu}, &\qquad
hhZZ&:\ \ i\frac{2M_Z^2}{v^2}g^{\mu\nu}
\end{aligned}
\label{eq:hVV-couplings}
\end{equation}
が読み取れる。結合はゲージボソンの質量の2乗に比例し、光子との直接の結合（$h\gamma\gamma$）はツリーレベルでは存在しない
（$h\to\gamma\gamma$ はループを通じて起こる。第12章）。
\begin{note}
$hW^+W^-$：$M_W^2\cdot\frac{2h}{v}W^+_\mu W^{-\mu}$ で、$W^+$ と $W^-$ は別の場なので縮約の因子は 1。頂点は $i\frac{2M_W^2}{v}g^{\mu\nu}$。
$M_W=gv/2$ より $\frac{2M_W^2}{v}=gM_W$。
$hZZ$：$\frac12M_Z^2\cdot\frac{2h}{v}Z^2$ で、同種の $Z$ 2つの縮約で因子 2。頂点は $i\frac{2M_Z^2}{v}g^{\mu\nu}$。
$hhVV$：$\frac{h^2}{v^2}$ の係数に、$h$ の縮約 $2$ と（$Z$ なら）$Z$ の縮約 $2\times\frac12$ がかかる。
\end{note}
\subsection{ゲージボソンの自己相互作用}
$\SU(2)_L$ は非可換なので、ゲージ場の運動項 $-\frac14W^i_{\mu\nu}W^{i\mu\nu}$ は3点・4点の自己相互作用を含む。
$W^i_{\mu\nu}=\partial_\mu W^i_\nu-\partial_\nu W^i_\mu+g\epsilon^{ijk}W^j_\mu W^k_\nu$ を代入すると、
\begin{equation}
-\frac14W^i_{\mu\nu}W^{i\mu\nu}
=-\frac14\left(\partial_\mu W^i_\nu-\partial_\nu W^i_\mu\right)^2
\underbrace{-g\epsilon^{ijk}(\partial_\mu W^i_\nu)W^{j\mu}W^{k\nu}}_{\Lag_3}
\underbrace{-\frac{g^2}{4}\left(\epsilon^{ijk}W^j_\mu W^k_\nu\right)^2}_{\Lag_4}
\end{equation}
である。これを質量固有状態で書き直す。
\paragraph{3点結合}
\begin{equation}
\boxed{
\Lag_{WWV}=ig_V\Bigl[\left(W^+_{\mu\nu}W^{-\mu}-W^-_{\mu\nu}W^{+\mu}\right)V^\nu+V_{\mu\nu}W^{+\mu}W^{-\nu}\Bigr],
\qquad g_\gamma=e,\quad g_Z=g\cos\theta_W
}
\label{eq:WWV}
\end{equation}
ここで $W^\pm_{\mu\nu}=\partial_\mu W^\pm_\nu-\partial_\nu W^\pm_\mu$, $V_{\mu\nu}=\partial_\mu V_\nu-\partial_\nu V_\mu$（$V=\gamma,Z$）。
\begin{note}
$\epsilon^{ijk}a^ib^jc^k=\det(a,b,c)$（行ベクトルを並べた行列式）である。
成分 $(1,2)$ を $(+,-)$ に変える変換 $X^1=\frac{1}{\sqrt2}(X^++X^-)$, $X^2=\frac{i}{\sqrt2}(X^+-X^-)$ の行列式は $\frac12(-i-i)=-i$ なので、
\[
\epsilon^{ijk}a^ib^jc^k=-i\det\begin{pmatrix}a^+&a^-&a^3\\ b^+&b^-&b^3\\ c^+&c^-&c^3\end{pmatrix}.
\]
$a=\partial_\mu W_\nu$, $b=W^\mu$, $c=W^\nu$ として展開し、$\mu\leftrightarrow\nu$ の反対称化で場の強さにまとめると
\[
\Lag_3=ig\Bigl[\left(W^+_{\mu\nu}W^{-\mu}-W^-_{\mu\nu}W^{+\mu}\right)W^{3\nu}+W^3_{\mu\nu}W^{+\mu}W^{-\nu}\Bigr].
\]
最後に $W^3=\cos\theta_WZ+\sin\theta_WA$ を代入し、$g\sin\theta_W=e$ を用いる。
$\Lag_{WWV}$ が実であることは、複素共役で $W^+\leftrightarrow W^-$ と $V_{\mu\nu}$ の反対称性から確かめられる。
\end{note}
\paragraph{4点結合}
同様に
\begin{align}
\Lag_4&=-\frac{g^2}{2}\Bigl[\left(W^+_\mu W^{-\mu}\right)^2-\left(W^+_\mu W^{+\mu}\right)\left(W^-_\nu W^{-\nu}\right)\Bigr]\notag\\
&\quad-g^2\Bigl[\left(W^+_\mu W^{-\mu}\right)\left(W^3_\nu W^{3\nu}\right)-\left(W^+_\mu W^{3\mu}\right)\left(W^-_\nu W^{3\nu}\right)\Bigr]
\label{eq:quartic-gauge}
\end{align}
となる。第2行に $W^3=\cos\theta_WZ+\sin\theta_WA$ を代入すると、
$WW\gamma\gamma$（$e^2$）、$WWZ\gamma$（$eg\cos\theta_W$）、$WWZZ$（$g^2\cos^2\theta_W$）の4点結合が得られる。
\begin{note}
$\SU(2)$ 空間の内積を $X\cdot Y=X^+Y^-+X^-Y^++X^3Y^3$ と書くと、
$\left(\epsilon^{ijk}W^j_\mu W^k_\nu\right)^2=(W_\mu\cdot W^\mu)^2-(W_\mu\cdot W_\nu)(W^\mu\cdot W^\nu)$（ベクトル積の2乗の公式）。
$W^3=0$ の部分は $W_\mu\cdot W^\mu=2W^+\!\cdot W^-$、
$(W_\mu\cdot W_\nu)^2=2(W^+\!\cdot W^+)(W^-\!\cdot W^-)+2(W^+\!\cdot W^-)^2$ から第1行が出る
（ここで「$\cdot$」はローレンツ添字の縮約）。$W^3$ を1対含む項から第2行が出る。
\end{note}
\paragraph{物理的な意味}
\begin{itemize}
    \item 光子との結合 \eqref{eq:WWV} の第1項は、電荷 $+1$ の場 $W^+$ の最小結合（$\partial_\mu\to\partial_\mu-ieA_\mu$）から来る。
    実際、$W^\pm$ を含む項は $\U(1)_{\mathrm{EM}}$ の共変微分 $\mathcal D_\mu W^\pm_\nu=(\partial_\mu\mp ieA_\mu)W^\pm_\nu$ を用いて
    \[
    -\frac12\left(\mathcal D_\mu W^+_\nu-\mathcal D_\nu W^+_\mu\right)^\dagger\left(\mathcal D^\mu W^{+\nu}-\mathcal D^\nu W^{+\mu}\right)
    +ieF_{\mu\nu}W^{+\mu}W^{-\nu}+\cdots
    \]
    とまとめられ、残る $ieF_{\mu\nu}W^{+\mu}W^{-\nu}$ が $W$ の\textbf{異常磁気モーメント}的な項になる。
    \item その結果、$W$ ボソンの磁気モーメントは $\mu_W=\frac{e}{M_W}$（g 因子 $g_W=2$）、
    電気四重極モーメントは $Q_W=-\frac{e}{M_W^2}$ と予言される。これらは LEP や Tevatron の $W$ 対生成で検証されている。
    \item 4点結合 $WW\gamma\gamma$ は、$\U(1)_{\mathrm{EM}}$ の最小結合の2次の項（スカラー QED の seagull 項に対応）である。
    \item $e^+e^-\to W^+W^-$ では、$\nu$ の $t$ チャネル交換と、$\gamma$, $Z$ の $s$ チャネル交換（\eqref{eq:WWV} の3点結合）
    が個別にはエネルギーとともに増大し、和をとると打ち消し合う。これは LEP2 で精密に確かめられた
    ゲージ構造の直接の証拠である（第12章）。
\end{itemize}
\begin{table}[htbp]
\centering
\caption{電弱ボソンの相互作用（ユニタリーゲージ）}
\label{tab:EW-boson-vertices}
\begin{tabular}{lll}
\toprule
頂点 & 結合の強さ & 起源\\
\midrule
$W^+W^-\gamma$, $W^+W^-Z$ & $e$, $g\cos\theta_W$ & $-\frac14W^i_{\mu\nu}W^{i\mu\nu}$（非可換性）\\
$W^+W^-W^+W^-$ & $g^2$ & 同上\\
$W^+W^-\gamma\gamma$, $W^+W^-Z\gamma$, $W^+W^-ZZ$ & $e^2$, $eg\cos\theta_W$, $g^2\cos^2\theta_W$ & 同上\\
$hW^+W^-$, $hZZ$ & $2M_W^2/v$, $2M_Z^2/v$ & $(D_\mu H)^\dagger D^\mu H$\\
$hhW^+W^-$, $hhZZ$ & $2M_W^2/v^2$, $2M_Z^2/v^2$ & 同上\\
$hhh$, $hhhh$ & $3m_h^2/v$, $3m_h^2/v^2$ & $V(H)$\\
\bottomrule
\end{tabular}
\end{table}
$ZZZ$, $ZZ\gamma$, $Z\gamma\gamma$, $\gamma\gamma\gamma$ のような中性ボソンだけの3点結合はツリーレベルでは存在しない
（$\epsilon^{ijk}$ により3つの $\SU(2)$ 添字が異なる必要があり、$W^3$ を2つ以上含む項がない）。
\subsection{共変微分の質量固有状態による表示}
フェルミオンとの結合を見るため、一般の表現（$T^a$, 超電荷 $Y$）に対する共変微分の電弱部分を
質量固有状態で書き直す。$W^1T^1+W^2T^2=\frac{1}{\sqrt2}(W^+T^++W^-T^-)$ と、\eqref{eq:Z-A-def} の逆変換
$W^3=\cos\theta_W Z+\sin\theta_W A$, $B=-\sin\theta_W Z+\cos\theta_W A$（$\theta_W$ は次節で定義）を用いると、
\begin{equation}
\boxed{
D_\mu=\partial_\mu-\frac{ig}{\sqrt2}\left(W^+_\mu T^++W^-_\mu T^-\right)
-\frac{ig}{\cos\theta_W}Z_\mu\left(T^3-\sin^2\theta_W\,Q\right)
-ieA_\mu Q
}
\label{eq:D-mass-basis}
\end{equation}
となる。ここで $e=g\sin\theta_W=g'\cos\theta_W$。
\begin{note}
中性部分：$gW^3T^3+g'B\frac Y2=Z\left(g\cos\theta_WT^3-g'\sin\theta_W\frac Y2\right)+A\left(g\sin\theta_WT^3+g'\cos\theta_W\frac Y2\right)$。
$A$ の係数は $e(T^3+\frac Y2)=eQ$。
$Z$ の係数は $\frac Y2=Q-T^3$ と $g'=g\tan\theta_W$ を使って
$g\cos\theta_WT^3-g\frac{\sin^2\theta_W}{\cos\theta_W}(Q-T^3)=\frac{g}{\cos\theta_W}\left[(\cos^2\theta_W+\sin^2\theta_W)T^3-\sin^2\theta_WQ\right]$。
荷電部分：$T^1W^1+T^2W^2$ に $W^1=\frac{1}{\sqrt2}(W^++W^-)$, $W^2=\frac{i}{\sqrt2}(W^+-W^-)$ を代入すると
$\frac{1}{\sqrt2}\left[W^+(T^1+iT^2)+W^-(T^1-iT^2)\right]$。
\end{note}
\eqref{eq:D-mass-basis} は第8章で荷電カレント・中性カレント・電磁カレントを導く出発点になる。
光子は電荷 $Q$ に結合定数 $e$ で結合し、$Z$ は $T^3-\sin^2\theta_WQ$ に結合定数 $g/\cos\theta_W$ で結合する。
%------------------------------------------------
\section{Weinberg 角}
\label{sec:weinberg-angle}
%------------------------------------------------
\subsection{定義}
\eqref{eq:Z-A-def} の回転角を\term{Weinberg 角}{Weinberg angle}（\term{弱混合角}{weak mixing angle}）$\theta_W$ と呼ぶ：
\begin{equation}
\cos\theta_W=\frac{g}{\sqrt{g^2+g'^2}},\qquad
\sin\theta_W=\frac{g'}{\sqrt{g^2+g'^2}},\qquad
\tan\theta_W=\frac{g'}{g}.
\label{eq:weinberg-angle}
\end{equation}
このとき
\begin{equation}
\begin{pmatrix}Z_\mu\\ A_\mu\end{pmatrix}
=\begin{pmatrix}\cos\theta_W&-\sin\theta_W\\ \sin\theta_W&\cos\theta_W\end{pmatrix}
\begin{pmatrix}W^3_\mu\\ B_\mu\end{pmatrix}
\end{equation}
であり、第1章 \eqref{eq:ZA-preview} と一致する。
Weinberg 角は2つのゲージ結合定数の比で決まる量であり、ヒッグス機構によって初めて物理的な意味（$Z$ と $\gamma$ の混合角）を持つ。
\subsection{質量の関係と $\rho$ パラメータ}
\eqref{eq:MW}, \eqref{eq:MZ} より
\begin{equation}
\boxed{M_W=M_Z\cos\theta_W}
\label{eq:MW-MZ}
\end{equation}
が成り立つ。これは標準模型の重要な予言であり、
\begin{equation}
\rho\equiv\frac{M_W^2}{M_Z^2\cos^2\theta_W}=1\qquad(\text{ツリーレベル})
\label{eq:rho}
\end{equation}
と表される。$\rho=1$ はカストディアル対称性（\S\ref{sec:higgs-potential}）の帰結である。
$g'\to0$ の極限で $W^3$ は $W^{1,2}$ と同じ質量を持ち、$g'\ne0$ による $W^3$--$B$ 混合が
$Z$ の質量を $1/\cos\theta_W$ 倍に押し上げる。
\begin{myex}{一般のヒッグス表現と $\rho$ パラメータ}{rho-general}
弱アイソスピン $T_i$、超電荷 $Y_i$ のスカラー場が、中性成分（$T_{3}=-Y_i/2$）に VEV $v_i$ を持つとする。
ツリーレベルの $\rho$ パラメータは
\[
\rho=\frac{\sum_i c_i\left[T_i(T_i+1)-\frac{Y_i^2}{4}\right]v_i^2}{\sum_i\frac{Y_i^2}{2}v_i^2}
\]
（複素場で $c_i=1$、$Y=0$ の実場で $c_i=\frac12$）で与えられる。二重項と3重項について調べよ。
\tcblower
\begin{itemize}
    \item 二重項 $T=\frac12$, $Y=1$：分子 $\frac34-\frac14=\frac12$、分母 $\frac12$ で $\rho=1$。二重項を何個加えても $\rho=1$。
    \item 複素3重項 $T=1$, $Y=2$：分子 $2-1=1$、分母 $2$ で $\rho=\frac12$。
    \item 実3重項 $T=1$, $Y=0$：分母に寄与せず、分子だけを増やす（$\rho>1$）。
\end{itemize}
実験値は $\rho\simeq1$（量子補正を除いて $10^{-3}$ 程度の精度）なので、3重項の VEV は数 GeV 以下に制限される。
二重項以外で $\rho=1$ を与える最小の表現は $T=3$, $Y=4$ の7重項である。
\end{myex}
\begin{note}
一般の場合、$W^\pm$ の質量は $\frac{g^2}{4}\sum_i2c_i[T_i(T_i+1)-T_{3i}^2]v_i^2$、
$Z$ の質量は $\frac{g^2+g'^2}{4}\sum_i4T_{3i}^2v_i^2$ から得られる（$T^\pm$ と $T^3$ の行列要素を計算する）。
$T_{3i}=-Y_i/2$ を代入して比をとると上式になる。二重項では $M_W^2=\frac{g^2}{4}\cdot2\cdot\frac12v^2=\frac{g^2v^2}{4}$ を再現する。
\end{note}
\subsection{$\sin^2\theta_W$ の定義と数値}
ツリーレベルでは \eqref{eq:weinberg-angle} と \eqref{eq:MW-MZ} は同値だが、量子補正を含めると両者は異なり、
$\sin^2\theta_W$ にはいくつかの定義がある。
\begin{itemize}
    \item \term{on-shell 定義}{on-shell definition}：$\sin^2\theta_W\equiv1-M_W^2/M_Z^2\simeq0.2232$（$M_W=80.37$, $M_Z=91.19\ \mathrm{GeV}$）。
    \item \term{$\overline{\mathrm{MS}}$ 定義}{$\overline{\mathrm{MS}}$ scheme}：$\sin^2\hat\theta_W(M_Z)\equiv\hat g'^2/(\hat g^2+\hat g'^2)\simeq0.2312$。
    \item \term{有効混合角}{effective mixing angle}：$Z$ 共鳴でのレプトンの前後非対称性などから定義され、$\sin^2\theta_{\mathrm{eff}}^{\mathrm{lept}}\simeq0.2315$。
\end{itemize}
これらの差（約 $3\%$）は主にトップクォークのループによる $\rho$ の補正
$\Delta\rho\simeq\frac{3G_Fm_t^2}{8\sqrt2\pi^2}\simeq0.0094$ から来る（第12章）。
この量子補正からトップクォークの質量が発見前に予言された。
\subsection{数値のまとめ}
測定値 $G_F$, $M_W$, $M_Z$ からツリーレベルの関係で求めた電弱パラメータを表\ref{tab:EW-parameters}にまとめる。
\begin{table}[htbp]
\centering
\caption{電弱パラメータ（$G_F$, $M_W$, $M_Z$ からツリーレベルで決めた値）}
\label{tab:EW-parameters}
\begin{tabular}{lll}
\toprule
量 & 定義 & 値\\
\midrule
$v$ & $(\sqrt2G_F)^{-1/2}$ & $246.22\ \mathrm{GeV}$\\
$g$ & $2M_W/v$ & $0.653$\quad($\alpha_2=g^2/4\pi\simeq1/29.5$)\\
$g'$ & $2\sqrt{M_Z^2-M_W^2}/v$ & $0.350$\quad($\alpha_Y=g'^2/4\pi\simeq1/103$)\\
$\sin^2\theta_W$ & $1-M_W^2/M_Z^2$ & $0.223$\\
$e$ & $g\sin\theta_W$ & $0.308$\quad($\alpha\simeq1/132$)\\
$\lambda$ & $m_h^2/2v^2$ & $0.129$\\
\bottomrule
\end{tabular}
\end{table}
こうして求めた $\alpha\simeq1/132$ は、Thomson 極限の $\alpha(0)=1/137.036$ とも $\alpha(M_Z)\simeq1/128$ とも異なる。
これはツリーレベルの関係に量子補正があるためで、どの量を入力に選ぶか（\term{繰り込みスキーム}{renormalization scheme}）によって
ツリーレベルの予言が数％ずれることを示している。
\begin{remark}
標準模型では $\sin^2\theta_W$ は自由なパラメータである。
これに対し $\SU(5)$ などの大統一理論では、$\U(1)_Y$ が単純群に埋め込まれるため、
規格化 $g_1=\sqrt{5/3}\,g'$ のもとで統一スケールで $g_1=g$、すなわち
\[
\sin^2\theta_W=\frac{g'^2}{g^2+g'^2}=\frac{\frac35g^2}{g^2+\frac35g^2}=\frac38
\]
が予言される。これを繰り込み群で電弱スケールまで走らせると $\sin^2\theta_W\approx0.21$ となり、
実測値 $0.231$ に近いが一致はしない（第14章）。
\end{remark}
%------------------------------------------------
\section{電弱相互作用の結合定数}
\label{sec:EW-couplings}
%------------------------------------------------
\subsection{電気素量}
\eqref{eq:D-mass-basis} の光子の結合から、電気素量は
\begin{equation}
\boxed{e=g\sin\theta_W=g'\cos\theta_W=\frac{gg'}{\sqrt{g^2+g'^2}}}
\qquad\Longleftrightarrow\qquad
\frac{1}{e^2}=\frac{1}{g^2}+\frac{1}{g'^2}
\label{eq:e-gg}
\end{equation}
で与えられる（第1章 \eqref{eq:e-preview}）。
$e<g$, $e<g'$ であり、電磁相互作用は $\SU(2)_L$ と $\U(1)_Y$ のどちらよりも弱い。
\subsection{3種類のカレント結合}
\eqref{eq:D-mass-basis} をフェルミオンの運動項 $\bar\psi i\gamma^\mu D_\mu\psi$ に代入すると、相互作用は
\begin{equation}
\Lag_{\mathrm{int}}=\frac{g}{\sqrt2}\left(W^+_\mu J^{\mu+}_W+W^-_\mu J^{\mu-}_W\right)
+\frac{g}{\cos\theta_W}Z_\mu J^\mu_Z+eA_\mu J^\mu_{\mathrm{EM}}
\label{eq:EW-currents}
\end{equation}
の形になる。それぞれ
\begin{itemize}
    \item \textbf{荷電カレント}\en{charged current}：結合定数 $g/\sqrt2$、左巻き二重項の上下を入れ替える。
    \item \textbf{中性カレント}\en{neutral current}：結合定数 $g/\cos\theta_W$、カレントは $J_Z=J^3-\sin^2\theta_WJ_{\mathrm{EM}}$。
    \item \textbf{電磁カレント}\en{electromagnetic current}：結合定数 $e$、ベクトル的（左右対称）。
\end{itemize}
である。カレントの具体形とその性質は第8章で詳しく調べる。
\subsection{入力パラメータ}
ゲージ・ヒッグス部分の4つのパラメータ $(g,g',v,\lambda)$ は、精度よく測られている4つの量
\begin{equation}
\alpha=\frac{e^2}{4\pi},\qquad G_F,\qquad M_Z,\qquad m_h
\end{equation}
で置き換えるのが標準的である（\term{入力パラメータ}{input parameters}）。
すると $M_W$ や $\sin^2\theta_W$ は予言になる。
\begin{myex}{$M_W$ の予言}{MW-prediction}
$\alpha$, $G_F$, $M_Z$ から $M_W$ をツリーレベルで求めよ。
\tcblower
$e=g\sin\theta_W$ と $\frac{G_F}{\sqrt2}=\frac{g^2}{8M_W^2}$ から $g$ を消去すると
\[
M_W^2\sin^2\theta_W=\frac{\pi\alpha}{\sqrt2G_F}\equiv A^2 .
\]
$\sin^2\theta_W=1-M_W^2/M_Z^2$ を代入して $M_W^4-M_Z^2M_W^2+A^2M_Z^2=0$。大きい方の解をとって
\[
M_W^2=\frac{M_Z^2}{2}\left(1+\sqrt{1-\frac{4A^2}{M_Z^2}}\right).
\]
数値は、
\begin{itemize}
    \item $\alpha(0)=1/137.036$ を使うと $A=37.28\ \mathrm{GeV}$、$M_W\simeq80.94\ \mathrm{GeV}$、$\sin^2\theta_W\simeq0.212$。
    \item 電弱スケールでの値 $\alpha(M_Z)\simeq1/127.95$ を使うと $A=38.58\ \mathrm{GeV}$、$M_W\simeq79.83\ \mathrm{GeV}$、$\sin^2\theta_W\simeq0.234$。
\end{itemize}
測定値 $M_W=80.37\ \mathrm{GeV}$ はそのどちらとも $0.5\ \mathrm{GeV}$ 程度ずれている。
このずれは量子補正（$\alpha$ の走り、トップクォークとヒッグス粒子のループ）で説明され、
その精密な比較が標準模型の検証になる（第12章）。
\end{myex}
\begin{note}
$e^2=g^2\sin^2\theta_W$ と $g^2=8M_W^2G_F/\sqrt2=4\sqrt2G_FM_W^2$ より $4\pi\alpha=4\sqrt2G_FM_W^2\sin^2\theta_W$、
したがって $M_W^2\sin^2\theta_W=\frac{\pi\alpha}{\sqrt2G_F}$。
\end{note}
測定値から求めたツリーレベルの結合定数は
\begin{equation}
g=\frac{2M_W}{v}\simeq0.653,\qquad
\sqrt{g^2+g'^2}=\frac{2M_Z}{v}\simeq0.741,\qquad
g'\simeq0.350
\end{equation}
であり、$e\simeq0.31$ と比べて確かに $g$ と $g'$ は電荷より大きい。
「弱い相互作用」が弱いのは結合定数が小さいからではなく、
$W$, $Z$ が重く、低エネルギーで $1/M_W^2$ だけ suppress されるからである（第1章、第8章）。
%------------------------------------------------
\section{ヒッグス粒子の物理的自由度}
\label{sec:higgs-dof}
%------------------------------------------------
\subsection{自由度の数え上げ}
ヒッグス二重項の4つの実スカラーのうち、3つ（$G^\pm$, $G^0$）は $W^\pm$, $Z$ に食べられ、
1つ（$h$）が物理的なヒッグス粒子として残る。電弱セクター全体では
\begin{equation}
\underbrace{4}_{H}+\underbrace{4\times2}_{W^{1,2,3},\,B}=12
\quad\longrightarrow\quad
\underbrace{1}_{h}+\underbrace{3\times3}_{W^\pm,\,Z}+\underbrace{2}_{\gamma}=12
\end{equation}
となる（第1章、第5章）。
\begin{table}[htbp]
\centering
\caption{ヒッグス二重項の成分とその行方}
\label{tab:higgs-components}
\begin{tabular}{cccc}
\toprule
成分 & 電荷 $Q$ & 実自由度 & 行方\\
\midrule
$G^\pm$（$=\phi^\pm$） & $\pm1$ & 2 & $W^\pm$ の縦偏極\\
$G^0$（$=\sqrt2\,\Im\phi^0$） & $0$ & 1 & $Z$ の縦偏極\\
$h$（$=\sqrt2\,\Re\phi^0-v$） & $0$ & 1 & ヒッグス粒子\\
\bottomrule
\end{tabular}
\end{table}
$R_\xi$ ゲージでは $G^\pm$, $G^0$ はラグランジアンに残り、それぞれ質量 $\sqrt\xi M_W$, $\sqrt\xi M_Z$ を持つ
（ゴーストも同様）。物理的な振幅ではこれらの寄与は $\xi$ 依存性とともに打ち消し合う。
\subsection{一般のゲージでの NG ボソンと $R_\xi$ ゲージ}
ユニタリーゲージを使わずに直交座標 \eqref{eq:H-cartesian} で展開すると、NG ボソンがどのゲージボソンに食べられるかが2次の項に現れる。
\eqref{eq:DH-components} で $h$ 以外の揺らぎの1次まで残すと、
\begin{equation}
D_\mu H\supset\begin{pmatrix}\partial_\mu G^+-iM_WW^+_\mu\\ \dfrac{1}{\sqrt2}\left[\partial_\mu h+i\left(\partial_\mu G^0+M_ZZ_\mu\right)\right]\end{pmatrix}
\end{equation}
なので、2次の項は
\begin{equation}
(D_\mu H)^\dagger D^\mu H\supset
\left|\partial_\mu G^+-iM_WW^+_\mu\right|^2+\frac12\left(\partial_\mu G^0+M_ZZ_\mu\right)^2+\frac12(\partial_\mu h)^2
\label{eq:SM-goldstone-mixing}
\end{equation}
となる。$G^\pm$ は $W^\pm$ と、$G^0$ は $Z$ と混合し、光子はどの NG ボソンとも混合しない。
これは第5章 \eqref{eq:quadratic-mixing} の構造が、各 massive ゲージボソンごとに現れたものである。
\begin{note}
上成分：$\partial_\mu G^+-\frac{ig}{\sqrt2}W^+_\mu\cdot\frac{v}{\sqrt2}=\partial_\mu G^+-i\frac{gv}{2}W^+_\mu$。
下成分：$\frac{1}{\sqrt2}(\partial_\mu h+i\partial_\mu G^0)+\frac i2(gW^3_\mu-g'B_\mu)\frac{v}{\sqrt2}$ で、
$\frac v2(gW^3-g'B)=\frac v2\sqrt{g^2+g'^2}\,Z=M_ZZ$。
\end{note}
混合項を打ち消すゲージ固定項は
\begin{equation}
\Lag_{\mathrm{gf}}=-\frac{1}{2\xi}\left(\partial^\mu A_\mu\right)^2
-\frac{1}{2\xi}\left(\partial^\mu Z_\mu-\xi M_ZG^0\right)^2
-\frac1\xi\left|\partial^\mu W^+_\mu-i\xi M_WG^+\right|^2
\label{eq:SM-Rxi}
\end{equation}
であり（$R_\xi$ ゲージ）、このとき
\begin{itemize}
    \item $W^\pm$, $Z$ の伝播関数は \eqref{eq:Rxi-propagator} で $m\to M_W,M_Z$ としたもの、光子は通常の共変ゲージの伝播関数。
    \item $G^\pm$ の質量は $\sqrt\xi M_W$、$G^0$ の質量は $\sqrt\xi M_Z$。
    \item ゴースト $c^\pm$, $c^Z$, $c^\gamma$ の質量はそれぞれ $\sqrt\xi M_W$, $\sqrt\xi M_Z$, $0$。
\end{itemize}
となる。ゴーストの相互作用を含む完全な形は付録Fで導出する。
\begin{note}
中性部分：$-\frac{1}{2\xi}(\partial Z-\xi M_ZG^0)^2\supset M_ZG^0\partial\cdot Z\simeq-M_ZZ^\mu\partial_\mu G^0$（部分積分）で、
\eqref{eq:SM-goldstone-mixing} の $+M_ZZ^\mu\partial_\mu G^0$ を打ち消す。残りは $-\frac\xi2M_Z^2(G^0)^2$。
荷電部分：$|\partial G^+-iM_WW^+|^2\supset-iM_WW^{+\mu}\partial_\mu G^-+iM_WW^{-\mu}\partial_\mu G^+$。
ゲージ固定項からは $iM_WG^+\partial\cdot W^--iM_WG^-\partial\cdot W^+$ が出て、部分積分すると上と逆符号になる。
第5章の Abelian Higgs model と $G$ の符号が逆なのは、$Z\propto gW^3-g'B$ と定義したためである。
\end{note}
\subsection{ヒッグス粒子の性質}
\begin{itemize}
    \item \textbf{質量}：$m_h^2=2\lambda v^2=2\mu^2$。$\lambda$ は他の観測量から決まらないため、
    標準模型は $m_h$ を予言しない。2012年に LHC（ATLAS, CMS）で発見され、$m_h\simeq125\ \mathrm{GeV}$。
    \item \textbf{量子数}：電荷 0、カラー一重項、スピン 0、$J^{PC}=0^{++}$（CP 偶）。
    \item \textbf{ゲージ不変性}：ユニタリーゲージで $H^\dagger H=\frac12(v+h)^2$ なので、
    $h=\sqrt{2H^\dagger H}-v$ はゲージ不変に定義される。
    \item \textbf{結合}：ゲージボソンとの結合は $M_V^2/v$ \eqref{eq:hVV-couplings}、
    フェルミオンとの結合は $m_f/v$（第7章）と、いずれも質量に比例する。
    \item \textbf{自己結合}：\eqref{eq:V-cartesian} で $G=0$ とおくと
    \begin{equation}
    V\supset\lambda vh^3+\frac\lambda4h^4=\frac{m_h^2}{2v}h^3+\frac{m_h^2}{8v^2}h^4
    \end{equation}
    であり、第4章 \eqref{eq:Z2-couplings} と同じ形になる（第7章）。
\end{itemize}
\begin{remark}
$m_h\simeq125\ \mathrm{GeV}$ は $v$ と同程度であり、$\lambda\simeq0.13$ は摂動論が使える大きさである。
一方、ヒッグス粒子の質量の2乗は量子補正で高いエネルギースケールの2乗に比例して大きくなりうる。
なぜ $m_h$ が電弱スケールにとどまるのかという問題（\term{階層性問題}{hierarchy problem}）は第14章で議論する。
\end{remark}
\subsection{等価定理とユニタリー性}
\label{subsec:unitarity}
食べられた NG ボソンは、高エネルギーで再び姿を現す。
縦偏極ベクトル $\epsilon_L^\mu=k^\mu/M+O(M/E)$（第5章 \eqref{eq:epsL-highE}）のため、
縦偏極の $W$, $Z$ を含む振幅は、$E\gg M_W$ で対応する NG ボソンの振幅に一致する：
\begin{equation}
\mathcal M\left(W_L^\pm,Z_L,\dots\right)=(\pm i)^n\,\mathcal M\left(G^\pm,G^0,\dots\right)\left[1+O\!\left(\frac{M_W}{E}\right)\right]
\label{eq:equivalence-theorem}
\end{equation}
（$n$ は外線の縦偏極ボソンの数、位相は規約による）。
これを\term{Goldstone ボソン等価定理}{Goldstone boson equivalence theorem}という（Cornwall--Levin--Tiktopoulos, Lee--Quigg--Thacker, Chanowitz--Gaillard）。
NG ボソンの振幅はポテンシャル \eqref{eq:V-cartesian} から簡単に計算できるので、縦偏極ボソンの高エネルギー散乱の振る舞いがすぐに分かる。
\begin{myex}{$W_L^+W_L^-\to W_L^+W_L^-$ とヒッグス粒子の質量の上限}{WLWL}
等価定理を用いて $W_L^+W_L^-\to W_L^+W_L^-$ の振幅を $s,|t|\gg M_W^2$ で求め、
(a) ヒッグス粒子がない場合と (b) ある場合について、部分波のユニタリー性 $|\Re a_0|\le\frac12$ から制限を導け。
\tcblower
\eqref{eq:V-cartesian} から $G^\pm$ に関係する項は $V\supset\lambda(G^+G^-)^2+2\lambda vhG^+G^-$。
Feynman 則は4点 $-4i\lambda=-2im_h^2/v^2$、$hG^+G^-$ 頂点 $-2i\lambda v=-im_h^2/v$。
$G^+G^-\to G^+G^-$ には4点結合と、$h$ の $s$ チャネル・$t$ チャネル交換が寄与し、
\[
\mathcal M=-\frac{m_h^2}{v^2}\left[2+\frac{m_h^2}{s-m_h^2}+\frac{m_h^2}{t-m_h^2}\right]
=-\frac{m_h^2}{v^2}\left[\frac{s}{s-m_h^2}+\frac{t}{t-m_h^2}\right].
\]
\textbf{(a) ヒッグス粒子がない場合}（$m_h\to\infty$ の極限に相当）：
$\mathcal M\to\frac{s+t}{v^2}$ となり、振幅はエネルギーの2乗に比例して増大する。
$J=0$ 部分波 $a_0=\frac{1}{32\pi}\int_{-1}^1d\cos\theta\,\mathcal M$ は、$t=-\frac s2(1-\cos\theta)$ の平均が $-\frac s2$ なので
\[
a_0=\frac{1}{16\pi}\cdot\frac{s-s/2}{v^2}=\frac{s}{32\pi v^2}.
\]
$|a_0|\le\frac12$ より $\sqrt s\lesssim\sqrt{16\pi}\,v\simeq1.7\ \mathrm{TeV}$。
すなわちヒッグス粒子がなければ、このエネルギーまでに摂動論的なユニタリー性が破れ、新しい物理が現れなければならない。
\textbf{(b) ヒッグス粒子がある場合}：$s,|t|\gg m_h^2$ では $\mathcal M\to-\frac{2m_h^2}{v^2}$ と定数に近づき、増大は止まる。
$a_0=-\frac{m_h^2}{8\pi v^2}$ で、$|a_0|\le\frac12$ より $m_h\lesssim\sqrt{4\pi}\,v\simeq870\ \mathrm{GeV}$。
$W_L^+W_L^-$, $Z_LZ_L$, $hh$, $hZ_L$ の結合したチャネルを考えるとより強く、
\[
m_h^2\le\frac{8\pi\sqrt2}{3G_F}\simeq(1.0\ \mathrm{TeV})^2
\]
となる（\term{Lee--Quigg--Thacker の上限}{Lee--Quigg--Thacker bound}）。
\end{myex}
\begin{note}
$s$ チャネル：$G^+G^-\to h\to G^+G^-$ で $(-im_h^2/v)^2\cdot\frac{i}{s-m_h^2}$。
$t$ チャネル：$G^+\to G^+$ の間で $h$ を交換し、同じ形で $s\to t$。$G^+G^-$ の交換による $u$ チャネルはない（$h$ は $G^+G^+$ と結合しない）。
4点結合：$\lambda(G^+G^-)^2$ の縮約は $2\times2=4$ 通りで $-4i\lambda$。
和をとると $i\mathcal M=-\frac{2im_h^2}{v^2}-\frac{im_h^4}{v^2}\left[\frac{1}{s-m_h^2}+\frac{1}{t-m_h^2}\right]$。
$2+\frac{m_h^2}{s-m_h^2}+\frac{m_h^2}{t-m_h^2}=\frac{s}{s-m_h^2}+\frac{t}{t-m_h^2}$ を用いた。
\end{note}
この例は、標準模型の中でヒッグス粒子が果たす役割を端的に示している。
\begin{itemize}
    \item ゲージボソンの質量だけを（例えば電弱カイラルラグランジアン \eqref{eq:EW-chiral-L} で $h$ を除いて）与えた理論は、
    $\sim4\pi v\sim3\ \mathrm{TeV}$ 以下で破綻する\textbf{有効理論}にすぎない。
    \item ヒッグス粒子の交換は、$W_LW_L$ 散乱の $E^2$ に比例する増大を\textbf{ちょうど}打ち消す。
    そのためには $hWW$ 結合がちょうど $gM_W$（\eqref{eq:hVV-couplings}）でなければならない。
    \item 実際、LHC 以前から「TeV スケール以下に、ヒッグス粒子か何らかの新しい物理が必ず存在する」ことが分かっていた。
    これが LHC の重要な設計目標の1つであった。
\end{itemize}
%------------------------------------------------
\section*{第6章のまとめ}
%------------------------------------------------
\begin{itemize}
    \item 標準模型のヒッグス場は $(\mathbf1,\mathbf2)_1$ の複素二重項で、ポテンシャルは $V=-\mu^2H^\dagger H+\lambda(H^\dagger H)^2$ に限られる。
    ポテンシャルはカストディアル対称性 $\SU(2)_L\times\SU(2)_R$ を持つ。
    高温では対称性が回復し、宇宙は $T\simeq160\ \mathrm{GeV}$ で電弱対称性の破れた相に移った（クロスオーバー）。
    \item VEV は $\langle H\rangle=(0,v/\sqrt2)^T$、$v=(\sqrt2G_F)^{-1/2}\simeq246\ \mathrm{GeV}$ にとれ、
    $Q=T_3+Y/2$ が真空を変えないので $\SU(2)_L\times\U(1)_Y\to\U(1)_{\mathrm{EM}}$ と破れる。
    \item ゲージボソンの質量固有状態と質量は
    $W^\pm=(W^1\mp iW^2)/\sqrt2$（$M_W=gv/2$）、$Z=\cos\theta_WW^3-\sin\theta_WB$（$M_Z=\sqrt{g^2+g'^2}\,v/2$）、
    $A=\sin\theta_WW^3+\cos\theta_WB$（massless）である。
    \item Weinberg 角は $\tan\theta_W=g'/g$ で定義され、$M_W=M_Z\cos\theta_W$（$\rho=1$）、$e=g\sin\theta_W=g'\cos\theta_W$ が成り立つ。
    \item 共変微分は $D_\mu=\partial_\mu-\frac{ig}{\sqrt2}(W^+T^++W^-T^-)-\frac{ig}{\cos\theta_W}Z(T^3-\sin^2\theta_WQ)-ieAQ$ と書ける。
    \item ゲージ場の運動項は $WW\gamma$, $WWZ$ の3点結合と4点結合を含み、$W$ の磁気モーメント（$g_W=2$）などを予言する。
    \item 4つの実スカラーのうち3つが $W^\pm,Z$ に食べられ、1つが CP 偶のヒッグス粒子 $h$ として残る。
    その結合はすべて質量に比例する。
    \item 等価定理により、縦偏極ボソンの高エネルギー散乱は NG ボソンの散乱で近似できる。
    ヒッグス粒子の交換は $W_LW_L$ 散乱の増大を打ち消し、ユニタリー性を保つ（$m_h\lesssim1\ \mathrm{TeV}$）。
\end{itemize}
\end{document}
